% Les solutions tampons — Chimie Tle E — Cours signé OnBuch+
% Programme officiel MINESEC, Terminale C, D, E. Compiler avec : tectonic L10-solutions-tampons.tex
\newcommand{\DOCMATIERE}{Chimie}
\newcommand{\DOCNIVEAU}{Tle E}
\newcommand{\DOCMODULE}{Module 2 · Acides et bases}
\newcommand{\DOCLECON}{10}
\newcommand{\DOCTITRE}{Les solutions tampons}
% OnBuch+ — préambule « cours signés » ENRICHI (compilé avec tectonic / XeLaTeX)
% Système visuel « Pop » d'OnBuch+ : fond crème (jamais blanc), bordures encre
% épaisses, ombres « dures » décalées, titres Archivo Black en capitales.
% Version enrichie pour les cours de chimie : chimie (mhchem, chemfig),
% graphes (pgfplots), boîtes pédagogiques supplémentaires (définition,
% propriété, méthode, attention, le savais-tu, exercice résolu, exercice,
% TP, bilan), page de garde refaite et signature OnBuch+ ancrée en bas.
%
% Macros attendues dans main.tex AVANT \input{preamble} :
%   \DOCMATIERE  \DOCNIVEAU  \DOCMODULE  \DOCLECON (numéro)  \DOCTITRE
\documentclass[11pt]{article}

\usepackage{fontspec}
\usepackage[french]{babel}
\usepackage[a4paper,margin=2cm,top=3.4cm,bottom=2.8cm,headheight=1.4cm,footskip=1.3cm]{geometry}
\usepackage{amsmath,amssymb,amsfonts}
\usepackage[table]{xcolor} % [table] : fournit aussi \rowcolors (alternance de lignes)
\usepackage{pagecolor}
\usepackage{tikz}
\usetikzlibrary{arrows.meta,positioning,calc,shapes.geometric,shapes.misc,shapes.arrows,decorations.pathmorphing,decorations.markings,patterns,shadows,mindmap,trees,fit,backgrounds,matrix,intersections}
\usepackage{pgfplots}
\pgfplotsset{compat=1.18}
\usepackage[version=4]{mhchem}
\usepackage{chemfig}
\usepackage{enumitem}
\usepackage{fancyhdr}
% [most] charge toutes les bibliothèques tcolorbox (listings, minted, raster,
% documentation...) et gonfle inutilement la mémoire TeX sur un document long
% (~300 boîtes) : on ne charge que ce qui est réellement utilisé.
\usepackage[skins,breakable]{tcolorbox}
\usepackage{graphicx}
\usepackage{colortbl}
\usepackage{booktabs}
\usepackage{tabularx}
\usepackage{array}
\usepackage{multicol}
\usepackage{siunitx}
% parse-numbers=false : accepte aussi \SI{2,5 \times 10^{-3}}{...}
\sisetup{locale=FR,output-decimal-marker={,},group-minimum-digits=4,parse-numbers=false}
\usepackage{qrcode}
% \degree : commande fréquemment utilisée par les modèles pour un angle ou une
% température en dehors de \SI{}{} (non fournie par les paquets chargés).
\providecommand{\degree}{^{\circ}}
\usepackage{lastpage}
\usepackage{hyperref}
\hypersetup{hidelinks}

% ============================================================
% Palette « Pop » OnBuch+ — mêmes hex que la classe P (plus_ui.dart)
% ============================================================
\definecolor{poporange}{HTML}{F59321}
\definecolor{popdark}{HTML}{E07A0C}
\definecolor{popcream}{HTML}{FFF6E8}
\definecolor{popink}{HTML}{1C1714}
\definecolor{poppurple}{HTML}{7A5AE0}
\definecolor{popgreen}{HTML}{2E9E6A}
\definecolor{poppink}{HTML}{E0457B}
\definecolor{popblue}{HTML}{2F7DE1}
\definecolor{popgold}{HTML}{C98A12}
\definecolor{popmuted}{HTML}{8A7F76}
\definecolor{popline}{HTML}{EADFD2}
\definecolor{popgray}{HTML}{8A7F76}
\colorlet{popgrey}{popgray}
% Alias tolérants (le modèle nomme parfois les couleurs différemment)
\colorlet{popwhite}{white}
\colorlet{popblack}{popink}
\colorlet{popred}{poppink}
\colorlet{popyellow}{popgold}
% Teintes claires pour fonds de tableaux / zones
\colorlet{popblueL}{popblue!10!white}
\colorlet{poporangeL}{poporange!14!white}
\colorlet{popgreenL}{popgreen!12!white}
\colorlet{poppurpleL}{poppurple!12!white}
\colorlet{poppinkL}{poppink!12!white}
\colorlet{popgoldL}{popgold!14!white}

\pagecolor{popcream}

% ============================================================
% Typographie
% ============================================================
\defaultfontfeatures{Ligatures=TeX}
\setmainfont{PlusJakartaSans-Medium.ttf}[Path=../../../fonts/,BoldFont=PlusJakartaSans-Bold.ttf]
% Maths en sans-serif (Fira Math) avec chiffres et lettres droites de Plus
% Jakarta, pour que les formules se fondent dans le texte.
\usepackage[math-style=ISO,bold-style=ISO]{unicode-math}
\setmathfont{FiraMath-Regular.otf}[Path=../../../fonts/]
\setmathfont{PlusJakartaSans-Medium.ttf}[Path=../../../fonts/,range={up/num,up/{Latin,latin},\mathup}]
\usepackage{icomma} % virgule décimale sans espace en mode math
% Caractères absents de la police de texte (sinon carré vide) : rendus par
% leur équivalent mathématique, en texte comme en formule.
\usepackage{newunicodechar}
\newunicodechar{α}{\ensuremath{\alpha}}
\newunicodechar{Α}{\ensuremath{\alpha}}
\newunicodechar{β}{\ensuremath{\beta}}
\newunicodechar{δ}{\ensuremath{\delta}}
\newunicodechar{✓}{\popcheck{popgreen}}
\newfontfamily\popdisplay{ArchivoBlack-Regular.ttf}[Path=../../../fonts/]
\newfontfamily\poplogo{SpaceGrotesk-Bold.ttf}[Path=../../../fonts/]
\newfontfamily\popsemi{PlusJakartaSans-SemiBold.ttf}[Path=../../../fonts/]
\newfontfamily\popxbold{PlusJakartaSans-ExtraBold.ttf}[Path=../../../fonts/]
\newfontfamily\popxboldls{PlusJakartaSans-ExtraBold.ttf}[Path=../../../fonts/,LetterSpace=3.0]

\setlist[enumerate]{leftmargin=1.7em,itemsep=5pt,topsep=4pt}
\setlist[itemize]{leftmargin=1.7em,itemsep=4pt,topsep=4pt,label={\color{poporange}\textbullet}}
\setlength{\parindent}{0pt}
\setlength{\parskip}{5pt}
\renewcommand{\arraystretch}{1.25}

% chemfig : liaisons un peu plus épaisses, encre Pop
\setchemfig{atom sep=2em,bond style={line width=0.9pt,popink},cram width=3pt}

% pgfplots : style Pop par défaut
\pgfplotsset{
  every axis/.append style={axis line style={popink,line width=1pt},tick style={popink},
    grid style={popline,line width=0.5pt},label style={font=\small},tick label style={font=\footnotesize}},
  popcurve/.style={poporange,line width=1.8pt,no markers,smooth},
}

% ============================================================
% Pill / squiggle
% ============================================================
\newcommand{\poppill}[3]{%
  \tikz[baseline=-0.55ex]\node[draw=popink,line width=1.2pt,rounded corners=2.6pt,%
    fill=#2,inner xsep=6.5pt,inner ysep=3pt]{%
    \popxboldls\fontsize{7}{7}\selectfont\color{#3}\MakeUppercase{#1}};%
}
\newcommand{\popsquiggle}[1]{%
  \tikz[baseline=-0.4ex]{
    \draw[#1,line width=2.6pt,line cap=round]
      plot [smooth] coordinates {(0,0.09) (0.34,0.01) (0.68,0.21) (1.02,0.01) (1.36,0.10)};
  }%
}
% Mot-clé surligné (marqueur orange sous le texte)
\newcommand{\cle}[1]{{\popxbold\color{popdark}#1}}

% ============================================================
% En-tête / pied
% ============================================================
\pagestyle{fancy}
\fancyhf{}
\renewcommand{\headrulewidth}{0pt}
\renewcommand{\footrulewidth}{0pt}
\lhead{\raisebox{-0.15ex}{\poplogo\fontsize{13}{13}\selectfont\color{popink}OnBuch\color{poporange}+}}
\rhead{\poppill{\DOCMATIERE\ \textbullet\ \DOCNIVEAU\ \textbullet\ Leçon \DOCLECON}{white}{popink}}
\lfoot{\popxbold\fontsize{7.6}{7.6}\selectfont\color{popmuted}\MakeUppercase{Cours signé OnBuch+}\ \textbullet\ {\popsemi\DOCTITRE}}
\rfoot{\tikz[baseline=-0.6ex]\node[rounded corners=8.5pt,fill=popink,minimum height=17pt,minimum width=17pt,inner xsep=5pt,inner sep=0]{\popxbold\fontsize{8}{8}\selectfont\color{popcream}\thepage\,/\,\pageref*{LastPage}};}

% ============================================================
% Titres
% ============================================================
\newcommand{\chaptitre}[2]{%
  \begin{tcolorbox}[enhanced,colback=poporange,colframe=popink,boxrule=2.6pt,arc=11pt,
      drop shadow=popink,left=15pt,right=15pt,top=12pt,bottom=11pt,before skip=3pt,after skip=18pt]
    \poppill{#2}{popink}{popcream}\par\vspace{9pt}
    {\raggedright\hyphenpenalty=10000\exhyphenpenalty=10000\popdisplay\fontsize{22}{24}\selectfont\color{popcream}\MakeUppercase{#1}\par}
  \end{tcolorbox}
}
% Section numérotée : pastille orange avec le numéro + titre Archivo
\newcounter{popsec}
\newcommand{\coursec}[1]{%
  \stepcounter{popsec}\setcounter{popsub}{0}%
  \par\vspace{16pt}\noindent
  \begin{minipage}{\linewidth}
  \tikz[baseline=-0.7ex]\node[draw=popink,line width=1.6pt,fill=poporange,rounded corners=4pt,minimum size=22pt,inner sep=0]{\popdisplay\fontsize{12}{12}\selectfont\color{popcream}\thepopsec};%
  \hspace{8pt}{\popdisplay\fontsize{15}{17}\selectfont\color{popink}\MakeUppercase{#1}}\par
  \vspace{4pt}\noindent{\color{poporange}\rule{36pt}{3.6pt}}
  \end{minipage}\par\nopagebreak\vspace{8pt}
}
\newcounter{popsub}
\newcommand{\courssub}[1]{%
  \stepcounter{popsub}%
  \par\vspace{9pt}\noindent{\popxbold\fontsize{11.5}{13}\selectfont\color{popink}{\color{poporange}\thepopsec.\thepopsub}\hspace{6pt}#1}\par\nopagebreak\vspace{4pt}
}

% ============================================================
% Boîtes pédagogiques — même recette : fond blanc, bordure épaisse colorée,
% ombre dure encre, badge plein. Titre optionnel [#1].
% ============================================================
\tcbset{popbase/.style={breakable,enhanced,colback=white,boxrule=2.3pt,arc=9pt,
  shadow={3pt}{-3pt}{0pt}{popink},left=14pt,right=14pt,top=11pt,bottom=11pt,before skip=11pt,after skip=15pt,
  fontupper=\popsemi\fontsize{10.6}{14.5}\selectfont\color{popink},
  fontlower=\popsemi\fontsize{10.6}{14.5}\selectfont\color{popink}}}
% Coche dessinée (le glyphe ✓ n'existe pas dans les polices du document)
\newcommand{\popcheck}[1]{\tikz[baseline=-0.5ex]\draw[#1,line width=1.6pt,line cap=round,line join=round](-0.13,0.0)--(-0.04,-0.09)--(0.13,0.1);}
\newcommand{\popboxhead}[3]{\poppill{#1}{#2}{white}\ifx\relax#3\relax\else\hspace{6pt}{\popxbold\fontsize{10.6}{12}\selectfont\color{popink}#3}\fi\par\vspace{7pt}}

\newtcolorbox{aretenir}[1][]{popbase,colframe=popgreen,before upper={\popboxhead{À retenir}{popgreen}{#1}}}
\newtcolorbox{exemplebox}[1][]{popbase,colframe=poporange,before upper={\popboxhead{Exemple}{poporange}{#1}}}
\newtcolorbox{definition}[1][]{popbase,colframe=popblue,before upper={\popboxhead{Définition}{popblue}{#1}}}
\newtcolorbox{propriete}[1][]{popbase,colframe=poppurple,before upper={\popboxhead{Propriété}{poppurple}{#1}}}
\newtcolorbox{methode}[1][]{popbase,colframe=popgold,before upper={\popboxhead{Méthode}{popgold}{#1}}}
\newtcolorbox{attention}[1][]{popbase,colframe=poppink,before upper={\popboxhead{Attention}{poppink}{#1}}}
\newtcolorbox{experience}[1][]{popbase,colframe=popblue,colback=popblueL,before upper={\popboxhead{Expérience}{popblue}{#1}}}
% « Le savais-tu ? » : carte violette pleine (contexte, culture, Cameroun)
\newtcolorbox{savaistu}[1][]{popbase,colback=poppurple,colframe=popink,
  fontupper=\popsemi\fontsize{10.4}{14.2}\selectfont\color{white},
  before upper={\poppill{Le savais-tu\,?}{white}{poppurple}\ifx\relax#1\relax\else\hspace{6pt}{\popxbold\color{white}#1}\fi\par\vspace{7pt}}}
% Exercice résolu : énoncé en haut, \tcblower puis la solution sur fond crème
\newcounter{popexo}\newcounter{popexr}
\newtcolorbox{exoresolu}[1][]{popbase,colframe=poporange,colbacklower=popcream,
  segmentation style={popink,line width=1.2pt,dashed},
  before upper={\stepcounter{popexr}\popboxhead{Exercice résolu \thepopexr}{poporange}{#1}},
  before lower={\poppill{Solution}{popink}{popcream}\par\vspace{6pt}}}
% Exercice (non résolu en place) — la correction est dans la partie Corrigés
\newtcolorbox{exercice}[1][]{popbase,colframe=popink,
  before upper={\stepcounter{popexo}\popboxhead{Exercice \thepopexo}{popink}{#1}}}
% Corrigé
\newtcolorbox{corrige}[1][]{popbase,colframe=popgreen,colback=popgreenL,
  before upper={\popboxhead{Corrigé}{popgreen}{#1}}}
% Objectifs / prérequis / bilan
\newtcolorbox{objectifs}{popbase,colframe=popink,colback=poporangeL,
  before upper={\popboxhead{Objectifs de la leçon}{popink}{}}}
\newtcolorbox{prerequis}{popbase,colframe=popink,colback=popblueL,
  before upper={\popboxhead{Prérequis}{popblue}{}}}
\newtcolorbox{bilan}{popbase,colframe=popink,colback=white,boxrule=2.8pt,
  before upper={\popboxhead{Fiche bilan}{poporange}{L'essentiel à connaître pour l'examen}}}
% Figure encadrée avec légende
\newtcolorbox{popfigure}[1][]{enhanced,breakable=false,colback=white,colframe=popline,boxrule=1.4pt,arc=8pt,
  left=10pt,right=10pt,top=10pt,bottom=8pt,before skip=10pt,after skip=12pt,halign=center,
  lower separated=false,halign lower=center,
  fontlower=\popsemi\fontsize{9}{11.5}\selectfont\color{popmuted},
  #1}
\newcommand{\legende}[1]{\par\vspace{6pt}{\popsemi\fontsize{9}{11.5}\selectfont\color{popmuted}#1}}

% ============================================================
% Signature OnBuch+
% ============================================================
\newcommand{\poplogoinline}[1]{{\poplogo\fontsize{#1}{#1}\selectfont\color{popink}OnBuch\color{poporange}+}}
\newcommand{\signatureonbuch}{%
  \begin{tcolorbox}[enhanced,colback=popink,colframe=popink,boxrule=2.2pt,arc=10pt,
      drop shadow=poporange,left=14pt,right=14pt,top=10pt,bottom=10pt,before skip=0pt,after skip=0pt]
    \tikz[baseline=-0.7ex]\node[circle,fill=popgreen,draw=popcream,line width=1.3pt,minimum size=22pt,inner sep=0]{\popcheck{white}};%
    \hspace{9pt}\begin{minipage}[c]{0.62\linewidth}
      {\popxbold\fontsize{12}{13}\selectfont\color{popcream}Signé\ \textbullet\ {\poplogo OnBuch\color{poporange}+}}\par\vspace{2pt}
      {\raggedright\popsemi\fontsize{8}{10.5}\selectfont\color{popline}\MakeUppercase{Équipe pédagogique OnBuch+}\\\MakeUppercase{Conforme au programme officiel MINESEC}\par}
    \end{minipage}\hfill
    {\poplogo\fontsize{22}{22}\selectfont\color{popcream}OnBuch\color{poporange}+}
  \end{tcolorbox}%
}

% ============================================================
% Page de garde
% #1 titre · #2 matière · #3 niveau · #4 module · #5 n° leçon · #6 durée ·
% #7 description / objectifs (texte ou itemize)
% ============================================================
\newcommand{\pagegarde}[7]{%
  \thispagestyle{empty}
  \newgeometry{a4paper,margin=1.7cm,top=1.6cm,bottom=1.4cm}%
  \begin{tikzpicture}[remember picture,overlay]
    \fill[poporange!18] ([xshift=-3.2cm,yshift=-3.6cm]current page.north east) circle (4.6cm);
    \fill[poppurple!14] ([xshift=2.4cm,yshift=7.6cm]current page.south west) circle (3.8cm);
  \end{tikzpicture}%
  {\poplogo\fontsize{30}{30}\selectfont\color{popink}OnBuch\color{poporange}+}\hfill
  \raisebox{0.6ex}{\poppill{Cours signé \textbullet\ Premium}{popink}{popcream}}\par
  \vspace{1pt}\popsquiggle{poppurple}\par
  \vspace{8pt}
  \begin{tcolorbox}[enhanced,colback=poporange,colframe=popink,boxrule=2.8pt,arc=13pt,
      drop shadow=popink,left=18pt,right=18pt,top=12pt,bottom=13pt,after skip=10pt]
    \poppill{#2 \textbullet\ #3}{popink}{popcream}\hspace{5pt}\poppill{#4}{white}{popink}\par\vspace{8pt}
    {\popdisplay\fontsize{14}{14}\selectfont\color{popink}LEÇON #5}\par\vspace{4pt}
    {\raggedright\hyphenpenalty=10000\exhyphenpenalty=10000\popdisplay\fontsize{25}{27}\selectfont\color{popcream}\MakeUppercase{#1}\par}
    \vspace{8pt}
    {\popxbold\fontsize{9.5}{9.5}\selectfont\color{popink}\MakeUppercase{Durée conseillée : #6}}
  \end{tcolorbox}
  \begin{tcolorbox}[enhanced,colback=white,colframe=popink,boxrule=2.2pt,arc=10pt,
      drop shadow=popink,left=16pt,right=16pt,top=10pt,bottom=10pt,after skip=8pt]
    \poppill{Dans cette leçon}{poppurple}{white}\par\vspace{5pt}
    \popsemi\fontsize{9.3}{12.6}\selectfont\color{popink}#7
  \end{tcolorbox}
  \noindent\begin{minipage}[t]{0.487\linewidth}
    \begin{tcolorbox}[enhanced,colback=popgreen,colframe=popink,boxrule=2.2pt,arc=10pt,
        drop shadow=popink,left=11pt,right=11pt,top=10pt,bottom=10pt,height=3.1cm,valign=center]
      \tcbox[colback=white,colframe=popink,boxrule=1.3pt,arc=5pt,boxsep=0pt,nobeforeafter,
        left=3pt,right=3pt,top=3pt,bottom=3pt,valign=center]{\qrcode[height=1.9cm]{https://play.google.com/store/apps/details?id=com.luvvix.onbuch&hl=fr}}%
      \hspace{8pt}\begin{minipage}[c]{0.5\linewidth}
        {\popdisplay\fontsize{11}{12.5}\selectfont\color{white}\MakeUppercase{Télécharge OnBuch}}\par\vspace{4pt}
        {\raggedright\popsemi\fontsize{8.4}{11}\selectfont\color{white}Gratuit \textbullet\ Google Play\\Tuteur IA, annales, concours, résultats\par}
      \end{minipage}
    \end{tcolorbox}
  \end{minipage}\hfill
  \begin{minipage}[t]{0.487\linewidth}
    \begin{tcolorbox}[enhanced,colback=poppurple,colframe=popink,boxrule=2.2pt,arc=10pt,
        drop shadow=popink,left=11pt,right=11pt,top=10pt,bottom=10pt,height=3.1cm,valign=center]
      \tikz[baseline=-0.7ex]\node[circle,fill=white,draw=popink,line width=1.3pt,minimum size=1.6cm,inner sep=0]{\popdisplay\fontsize{24}{24}\selectfont\color{poppurple}+};%
      \hspace{8pt}\begin{minipage}[c]{0.6\linewidth}
        {\popdisplay\fontsize{11}{12.5}\selectfont\color{white}\MakeUppercase{Passe à OnBuch+}}\par\vspace{4pt}
        {\raggedright\popsemi\fontsize{8.4}{11}\selectfont\color{white}Tous les cours signés, Tuteur IA illimité, fiches PDF — onglet OnBuch+ de l'app\par}
      \end{minipage}
    \end{tcolorbox}
  \end{minipage}
  \vspace{8pt}
  \popcontenu
  \vfill
  \signatureonbuch
  \restoregeometry
  \newpage
}

% Rangée « Ce cours contient » (page de garde)
\newcommand{\poptile}[3]{% #1 couleur #2 gros texte #3 légende
  \begin{tcolorbox}[enhanced,colback=#1,colframe=popink,boxrule=2pt,arc=9pt,shadow={2.5pt}{-2.5pt}{0pt}{popink},
      left=6pt,right=6pt,top=9pt,bottom=9pt,halign=center,valign=center,height=1.75cm,width=\linewidth,nobeforeafter]
    {\popdisplay\fontsize{12.5}{13}\selectfont\color{white}\MakeUppercase{#2}}\par\vspace{3pt}
    {\centering\popsemi\fontsize{7.8}{9.5}\selectfont\color{white}#3\par}
  \end{tcolorbox}}
\newcommand{\popcontenu}{%
  \par\noindent{\popxboldls\fontsize{7.5}{7.5}\selectfont\color{popmuted}CE COURS SIGNÉ CONTIENT}\par\vspace{7pt}
  \noindent\begin{minipage}[t]{0.235\linewidth}\poptile{popblue}{Cours}{complet, illustré, pas à pas}\end{minipage}\hfill
  \begin{minipage}[t]{0.235\linewidth}\poptile{poporange}{Méthodes}{et exercices résolus}\end{minipage}\hfill
  \begin{minipage}[t]{0.235\linewidth}\poptile{poppink}{Exercices}{type examen + corrigés}\end{minipage}\hfill
  \begin{minipage}[t]{0.235\linewidth}\poptile{popgreen}{Fiche bilan}{l'essentiel à retenir}\end{minipage}\par}

% Sommaire
\newtcolorbox{sommaire}{enhanced,colback=white,colframe=popink,boxrule=2.2pt,arc=10pt,
  drop shadow=popink,left=18pt,right=18pt,top=13pt,bottom=13pt,before skip=6pt,after skip=20pt,
  before upper={\poppill{Sommaire}{poppurple}{white}\par\vspace{9pt}\popsemi\fontsize{10.6}{14.5}\selectfont\color{popink}}}

% Signature de fin de cours — poussée tout en bas de la dernière page
\newcommand{\findecours}{%
  \par\vfill\nopagebreak
  \begin{center}\popsquiggle{poporange}\end{center}\vspace{4pt}
  \signatureonbuch
}

%% FIGFIX-BEGIN v3 — correctifs globaux des figures (docs/onbuch-generation/tools/figfix)
\usepackage{adjustbox}\usepackage{etoolbox}
% toute figure TikZ/pgfplots plus large que la ligne est réduite à la largeur disponible
\BeforeBeginEnvironment{tikzpicture}{\begin{adjustbox}{max width=\linewidth}}
\AfterEndEnvironment{tikzpicture}{\end{adjustbox}}
% légendes pgfplots lisibles par-dessus les courbes
\pgfplotsset{every axis/.append style={legend style={fill=white,fill opacity=0.92,text opacity=1,draw=black!15,inner sep=3pt}}}
% étiquettes d'arêtes (syntaxe edge["..."]) avec fond blanc
\tikzset{every edge quotes/.append style={fill=white,inner sep=1.2pt}}
% bulles de cartes mentales : texte à la ligne dans la bulle, bulle un peu plus grande
\tikzset{every concept/.append style={text width=2.0cm,align=center,minimum size=2.3cm,inner sep=2pt}}
%% FIGFIX-END
\begin{document}

\pagegarde{Les solutions tampons}{Chimie}{Tle E}{Module 2 · Acides et bases}{10}{2 h (cours + TP)}{Cette leçon définit les solutions tampons, solutions dont le pH varie peu par addition modérée d'acide, de base ou par dilution. Elle montre comment les préparer (mélange acide faible/base conjuguée, demi-neutralisation), comment calculer leur pH, comment repérer la zone tampon sur une courbe de dosage, et pourquoi elles sont indispensables en biologie, en médecine et dans l'industrie.
\vspace{4pt}
\begin{multicols}{2}\small
\begin{itemize}
\item À la fin de cette leçon, je sais définir une solution tampon et énoncer ses trois propriétés fondamentales.
\item À la fin de cette leçon, je sais justifier qu'un mélange acide faible/base conjuguée en quantités voisines constitue une solution tampon.
\item À la fin de cette leçon, je sais repérer et interpréter la zone tampon (autour de la demi-équivalence) sur une courbe de dosage pH-métrique.
\item À la fin de cette leçon, je sais préparer une solution tampon de pH = pKa par mélange équimolaire d'un acide faible et de sa base conjuguée, en calculant les volumes à mélanger.
\item À la fin de cette leçon, je sais préparer une solution tampon par demi-neutralisation d'un acide faible par une base forte (ou d'une base faible par un acide fort).
\item À la fin de cette leçon, je sais calculer le pH d'un mélange acide faible/base conjuguée à l'aide de la relation de Henderson-Hasselbalch.
\end{itemize}\end{multicols}}

\begin{sommaire}
\begin{multicols}{2}
\begin{enumerate}
\item Définition et propriétés d'une solution tampon
\item pH d'un mélange acide faible / base conjuguée
\item Préparation d'une solution tampon de pH = pKa
\item Zone tampon et pouvoir tampon sur une courbe de dosage
\item Applications des solutions tampons
\item Travaux pratiques
\item Méthodes et exercices résolus
\item Exercices
\item Corrigés des exercices
\item Fiche bilan
\end{enumerate}
\end{multicols}
\end{sommaire}

\begin{objectifs}
\begin{itemize}
\item À la fin de cette leçon, je sais définir une solution tampon et énoncer ses trois propriétés fondamentales.
\item À la fin de cette leçon, je sais justifier qu'un mélange acide faible/base conjuguée en quantités voisines constitue une solution tampon.
\item À la fin de cette leçon, je sais repérer et interpréter la zone tampon (autour de la demi-équivalence) sur une courbe de dosage pH-métrique.
\item À la fin de cette leçon, je sais préparer une solution tampon de pH = pKa par mélange équimolaire d'un acide faible et de sa base conjuguée, en calculant les volumes à mélanger.
\item À la fin de cette leçon, je sais préparer une solution tampon par demi-neutralisation d'un acide faible par une base forte (ou d'une base faible par un acide fort).
\item À la fin de cette leçon, je sais calculer le pH d'un mélange acide faible/base conjuguée à l'aide de la relation de Henderson-Hasselbalch.
\item À la fin de cette leçon, je sais définir le pouvoir tampon et identifier les conditions qui le maximisent.
\item À la fin de cette leçon, je sais citer des applications des solutions tampons : sang, étalonnage des pH-mètres, industrie pharmaceutique et agroalimentaire.
\end{itemize}
\end{objectifs}

\begin{prerequis}
\begin{itemize}
\item Couples acide/base conjugués et réactions acido-basiques (Leçons 7 et 8)
\item Constante d'acidité Ka, pKa et relation pH = pKa + log([base]/[acide]) (Leçon 8)
\item Courbes de dosage pH-métrique, équivalence et demi-équivalence (Leçon 9)
\item Calculs de quantités de matière, concentrations et dilutions (Première)
\item Force des acides et des bases : acides forts, acides faibles (Leçon 7)
\end{itemize}
\end{prerequis}

\newpage

\coursec{Définition et propriétés d'une solution tampon}

\textbf{Situation d'accroche.} Aminatou, élève en classe de Terminale à Yaoundé, prépare son thé du matin. En y versant quelques gouttes de jus de citron, elle observe que l'infusion vire brutalement de couleur : l'acidité a changé d'un coup. Le soir même, son père revient de l'hôpital avec ses résultats d'analyse : le pH de son sang vaut \num{7,40}, exactement comme le mois précédent, alors qu'il a pourtant mangé des plats acides et bu du vin de palme. Aminatou s'interroge : comment le sang résiste-t-il aux variations de pH ? Pourquoi les médicaments injectables préparés à la pharmacie centrale doivent-ils être \og tamponnés \fg{} avant injection ? Et comment le technicien du laboratoire étalonne-t-il son pH-mètre avant de contrôler l'acidité du lait caillé ? Toutes ces situations reposent sur une même notion, que tu vas découvrir dans cette leçon : la \cle{solution tampon}.

\textbf{Problématique.} Qu'est-ce qu'une solution tampon ? Quelles sont ses propriétés et comment expliquer, au niveau microscopique, sa résistance aux variations de pH ?

\courssub{Expérience introductive : deux liquides face au même acide}

Avant de définir quoi que ce soit, faisons parler l'expérience. Prenons deux béchers identiques et ajoutons-leur exactement la même quantité d'acide chlorhydrique dilué.

\begin{experience}[L'eau pure et le mélange acétique face à un ajout d'acide]
\textbf{Dispositif et protocole.}
\begin{itemize}
\item Bécher n\textsuperscript{o}1 : \SI{100}{mL} d'eau pure distillée, de pH initial égal à \num{7,0}.
\item Bécher n\textsuperscript{o}2 : \SI{100}{mL} d'un mélange contenant de l'acide éthanoïque \ce{CH3COOH} et de l'éthanoate de sodium \ce{CH3COONa} (qui libère les ions \ce{CH3COO-}) en concentrations voisines, de pH initial égal à \num{4,8}.
\item Dans chaque bécher, on ajoute goutte à goutte \SI{1}{mL} d'une solution d'acide chlorhydrique de concentration \SI{0,1}{mol.L^{-1}}, en agitant, et on suit le pH avec un pH-mètre.
\end{itemize}
\textbf{Observations.}
\begin{itemize}
\item Bécher n\textsuperscript{o}1 : le pH chute brutalement de \num{7,0} à \num{3,0}, soit une variation $\Delta\text{pH} = 4$ unités.
\item Bécher n\textsuperscript{o}2 : le pH ne passe que de \num{4,8} à \num{4,7}, soit une variation $\Delta\text{pH} \approx 0,1$ unité.
\end{itemize}
\textbf{Interprétation.} Le même ajout d'acide produit des effets très différents : l'eau pure est sans défense face aux ions \ce{H3O+}, tandis que le mélange acide éthanoïque / ion éthanoate les \og absorbe \fg{} presque entièrement. Ce mélange possède une propriété remarquable : c'est une solution tampon.
\end{experience}

\begin{popfigure}
\begin{tikzpicture}[scale=0.9, every node/.style={font=\small}]
% Bécher 1 : eau pure
\draw[thick, popblue] (0,0) -- (0,2.6) (2.4,0) -- (2.4,2.6) (0,0) -- (2.4,0);
\draw[popblue!50, fill=popblueL] (0.05,0.05) rectangle (2.35,1.5);
\node[popdark] at (1.2,0.8) {eau pure};
\node[popdark, font=\footnotesize] at (1.2,0.35) {pH initial $=7{,}0$};
% pipette 1
\draw[thick, popink] (1.2,3.6) -- (1.2,2.9);
\draw[thick, popink] (1.05,3.6) -- (1.35,3.6);
\draw[-{Stealth[length=2.5mm]}, thick, popink] (1.2,2.85) -- (1.2,2.0);
\node[popink, font=\footnotesize, align=center] at (1.2,4.0) {\SI{1}{mL} de \ce{HCl}};
\node[popink, font=\footnotesize] at (1.2,3.8) {\SI{0,1}{mol.L^{-1}}};
% flèche résultat 1
\draw[-{Stealth[length=3mm]}, very thick, poporange] (2.8,1.3) -- (4.0,1.3);
\node[poporange, font=\footnotesize, align=center] at (3.4,1.75) {chute};
\node[poporange, font=\footnotesize, align=center] at (3.4,0.85) {brutale};
% état final 1
\draw[thick, poporange] (4.4,0) -- (4.4,2.6) (6.8,0) -- (6.8,2.6) (4.4,0) -- (6.8,0);
\draw[poporange!60, fill=poporangeL] (4.45,0.05) rectangle (6.75,1.5);
\node[popdark] at (5.6,0.8) {pH $=3{,}0$};
\node[popdark, font=\footnotesize] at (5.6,0.35) {$\Delta\text{pH}=4$};
% Bécher 2 : tampon
\draw[thick, popgreen] (8.2,0) -- (8.2,2.6) (10.6,0) -- (10.6,2.6) (8.2,0) -- (10.6,0);
\draw[popgreen!40, fill=popgreenL] (8.25,0.05) rectangle (10.55,1.5);
\node[popdark, font=\footnotesize] at (9.4,0.9) {\ce{CH3COOH}/\ce{CH3COO-}};
\node[popdark, font=\footnotesize] at (9.4,0.4) {pH initial $=4{,}8$};
% pipette 2
\draw[thick, popink] (9.4,3.6) -- (9.4,2.9);
\draw[thick, popink] (9.25,3.6) -- (9.55,3.6);
\draw[-{Stealth[length=2.5mm]}, thick, popink] (9.4,2.85) -- (9.4,2.0);
\node[popink, font=\footnotesize, align=center] at (9.4,4.0) {\SI{1}{mL} de \ce{HCl}};
\node[popink, font=\footnotesize] at (9.4,3,8) {\SI{0,1}{mol.L^{-1}}};
% flèche résultat 2
\draw[-{Stealth[length=3mm]}, very thick, popgreen] (11.0,1.3) -- (12.2,1.3);
\node[popgreen, font=\footnotesize, align=center] at (11.6,1.75) {variation};
\node[popgreen, font=\footnotesize, align=center] at (11.6,0.85) {faible};
% état final 2
\draw[thick, popgreen] (12.6,0) -- (12.6,2.6) (15.0,0) -- (15.0,2.6) (12.6,0) -- (15.0,0);
\draw[popgreen!40, fill=popgreenL] (12.65,0.05) rectangle (14.95,1.5);
\node[popdark] at (13.8,0.8) {pH $=4{,}7$};
\node[popdark, font=\footnotesize] at (13.8,0.35) {$\Delta\text{pH}=0{,}1$};
\end{tikzpicture}
\legende{Le même ajout d'acide chlorhydrique provoque une chute de pH de 4 unités dans l'eau pure, mais seulement de 0,1 unité dans le mélange tampon acide éthanoïque / ion éthanoate.}
\end{popfigure}

\courssub{Définition d'une solution tampon}

L'expérience précédente nous donne l'intuition : certaines solutions se comportent comme des \og amortisseurs \fg{} de pH, à la manière des amortisseurs d'une moto-taxi qui absorbent les chocs des nids-de-poule sur la route de Douala. Le choc existe toujours, mais ses effets sont fortement atténués. Traduisons cette intuition en langage rigoureux.

\begin{definition}[Solution tampon]
Une \cle{solution tampon} est une solution dont le \cle{pH varie peu} lors de :
\begin{itemize}
\item l'addition \cle{modérée} d'un acide ;
\item l'addition \cle{modérée} d'une base ;
\item une \cle{dilution modérée} (ajout d'eau distillée).
\end{itemize}
On dit aussi que la solution \og tamponne \fg{} le pH, ou qu'elle possède un \emph{effet tampon} (ou \cle{pouvoir tampon}).
\end{definition}

Deux mots de cette définition méritent que tu t'y arrêtes. D'abord \og varie \emph{peu} \fg{} : le pH n'est pas figé, il bouge légèrement. Ensuite \og \emph{modérée} \fg{} : la résistance du tampon a des limites, que nous préciserons à la fin de cette section.

\begin{aretenir}[Les trois propriétés fondamentales d'une solution tampon]
Une solution tampon possède trois propriétés à connaître par cœur :
\begin{enumerate}
\item son pH \textbf{varie peu} par addition modérée d'\textbf{acide} ;
\item son pH \textbf{varie peu} par addition modérée de \textbf{base} ;
\item son pH est \textbf{quasi invariant par dilution modérée}.
\end{enumerate}
\end{aretenir}

\courssub{Nature chimique d'une solution tampon}

Qu'est-ce qui, dans le bécher n\textsuperscript{o}2, confère cette étonnante résistance ? La réponse tient en une phrase : la solution contient \emph{à la fois} un acide faible \emph{et} sa base conjuguée.

\begin{propriete}[Composition d'une solution tampon]
Une solution tampon est constituée d'un mélange d'un \cle{acide faible} \ce{AH} et de sa \cle{base conjuguée} \ce{A-}, présents en \cle{quantités comparables} (aucun des deux ne doit être négligeable devant l'autre). On dit que le tampon est associé au couple acido-basique \ce{AH}/\ce{A-}.
\end{propriete}

Quelques exemples de tampons classiques :
\begin{itemize}
\item le \cle{tampon éthanoïque} (ou acétique) : couple \ce{CH3COOH}/\ce{CH3COO-}, de $\text{p}K_a = 4{,}8$, qui tamponne autour de pH $= 4{,}8$ ;
\item le \cle{tampon ammoniacal} : couple \ce{NH4+}/\ce{NH3}, de $\text{p}K_a = 9{,}2$, qui tamponne autour de pH $= 9{,}2$ ;
\item le \cle{tampon carbonique} du sang : couple \ce{H2CO3}/\ce{HCO3-}, qui maintient le pH sanguin voisin de \num{7,4} (pK$'_a$ apparent = 6,1).
\end{itemize}

\begin{attention}[Ce qu'une solution tampon n'est pas]
Une solution d'acide fort seul, une solution de base forte seule, ou un mélange acide faible / base \emph{non conjuguée} ne constituent \textbf{pas} des tampons. Il faut impérativement un \textbf{couple conjugué} \ce{AH}/\ce{A-}, avec les deux espèces présentes en quantités comparables. Une solution d'acide éthanoïque pur n'est pas un tampon : il lui manque l'ion éthanoate.
\end{attention}

\courssub{Justification qualitative de l'effet tampon}

Pourquoi ce couple fonctionne-t-il comme un amortisseur ? Parce qu'il dispose de \emph{deux pièges complémentaires}, un pour chaque type d'agresseur.

\textbf{Piège n\textsuperscript{o}1 : la base conjuguée \ce{A-} consomme les ions \ce{H3O+} ajoutés.} Si l'on ajoute un acide (qui apporte des ions \ce{H3O+}), la base \ce{A-} présente dans le tampon réagit selon la réaction quasi totale :
\[ \ce{A- + H3O+ -> AH + H2O} \]
Cette réaction consomme presque tous les ions \ce{H3O+} introduits : au lieu de s'accumuler dans la solution et de faire chuter le pH, ils sont transformés en acide faible \ce{AH}, qui ne s'ionise que très peu. Le pH ne baisse donc que légèrement.

\textbf{Piège n\textsuperscript{o}2 : l'acide faible \ce{AH} consomme les ions \ce{HO-} ajoutés.} Si l'on ajoute une base (qui apporte des ions hydroxyde \ce{HO-}), l'acide \ce{AH} présent dans le tampon réagit selon la réaction quasi totale :
\[ \ce{AH + HO- -> A- + H2O} \]
Les ions \ce{HO-} agresseurs sont neutralisés par l'acide faible : le pH ne monte que légèrement.

\begin{popfigure}
\begin{tikzpicture}[scale=0.95, every node/.style={font=\small}]
% Tampon central
\draw[rounded corners=6pt, thick, poppurple, fill=poppurpleL] (4.6,0.4) rectangle (9.4,3.0);
\node[popdark, font=\bfseries] at (7,2.6) {Solution tampon \ce{AH}/\ce{A-}};
\node[popdark, align=center] at (7,1.6) {\ce{AH} : piège à \ce{HO-}\\[2pt] \ce{A-} : piège à \ce{H3O+}};
% Attaque acide (gauche)
\draw[rounded corners=4pt, thick, popink, fill=poppinkL] (0,1.0) rectangle (3.2,2.2);
\node[popdark, align=center] at (1.6,1.6) {Ajout d'acide\\ \ce{H3O+}};
\draw[-{Stealth[length=3mm]}, very thick, popink] (3.3,1.6) -- (4.5,1.6);
\node[popink, font=\footnotesize, align=center] at (1.6,3.1) {\ce{A- + H3O+ -> AH + H2O}};
% Attaque base (droite)
\draw[rounded corners=4pt, thick, poporange, fill=poporangeL] (10.8,1.0) rectangle (14.0,2.2);
\node[popdark, align=center] at (12.4,1.6) {Ajout de base\\ \ce{HO-}};
\draw[-{Stealth[length=3mm]}, very thick, poporange] (10.7,1.6) -- (9.5,1.6);
\node[poporange, font=\footnotesize, align=center] at (12.4,3.1) {\ce{AH + HO- -> A- + H2O}};
% Résultat
\draw[-{Stealth[length=3mm]}, very thick, popgreen] (7,0.3) -- (7,-0.8);
\node[popgreen, font=\bfseries, align=center] at (7,-1.4) {pH quasi stable};
\end{tikzpicture}
\legende{Double défense du tampon : la base conjuguée \ce{A-} piège les ions \ce{H3O+} ajoutés, l'acide faible \ce{AH} piège les ions \ce{HO-} ajoutés. Le pH reste quasi stable.}
\end{popfigure}

Remarque au passage que chaque attaque transforme une espèce du couple en l'autre : ajouter de l'acide consomme \ce{A-} et forme \ce{AH} ; ajouter de la base consomme \ce{AH} et forme \ce{A-}. Tant que les deux espèces restent présentes en quantités notables, la défense tient. C'est exactement ce que traduit la relation d'Henderson-Hasselbalch, que tu étudieras en détail dans la section suivante :
\[ \text{pH} = \text{p}K_a + \log\frac{[\ce{A-}]}{[\ce{AH}]} \]
Un ajout modéré ne modifie que peu le rapport $\dfrac{[\ce{A-}]}{[\ce{AH}]}$, et comme le logarithme d'un nombre proche de 1 varie très lentement, le pH varie peu.

\textbf{Effet de la dilution.} Que se passe-t-il si l'on ajoute de l'eau distillée au tampon ? Les deux concentrations $[\ce{AH}]$ et $[\ce{A-}]$ sont divisées \emph{par le même facteur} : le rapport $\dfrac{[\ce{A-}]}{[\ce{AH}]}$ reste inchangé. D'après la relation précédente, le pH reste donc quasi constant. C'est la troisième propriété du tampon, et elle est précieuse en pratique : on peut diluer une solution tampon commerciale concentrée sans craindre de déplacer son pH (tant que la dilution reste modérée, c'est-à-dire que les concentrations restent suffisamment élevées pour que l'autoprotolyse de l'eau reste négligeable).

\begin{exemplebox}[Comparaison chiffrée : eau pure contre tampon acétique]
Ajoutons \SI{1}{mL} d'acide chlorhydrique à \SI{0,1}{mol.L^{-1}} (soit $n(\ce{H3O+}) = \SI{0,1}{mol.L^{-1}} \times \SI{1e-3}{L} = \SI{1e-4}{mol}$) dans deux récipients.
\begin{itemize}
\item \textbf{Dans \SI{100}{mL} d'eau pure} : les ions \ce{H3O+} s'accumulent librement. Leur concentration devient $[\ce{H3O+}] \approx \dfrac{\num{1e-4}}{\num{0,101}} \approx \SI{1e-3}{mol.L^{-1}}$, d'où $\text{pH} \approx 3$. Le pH est passé de \num{7,0} à \num{3,0} : une chute de \textbf{4 unités}.
\item \textbf{Dans \SI{100}{mL} de tampon acétique} (par exemple \ce{CH3COOH} et \ce{CH3COO-} à \SI{0,1}{mol.L^{-1}} chacun, pH initial $= \text{p}K_a = 4{,}8$) : les \SI{1e-4}{mol} de \ce{H3O+} sont consommés par \ce{CH3COO-}. Le rapport $\dfrac{[\ce{CH3COO-}]}{[\ce{CH3COOH}]}$ passe de $\dfrac{0,1}{0,1} = 1$ à environ $\dfrac{0,099}{0,101} \approx 0{,}98$, d'où $\text{pH} \approx 4{,}8 + \log(0{,}98) \approx 4{,}79$. La variation n'est que de \textbf{0,1 unité} environ.
\end{itemize}
Le tampon a divisé la variation de pH par un facteur 40 !
\end{exemplebox}

\begin{attention}[Erreur fréquente : le pH d'un tampon ne varie JAMAIS]
Non ! Le pH d'une solution tampon \textbf{varie peu}, il ne varie \textbf{pas du tout}. Deux précisions essentielles :
\begin{enumerate}
\item chaque ajout d'acide ou de base, même petit, déplace légèrement le pH ;
\item la protection n'est efficace que pour des ajouts \textbf{modérés}. Si la quantité d'acide ajoutée dépasse la quantité de base conjuguée \ce{A-} disponible (ou si la base ajoutée dépasse l'acide \ce{AH} disponible), le piège est saturé : le tampon est \textbf{détruit} et le pH chute (ou s'envole) brutalement. On dit que la \emph{capacité tampon} est dépassée.
\end{enumerate}
Au baccalauréat, écrire \og le pH d'un tampon ne change pas \fg{} sans nuance coûte des points : écris toujours \og varie peu \fg{} et \og pour des ajouts modérés \fg{}.
\end{attention}

\begin{exoresolu}[Un tampon peut-il résister à n'importe quel ajout ?]
On dispose de \SI{100}{mL} d'une solution tampon contenant $n(\ce{CH3COOH}) = \SI{10}{mmol}$ et $n(\ce{CH3COO-}) = \SI{10}{mmol}$. On ajoute $n = \SI{2}{mmol}$ d'ions \ce{H3O+} (acide chlorhydrique), puis, dans une seconde expérience, $n' = \SI{15}{mmol}$ d'ions \ce{H3O+}. Le tampon résiste-t-il dans les deux cas ?
\tcblower
\textbf{Premier ajout (\SI{2}{mmol}).} La réaction \ce{CH3COO- + H3O+ -> CH3COOH + H2O} est quasi totale. Dressons le bilan en quantités de matière (en mmol) :
\begin{center}
\begin{tabular}{lcccc}
\toprule
 & \ce{CH3COO-} & \ce{H3O+} & \ce{CH3COOH} & \\
\midrule
Initial & 10 & 2 & 10 & \\
Final & $10-2=8$ & $\approx 0$ & $10+2=12$ & \\
\bottomrule
\end{tabular}
\end{center}
Il reste \SI{8}{mmol} de base conjuguée : le piège n'est pas saturé. Le pH passe de $\text{pH} = \text{p}K_a = 4{,}8$ à $\text{pH} = 4{,}8 + \log\dfrac{8}{12} \approx 4{,}8 - 0{,}18 = 4{,}62$. Variation faible : le tampon \textbf{résiste}.

\textbf{Second ajout (\SI{15}{mmol}).} Les \SI{15}{mmol} de \ce{H3O+} dépassent les \SI{10}{mmol} de \ce{CH3COO-} disponibles : toute la base conjuguée est consommée et il reste environ \SI{5}{mmol} de \ce{H3O+} en excès, libres dans la solution. Le pH chute alors vers des valeurs très acides ($[\ce{H3O+}] \approx \SI{0,05}{mol.L^{-1}}$, pH $\approx 1{,}3$) : le tampon est \textbf{détruit}. La capacité du couple a été dépassée.
\end{exoresolu}

\begin{savaistu}[Des tampons naturels dans la calebasse et la marmite]
Le \cle{vin de palme}, cette boisson blanche récoltée au sommet des raphias dans le littoral camerounais, contient naturellement des couples acide faible / base conjuguée (acides organiques et leurs sels) qui freinent les variations d'acidité pendant les premières heures de fermentation. De même, le \cle{lait caillé} (\og kindirmou \fg{} ou \og pendidam \fg{} chez les éleveurs peuls de l'Adamaoua) doit sa texture et sa conservation à des tampons naturels : les protéines du lait et le couple phosphate maintiennent le pH suffisamment stable pour que les ferments lactiques travaillent sans acidifier brutalement le milieu. Sans ces tampons, le lait tournerait de façon incontrôlée. La chimie des solutions tampons n'est donc pas une invention de laboratoire : la nature l'utilise depuis toujours, et ton organisme aussi — ton sang en est la plus belle preuve, avec son pH verrouillé autour de \num{7,4} grâce au couple \ce{H2CO3}/\ce{HCO3-} (pK$'_a$ apparent = 6,1, rapport $[\ce{HCO3-}]/[\ce{H2CO3}] \approx 20$).
\end{savaistu}

\begin{aretenir}[L'essentiel de la section]
\begin{itemize}
\item Une \cle{solution tampon} est une solution dont le pH varie peu par addition modérée d'acide, de base, ou par dilution modérée.
\item Elle est constituée d'un \cle{acide faible} et de sa \cle{base conjuguée} en quantités comparables (couple \ce{AH}/\ce{A-}).
\item Mécanisme : \ce{A-} consomme les \ce{H3O+} ajoutés (\ce{A- + H3O+ -> AH + H2O}) ; \ce{AH} consomme les \ce{HO-} ajoutés (\ce{AH + HO- -> A- + H2O}).
\item Par dilution, le rapport $\dfrac{[\ce{A-}]}{[\ce{AH}]}$ est inchangé : le pH reste quasi constant.
\item Le tampon a une \cle{capacité limitée} (pouvoir tampon) : un ajout trop important d'acide ou de base le détruit.
\item La relation \textbf{pH = p$K_a$ + log$\dfrac{[\ce{A-}]}{[\ce{AH}]}$} (Henderson-Hasselbalch) permet de quantifier le pH et sa variation.
\end{itemize}
\end{aretenir}

Maintenant que tu sais ce qu'est un tampon et comment il se défend, une question se pose naturellement : comment calculer précisément le pH d'un tel mélange ? C'est l'objet de la section suivante, où la relation $\text{pH} = \text{p}K_a + \log\dfrac{[\ce{A-}]}{[\ce{AH}]}$ deviendra ton outil de travail quotidien.

\coursec{pH d'un mélange acide faible / base conjuguée}

Dans la section précédente, nous avons découvert qu'une \cle{solution tampon} résiste aux variations de pH grâce à la présence simultanée d'un acide faible et de sa base conjuguée en quantités comparables. Mais une question demeure : comment prévoir \emph{la valeur exacte} du pH d'un tel mélange ? C'est toute l'utilité de la relation de Henderson-Hasselbalch, que tu dois maîtriser parfaitement pour le Baccalauréat. Elle transforme un problème de chimie en un simple calcul de logarithme.

\courssub{La relation de Henderson-Hasselbalch}

\textbf{Intuition : une balance entre deux partenaires}\par
Imagine un couple acide/base conjugués \ce{AH}/\ce{A-} comme deux équipes qui se partagent les ions \ce{H3O+} du milieu. Si la base \ce{A-} domine, elle « capture » les ions \ce{H3O+} et le pH monte ; si l'acide \ce{AH} domine, il en libère et le pH descend. Le pH du mélange reflète donc directement le \emph{rapport} entre les deux partenaires. C'est exactement ce que traduit mathématiquement la relation de Henderson-Hasselbalch.

\textbf{Établissement de la relation}\par
Considérons un acide faible \ce{AH} en solution aqueuse. Il réagit partiellement avec l'eau selon l'équilibre (réaction limitée, non totale) :

\[ \ce{AH + H2O <=> A- + H3O+} \]

La constante d'acidité du couple \ce{AH}/\ce{A-} s'écrit :

\[ K_a = \dfrac{[\ce{A-}][\ce{H3O+}]}{[\ce{AH}]} \]

Isolons la concentration en ions oxonium :

\[ [\ce{H3O+}] = K_a \times \dfrac{[\ce{AH}]}{[\ce{A-}]} \]

Appliquons l'opérateur $-\log$ aux deux membres, en utilisant les propriétés du logarithme ($-\log(ab) = -\log a - \log b$ et $-\log(1/x) = \log x$) :

\[ -\log[\ce{H3O+}] = -\log K_a - \log\dfrac{[\ce{AH}]}{[\ce{A-}]} = -\log K_a + \log\dfrac{[\ce{A-}]}{[\ce{AH}]} \]

En reconnaissant $\text{pH} = -\log[\ce{H3O+}]$ et $\text{p}K_a = -\log K_a$, on obtient la relation fondamentale.

\begin{aretenir}[Relation de Henderson-Hasselbalch]
Pour tout couple acide faible / base conjuguée \ce{AH}/\ce{A-} présents simultanément en solution :
\[ \boxed{\ \text{pH} = \text{p}K_a + \log\dfrac{[\ce{A-}]}{[\ce{AH}]}\ } \]
Le numérateur contient toujours la \cle{base} conjuguée, le dénominateur toujours l'\cle{acide}.
\end{aretenir}

\begin{attention}[Le sens du rapport : base sur acide, jamais l'inverse !]
L'erreur la plus fréquente au Baccalauréat consiste à inverser le rapport et à écrire $\log([\ce{AH}]/[\ce{A-}])$. Conséquence mathématique : comme $\log(1/x) = -\log x$, inverser le rapport change le \emph{signe} du terme correctif. Si la correction vaut $-0{,}3$ avec le bon rapport, elle vaut $+0{,}3$ avec le rapport inversé : l'erreur finale sur le pH est donc de $2 \times 0{,}3 = 0{,}6$ unité de pH, soit le \textbf{double} de l'écart. Astuce de mémorisation : plus il y a de base, plus le pH est élevé ; le rapport doit donc \emph{croître} quand la base croît, ce qui impose la base au numérateur.
\end{attention}

\textbf{Le cas particulier fondamental : pH = pKa}\par
Que se passe-t-il lorsque le mélange contient exactement autant d'acide que de base conjuguée ? Si $[\ce{A-}] = [\ce{AH}]$, alors le rapport vaut 1, et comme $\log 1 = 0$ :

\[ \text{pH} = \text{p}K_a + \log 1 = \text{p}K_a \]

\begin{propriete}[Cas de la demi-équivalence]
Lorsque les concentrations de l'acide faible et de sa base conjuguée sont égales, le pH de la solution est égal au p$K_a$ du couple :
\[ [\ce{A-}] = [\ce{AH}] \quad \Longleftrightarrow \quad \text{pH} = \text{p}K_a \]
C'est précisément la situation rencontrée à la \cle{demi-équivalence} d'un dosage acide faible -- base forte, et c'est la recette de base pour préparer un tampon de pH choisi (section suivante).
\end{propriete}

\begin{savaistu}[Pourquoi ce pH est-il si stable ?]
Au voisinage de $[\ce{A-}] = [\ce{AH}]$, la courbe $\text{pH} = f\!\left(\log\frac{[\ce{A-}]}{[\ce{AH}]}\right)$ montre qu'il faut des variations \emph{énormes} du rapport (d'un facteur 10 !) pour déplacer le pH d'une seule unité. C'est ce « verrou logarithmique » qui fait la force des tampons : dans ton sang, le couple \ce{H2CO3}/\ce{HCO3-} maintient le pH à $7{,}4 \pm 0{,}05$ malgré l'acide lactique produit par tes muscles lors d'un match de football.
\end{savaistu}

\textbf{Approximation fondamentale : concentrations apportées = concentrations à l'équilibre}\par
Pour appliquer la relation de Henderson-Hasselbalch, il faut connaître les concentrations \emph{à l'équilibre}. Or, en toute rigueur, l'acide \ce{AH} se dissocie un peu et la base \ce{A-} capte un peu de protons : les concentrations réelles diffèrent légèrement des concentrations apportées.

\begin{propriete}[Approximation admise]
Lorsque l'acide et sa base conjuguée sont présents en quantités \textbf{comparables et non négligeables} (c'est-à-dire dans la zone tampon), chacun des deux partenaires « bloque » la réaction de l'autre avec l'eau (effet d'ion commun). Les dissociations deviennent négligeables, et l'on peut assimiler les concentrations à l'équilibre aux concentrations apportées :
\[ [\ce{AH}]_{\text{eq}} \approx [\ce{AH}]_{\text{apporté}} \qquad \text{et} \qquad [\ce{A-}]_{\text{eq}} \approx [\ce{A-}]_{\text{apporté}} \]
Cette approximation est toujours admise au programme de Terminale.
\end{propriete}

\begin{attention}[Limites de l'approximation]
L'approximation cesse d'être valable si l'un des deux partenaires est en quantité très faible devant l'autre (rapport supérieur à 10 ou inférieur à 0,1) ou si les solutions sont extrêmement diluées. Dans ces cas, on sort de la zone tampon et le calcul rigoureux devient nécessaire — mais il n'est pas exigible au Baccalauréat.
\end{attention}

\textbf{Domaine d'efficacité d'une solution tampon}\par
Analysons la relation de Henderson-Hasselbalch pour deux rapports extrêmes :

\begin{itemize}
\item Si $[\ce{A-}] = 10\,[\ce{AH}]$ : $\text{pH} = \text{p}K_a + \log 10 = \text{p}K_a + 1$ ;
\item Si $[\ce{AH}] = 10\,[\ce{A-}]$ : $\text{pH} = \text{p}K_a + \log 0{,}1 = \text{p}K_a - 1$.
\end{itemize}

\begin{aretenir}[Domaine du pouvoir tampon]
Une solution tampon n'est efficace que si le rapport base/acide reste compris entre $0{,}1$ et $10$, c'est-à-dire si :
\[ \text{p}K_a - 1 \leq \text{pH} \leq \text{p}K_a + 1 \]
Au-delà, l'un des deux partenaires est quasiment épuisé et la solution ne peut plus « absorber » les apports d'acide ou de base : le pouvoir tampon s'effondre.
\end{aretenir}

La figure ci-dessous illustre cette propriété : le pH varie linéairement avec le logarithme du rapport, avec une pente égale à 1, et la zone grisée délimite le domaine d'efficacité du tampon.

\begin{popfigure}
\begin{center}
\begin{tikzpicture}
\begin{axis}[
  width=0.92\linewidth, height=7.2cm,
  xlabel={$\log\dfrac{[\ce{A-}]}{[\ce{AH}]}$},
  ylabel={$\text{pH} - \text{p}K_a$},
  xmin=-2.2, xmax=2.2, ymin=-2.2, ymax=2.2,
  xtick={-2,-1,0,1,2}, ytick={-2,-1,0,1,2},
  grid=both, grid style={popline!40},
  axis line style={popdark},
  label style={popdark},
  clip=false
]
% Zone tampon : bande verticale pour |x| < 1 (équivalent à |pH-pKa| < 1 car y=x)
\addplot[fill=popblueL, draw=none] coordinates {(-1, -2.2) (1, -2.2) (1, 2.2) (-1, 2.2)} \closedcycle;

\addplot[popcurve, domain=-2:2, samples=100] {x};
\addplot[only marks, mark=*, mark size=2.5pt, poporange] coordinates {(0,0)};
\addplot[only marks, mark=*, mark size=2pt, popgreen] coordinates {(-1,-1) (1,1)};
\draw[dashed, popmuted] (axis cs:-1,-2.2) -- (axis cs:-1,2.2);
\draw[dashed, popmuted] (axis cs:1,-2.2) -- (axis cs:1,2.2);
\node[poporange, font=\small, anchor=south west] at (axis cs:0.08,0.05) {$[\ce{A-}]=[\ce{AH}]$ : pH = p$K_a$};
\node[popblue, font=\small, anchor=south] at (axis cs:0,1.9) {zone tampon : p$K_a \pm 1$};
\node[popmuted, font=\footnotesize, anchor=west] at (axis cs:-2.15,-1.6) {pente $=1$};
\end{axis}
\end{tikzpicture}
\end{center}
\legende{Évolution du pH en fonction du logarithme du rapport base/acide (relation de Henderson-Hasselbalch). La droite de pente 1 passe par pH = p$K_a$ lorsque le rapport vaut 1 ; la zone colorée délimite le domaine d'efficacité du tampon, p$K_a - 1 \leq$ pH $\leq$ p$K_a + 1$.}
\end{popfigure}

\textbf{Méthode générale de calcul du pH d'un mélange}\par
Dans la pratique, on ne mélange pas toujours directement l'acide et sa base conjuguée : on peut aussi préparer le tampon en faisant réagir partiellement un acide faible avec une base forte (ou l'inverse). La méthode ci-dessous couvre tous les cas.

\begin{methode}[Calculer le pH d'un mélange tampon en trois étapes]
\begin{enumerate}
\item \textbf{Bilan des quantités.} Écrire l'équation de la réaction éventuelle entre les espèces mélangées (réaction acide fort/base faible ou base forte/acide faible, quasi totale), puis dresser un tableau d'avancement pour déterminer les quantités $n(\ce{AH})$ et $n(\ce{A-})$ présentes \emph{après} réaction. Si l'on mélange directement \ce{AH} et \ce{A-} (ex. acide éthanoïque + éthanoate de sodium), il n'y a pas de réaction prépondérante : les quantités apportées sont conservées.
\item \textbf{Calculer le rapport base/acide.} Former le rapport $\dfrac{n(\ce{A-})}{n(\ce{AH})}$. Comme les deux espèces sont dans le même volume final, les volumes se simplifient : on peut utiliser directement les quantités de matière.
\item \textbf{Appliquer la formule.} Calculer $\text{pH} = \text{p}K_a + \log\dfrac{n(\ce{A-})}{n(\ce{AH})}$, puis vérifier que le résultat est cohérent (compris entre p$K_a - 1$ et p$K_a + 1$ pour un bon tampon).
\end{enumerate}
\end{methode}

\begin{exemplebox}[Remarque utile au Bac : les volumes se simplifient]
Puisque l'acide et la base conjuguée se trouvent dans le \emph{même} volume final $V$, on a :
\[ \dfrac{[\ce{A-}]}{[\ce{AH}]} = \dfrac{n(\ce{A-})/V}{n(\ce{AH})/V} = \dfrac{n(\ce{A-})}{n(\ce{AH})} \]
Inutile donc de calculer les concentrations : le rapport des quantités de matière suffit. De même, si les deux solutions mélangées ont le même volume, le rapport des quantités se réduit au rapport des concentrations initiales. Ce raccourci fait gagner un temps précieux le jour de l'épreuve.
\end{exemplebox}

\textbf{Exemple complet : le tampon acétique}\par

\begin{exoresolu}[pH d'un mélange acide éthanoïque / éthanoate de sodium]
On mélange $V_1 = \SI{50}{mL}$ d'une solution d'acide éthanoïque \ce{CH3COOH} de concentration $C_1 = \SI{0,2}{mol.L^{-1}}$ et $V_2 = \SI{50}{mL}$ d'une solution d'éthanoate de sodium \ce{CH3COONa} de concentration $C_2 = \SI{0,1}{mol.L^{-1}}$. On donne $\text{p}K_a(\ce{CH3COOH}/\ce{CH3COO-}) = 4{,}8$.

Calculer le pH du mélange obtenu.
\tcblower
\textbf{Étape 1 : bilan des quantités.}

L'éthanoate de sodium est un sel totalement dissocié dans l'eau :
\[ \ce{CH3COONa ->[eau] CH3COO- + Na+} \]
Les ions \ce{Na+} sont spectateurs. L'acide éthanoïque et l'ion éthanoate sont les deux partenaires d'un même couple : il n'y a pas de réaction prépondérante entre eux (leur éventuelle réaction \ce{CH3COOH + CH3COO- <=> CH3COO- + CH3COOH} ne change rien aux bilans). Les quantités apportées sont donc conservées :
\begin{align*}
n(\ce{CH3COOH}) &= C_1 V_1 = 0{,}2 \times 50 \times 10^{-3} = \SI{1,0e-2}{mol} \\
n(\ce{CH3COO-}) &= C_2 V_2 = 0{,}1 \times 50 \times 10^{-3} = \SI{5,0e-3}{mol}
\end{align*}

\textbf{Étape 2 : rapport base/acide.}

Les deux espèces partagent le même volume final ($\SI{100}{mL}$), donc les volumes se simplifient :
\[ \dfrac{[\ce{CH3COO-}]}{[\ce{CH3COOH}]} = \dfrac{n(\ce{CH3COO-})}{n(\ce{CH3COOH})} = \dfrac{5{,}0 \times 10^{-3}}{1{,}0 \times 10^{-2}} = 0{,}5 \]

\textbf{Étape 3 : application de la relation de Henderson-Hasselbalch.}
\[ \text{pH} = 4{,}8 + \log 0{,}5 = 4{,}8 - 0{,}3 = \boxed{4{,}5} \]

\textbf{Vérification de cohérence.} Le mélange contient deux fois plus d'acide que de base : le pH doit être \emph{inférieur} au p$K_a$, ce qui est bien le cas ($4{,}5 < 4{,}8$). De plus, $4{,}5 \in [3{,}8\,;\,5{,}8] = [\text{p}K_a - 1\,;\,\text{p}K_a + 1]$ : le mélange se situe dans la zone tampon, c'est une bonne solution tampon.
\end{exoresolu}

\begin{attention}[Pièges classiques de cet exercice]
\begin{itemize}
\item \textbf{Inverser le rapport} donnerait $\text{pH} = 4{,}8 + \log 2 = 5{,}1$ : résultat faux, et pourtant fréquent en copie. Vérifie toujours le sens physique : plus d'acide que de base $\Rightarrow$ pH $<$ p$K_a$.
\item \textbf{Oublier la dilution} n'a ici aucune conséquence car les volumes se simplifient dans le rapport ; mais si l'énoncé demandait les concentrations finales, il faudrait diviser par le volume total $\SI{100}{mL}$.
\item \textbf{Confondre les quantités} : l'éthanoate de sodium apporte la base conjuguée \ce{CH3COO-}, pas de l'acide. Identifie toujours qui est l'acide et qui est la base avant tout calcul.
\end{itemize}
\end{attention}

\textbf{Ce qu'il faut retenir de cette section}\par
Tu sais désormais calculer le pH de n'importe quel mélange acide faible / base conjuguée grâce à la relation de Henderson-Hasselbalch, à condition de respecter trois règles d'or : la base au numérateur, les quantités après réaction éventuelle, et la vérification que le pH obtenu reste dans la zone p$K_a \pm 1$. Dans la section suivante, nous inverserons la démarche : au lieu de calculer le pH d'un mélange donné, nous apprendrons à \emph{fabriquer} un tampon de pH imposé, compétence expérimentale directement mobilisable en TP et très prisée au Baccalauréat.

\coursec{Préparation d'une solution tampon de pH = pKa}

Dans la section précédente, nous avons établi un résultat capital : pour un mélange contenant un acide faible \ce{AH} et sa base conjuguée \ce{A-}, le pH est donné par la relation de Henderson-Hasselbalch :

\[ \text{pH} = \text{p}K_a + \log\frac{[\ce{A-}]}{[\ce{AH}]} \]

Nous avons aussi vu que ce pH varie très peu lorsqu'on ajoute de petites quantités d'acide, de base ou d'eau : c'est précisément le comportement d'une \cle{solution tampon}. Mais une question pratique demeure : \textbf{comment fabriquer concrètement une telle solution au laboratoire ?} Comment faire en sorte, par exemple, d'obtenir un tampon de pH = 4,8 pour étalonner un pH-mètre, ou un tampon de pH = 7,4 pour une expérience de biologie ? C'est l'objet de cette section : passer de la théorie à la \textbf{préparation effective} d'une solution tampon.

\courssub{Le principe général : égaliser les quantités d'acide et de base conjuguée}

\begin{definition}[Condition d'un tampon optimal]
Une solution tampon est d'autant plus efficace que les concentrations de l'acide faible et de sa base conjuguée sont \textbf{proches}. Le cas idéal correspond à :
\[ [\ce{A-}] = [\ce{AH}] \quad \Longleftrightarrow \quad \text{pH} = \text{p}K_a \]
On dit alors que le tampon est préparé \textbf{au pH = pKa du couple}. C'est dans cette situation que le pouvoir tampon est maximal.
\end{definition}

L'intuition est simple : un tampon doit pouvoir « absorber » aussi bien les attaques acides que les attaques basiques. Si le mélange contient autant de \ce{AH} (capable de neutraliser une base ajoutée) que de \ce{A-} (capable de neutraliser un acide ajouté), la défense est symétrique et optimale. Tout le problème de la préparation se ramène donc à une question : \textbf{comment obtenir, dans le même récipient, autant de \ce{AH} qu'il y a de \ce{A-} ?} Il existe trois méthodes classiques.

\courssub{Méthode 1 : mélange équimolaire direct acide faible / sel de la base conjuguée}

\textbf{Idée.} On dispose de deux solutions : l'une contient l'acide faible (par exemple l'acide éthanoïque \ce{CH3COOH}), l'autre contient un sel soluble de la base conjuguée (par exemple l'éthanoate de sodium \ce{CH3COONa}, qui se dissocie totalement en \ce{CH3COO-} et \ce{Na+}). Il suffit de mélanger des \textbf{quantités de matière égales} des deux espèces.

\begin{exemplebox}[Tampon acétique à \SI{0,10}{mol.L^{-1}} par mélange direct]
On veut préparer \SI{100}{mL} de tampon acétique de pH = pKa = 4,8, à partir de deux solutions de même concentration $C = \SI{0,10}{mol.L^{-1}}$ : acide éthanoïque et éthanoate de sodium.

Comme les concentrations sont égales, il suffit de mélanger des \textbf{volumes égaux} :
\[ V(\ce{CH3COOH}) = V(\ce{CH3COONa}) = \SI{50}{mL} \]
Vérification : $n(\ce{CH3COOH}) = 0{,}10 \times 0{,}050 = \SI{5,0e-3}{mol}$ et $n(\ce{CH3COO-}) = 0{,}10 \times 0{,}050 = \SI{5,0e-3}{mol}$. Les quantités sont égales, donc :
\[ \text{pH} = \text{p}K_a + \log\frac{5{,}0 \times 10^{-3}}{5{,}0 \times 10^{-3}} = \text{p}K_a + \log 1 = 4{,}8 \]
Le tampon est prêt : \SI{100}{mL} de solution à pH = 4,8.
\end{exemplebox}

\textbf{Cas général : concentrations différentes.} Si les deux solutions n'ont pas la même concentration, on ne mélange plus des volumes égaux : il faut égaliser les \textbf{quantités de matière} :

\[ C_a V_a = C_b V_b \quad \Longrightarrow \quad \frac{V_a}{V_b} = \frac{C_b}{C_a} \]

où $C_a$ et $V_a$ concernent la solution d'acide faible, $C_b$ et $V_b$ la solution de sel de la base conjuguée.

\begin{exemplebox}[Mélange direct avec concentrations différentes]
On dispose d'acide éthanoïque $C_a = \SI{0,20}{mol.L^{-1}}$ et d'éthanoate de sodium $C_b = \SI{0,10}{mol.L^{-1}}$. On veut \SI{100}{mL} de tampon pH = 4,8.
On a $C_a V_a = C_b V_b \Rightarrow 0{,}20 V_a = 0{,}10 V_b \Rightarrow V_b = 2 V_a$.
Avec $V_a + V_b = \SI{100}{mL}$, on trouve $V_a = \SI{33,3}{mL}$ et $V_b = \SI{66,7}{mL}$.
Vérification : $n(\ce{AH}) = 0{,}20 \times 0{,}0333 = \SI{6,66e-3}{mol} = n(\ce{A-})$. pH = 4,8.
\end{exemplebox}

\begin{attention}[Volumes égaux seulement si concentrations égales]
L'égalité des volumes n'est valable que si $C_a = C_b$. Si l'acide est à \SI{0,20}{mol.L^{-1}} et le sel à \SI{0,10}{mol.L^{-1}}, il faudra deux fois plus de volume de sel que d'acide. Toujours raisonner sur les \textbf{moles}, jamais directement sur les volumes !
\end{attention}

\courssub{Méthode 2 : demi-neutralisation d'un acide faible par une base forte}

\textbf{Idée.} On part d'une seule solution : l'acide faible \ce{AH}. On y verse de la soude (base forte) en quantité \textbf{insuffisante pour tout neutraliser}. La réaction, quasi totale, est :

\[ \ce{CH3COOH + HO- -> CH3COO- + H2O} \]

Chaque mole d'ions \ce{HO-} versée transforme une mole de \ce{AH} en une mole de \ce{A-}. Si l'on neutralise exactement la \textbf{moitié} de l'acide initial, il reste autant de \ce{AH} qu'il s'est formé de \ce{A-} : c'est la \cle{demi-équivalence}, et le mélange obtenu est un tampon de pH = pKa.

\begin{methode}[Calcul des volumes pour une demi-neutralisation]
\begin{enumerate}
\item Choisir un volume $V_a$ de solution d'acide faible de concentration $C_a$. Calculer $n_a = C_a V_a$.
\item Calculer le volume équivalent $V_{bE}$ de base forte $C_b$ : $C_b V_{bE} = C_a V_a$.
\item Le volume à verser pour la demi-équivalence est $V_b = \dfrac{V_{bE}}{2}$.
\item Vérifier par un tableau d'avancement que $n_{\text{final}}(\ce{AH}) = n_{\text{final}}(\ce{A-}) = \dfrac{n_a}{2}$.
\end{enumerate}
\end{methode}

\begin{exoresolu}[Préparer 100 mL de tampon pH = 4,8 par demi-neutralisation]
On dispose d'une solution d'acide éthanoïque à $C_a = \SI{0,10}{mol.L^{-1}}$ et de soude à $C_b = \SI{0,10}{mol.L^{-1}}$. On veut préparer \SI{100}{mL} de tampon de pH = 4,8 (pKa du couple \ce{CH3COOH}/\ce{CH3COO-} = 4,8).
\tcblower
\textbf{Étape 1 : choix du volume d'acide.} Prenons $V_a = \SI{100}{mL}$ d'acide éthanoïque : $n_a = 0{,}10 \times 0{,}100 = \SI{1,0e-2}{mol}$.

\textbf{Étape 2 : volume équivalent.} À l'équivalence, $C_b V_{bE} = C_a V_a$, donc :
\[ V_{bE} = \frac{C_a V_a}{C_b} = \frac{0{,}10 \times 100}{0{,}10} = \SI{100}{mL} \]

\textbf{Étape 3 : demi-équivalence.} On verse la moitié du volume équivalent :
\[ V_b = \frac{V_{bE}}{2} = \SI{50}{mL} \quad \text{soit} \quad n_b = 0{,}10 \times 0{,}050 = \SI{5,0e-3}{mol} \]

\textbf{Étape 4 : tableau d'avancement} (réaction quasi totale) :
\begin{center}
\begin{tabular}{lcccc}
\toprule
 & \ce{CH3COOH} & \ce{HO-} & \ce{CH3COO-} & \ce{H2O} \\
\midrule
État initial (mol) & $1{,}0\times10^{-2}$ & $5{,}0\times10^{-3}$ & $0$ & excès \\
État final (mol) & $5{,}0\times10^{-3}$ & $0$ & $5{,}0\times10^{-3}$ & excès \\
\bottomrule
\end{tabular}
\end{center}

\textbf{Étape 5 : pH et volume final.} Il reste autant d'acide qu'il s'est formé de base conjuguée :
\[ \text{pH} = 4{,}8 + \log\frac{5{,}0\times10^{-3}}{5{,}0\times10^{-3}} = 4{,}8 \]
Le volume final est $V_a + V_b = \SI{150}{mL}$. On peut en prélever \SI{100}{mL} ou, mieux, ajuster $V_a$ initial (ici $V_a = \SI{66,7}{mL}$) pour obtenir exactement \SI{100}{mL} final.
\end{exoresolu}

\begin{attention}[Le piège classique de la demi-équivalence]
À la demi-équivalence, le volume de base versé est la \textbf{moitié du volume équivalent} $V_{bE}$, et \textbf{pas la moitié du volume d'acide initial} ! Les deux ne coïncident que si $C_a = C_b$. Si la soude est deux fois plus concentrée que l'acide ($C_b = 2 C_a$), le volume équivalent vaut déjà $V_a/2$, et la demi-équivalence n'exige que $V_a/4$ de soude. Calcule toujours $V_{bE}$ d'abord.
\end{attention}

\courssub{Méthode 3 : demi-neutralisation d'une base faible par un acide fort}

Le raisonnement est parfaitement symétrique. On part d'une solution de base faible, par exemple l'ammoniac \ce{NH3}, et on y verse de l'acide chlorhydrique (acide fort, totalement dissocié en \ce{H3O+} et \ce{Cl-}) :

\[ \ce{NH3 + H3O+ -> NH4+ + H2O} \]

Cette réaction est quasi totale. À la demi-équivalence, la moitié de l'ammoniac a été transformée en ion ammonium : $[\ce{NH4+}] = [\ce{NH3}]$, donc :

\[ \text{pH} = \text{p}K_a(\ce{NH4+}/\ce{NH3}) = 9{,}2 \]

On obtient ainsi un tampon de pH = 9,2, très utilisé en biochimie et dans l'industrie.

\begin{exemplebox}[Tampon ammoniacal par demi-neutralisation]
On veut \SI{100}{mL} de tampon \ce{NH4+}/\ce{NH3} de pH = 9,2. On a de l'ammoniac \SI{0,10}{mol.L^{-1}} et de l'acide chlorhydrique \SI{0,10}{mol.L^{-1}}.
Prenons $V_{\ce{NH3}} = \SI{100}{mL}$ ($n = \SI{1,0e-2}{mol}$). $V_{E} = \SI{100}{mL}$ d'HCl. Demi-équivalence : $V_{\ce{HCl}} = \SI{50}{mL}$.
Au final : $n(\ce{NH3}) = n(\ce{NH4+}) = \SI{5,0e-3}{mol}$. pH = 9,2. Volume final \SI{150}{mL}.
\end{exemplebox}

\courssub{Illustration synthétique des trois méthodes}

\begin{popfigure}
\begin{tikzpicture}[scale=0.75, every node/.style={transform shape}]
% Méthode 1
\begin{scope}[shift={(0,0)}]
\draw[thick,popblue] (-0.7,0) -- (-0.7,1.5) -- (-0.2,2.4) -- (-0.2,3.1);
\draw[thick,popblue] (0.7,0) -- (0.7,1.5) -- (0.2,2.4) -- (0.2,3.1);
\draw[thick,popblue] (-0.7,0) arc (180:360:0.7);
\draw[popblue!40,fill=popblue!15] (-0.65,0.1) -- (-0.65,1.2) -- (0.65,1.2) -- (0.65,0.1) arc (360:180:0.65);
\draw[popblue] (-0.4,2.7) -- (0.4,2.7);
\node[align=center,font=\small] at (0,-1.0) {\textbf{Méthode 1}\\ Mélange direct\\ \ce{AH} ($C_a,V_a$) + Sel \ce{A-} ($C_b,V_b$)\\ $C_aV_a = C_bV_b$};
\end{scope}
% Méthode 2
\begin{scope}[shift={(5.5,0)}]
\draw[thick,popgreen] (-0.7,0) -- (-0.7,1.5) -- (-0.2,2.4) -- (-0.2,3.1);
\draw[thick,popgreen] (0.7,0) -- (0.7,1.5) -- (0.2,2.4) -- (0.2,3.1);
\draw[thick,popgreen] (-0.7,0) arc (180:360:0.7);
\draw[popgreen!40,fill=popgreen!15] (-0.65,0.1) -- (-0.65,1.2) -- (0.65,1.2) -- (0.65,0.1) arc (360:180:0.65);
\draw[popgreen] (-0.4,2.7) -- (0.4,2.7);
\node[align=center,font=\small] at (0,-1.0) {\textbf{Méthode 2}\\ \ce{AH} + \ce{HO-}\\ Demi-équivalence\\ $V_b = V_{bE}/2$};
\end{scope}
% Méthode 3
\begin{scope}[shift={(11,0)}]
\draw[thick,poporange] (-0.7,0) -- (-0.7,1.5) -- (-0.2,2.4) -- (-0.2,3.1);
\draw[thick,poporange] (0.7,0) -- (0.7,1.5) -- (0.2,2.4) -- (0.2,3.1);
\draw[thick,poporange] (-0.7,0) arc (180:360:0.7);
\draw[poporange!40,fill=poporange!15] (-0.65,0.1) -- (-0.65,1.2) -- (0.65,1.2) -- (0.65,0.1) arc (360:180:0.65);
\draw[poporange] (-0.4,2.7) -- (0.4,2.7);
\node[align=center,font=\small] at (0,-1.0) {\textbf{Méthode 3}\\ Base faible \ce{B} + \ce{H3O+}\\ Demi-équivalence\\ $V_a = V_{aE}/2$};
\end{scope}
\end{tikzpicture}
\legende{Les trois méthodes de préparation d'une solution tampon de pH = pKa. Dans les trois cas, l'objectif est identique : obtenir $n(\ce{A-}) = n(\ce{AH})$ (ou $n(\ce{BH+}) = n(\ce{B})$) dans la fiole.}
\end{popfigure}

\courssub{Choisir le bon couple acide/base}

On ne choisit pas un couple au hasard ! Pour préparer un tampon de pH donné, on sélectionne le couple dont le \textbf{pKa est le plus proche du pH désiré} (idéalement à $\pm 1$ unité près). En effet, hors de l'intervalle $[\text{p}K_a - 1 \,;\, \text{p}K_a + 1]$, l'une des deux espèces devient négligeable (moins de 10\% ou plus de 90\%) et le tampon perd son efficacité.

\begin{center}
\begin{tabularx}{\linewidth}{|l|X|c|c|}
\hline
\rowcolor{popblueL} \textbf{Couple acide/base} & \textbf{Acide / Base conjuguée} & $\text{p}K_a$ (25°C) & \textbf{Usage typique} \\
\hline
Éthanoïque & \ce{CH3COOH}/\ce{CH3COO-} & 4,8 & Étalonnage pH-mètre, agroalimentaire (conserves, vinaigre) \\
\hline
Carbonique & \ce{H2CO3}/\ce{HCO3-} & 6,4 & Tampon sanguin (système ouvert \ce{CO2}/\ce{HCO3-}), boissons gazeuses \\
\hline
Dihydrogénophosphate & \ce{H2PO4-}/\ce{HPO4^2-} & 7,2 & Tampon phosphate biologique, cultures cellulaires \\
\hline
Ammonium & \ce{NH4+}/\ce{NH3} & 9,2 & Biochimie, détergents, traitements de surface \\
\hline
\end{tabularx}
\end{center}

\begin{methode}[Préparer un tampon de pH donné]
\begin{enumerate}
\item \textbf{Choisir le couple} dont le pKa est le plus proche du pH visé (tableau ci-dessus).
\item \textbf{Choisir la méthode} : mélange direct acide + sel, ou demi-neutralisation de l'acide (ou de la base).
\item \textbf{Calculer les quantités} pour obtenir $n(\ce{A-}) = n(\ce{AH})$ : volumes égaux si les concentrations sont égales, sinon $C_aV_a = C_bV_b$ ; à la demi-équivalence, verser $V_{E}/2$.
\item \textbf{Mélanger} dans un bécher, transvaser en fiole jaugée si volume imposé, \textbf{compléter au trait de jauge} si nécessaire, homogénéiser.
\item \textbf{Vérifier le pH} au pH-mètre étalonné (la valeur théorique pH = pKa est approchée ; l'activité ionique et la température peuvent la décaler de quelques centièmes).
\end{enumerate}
\end{methode}

\begin{aretenir}[L'essentiel]
\begin{itemize}
\item Un tampon optimal vérifie $[\ce{A-}] = [\ce{AH}]$, donc pH = pKa.
\item Trois préparations possibles : mélange équimolaire direct, demi-neutralisation d'un acide faible, demi-neutralisation d'une base faible.
\item On choisit le couple dont le pKa est le plus proche du pH désiré (intervalle pKa $\pm$ 1).
\item Tous les calculs se font en \textbf{quantités de matière}, jamais directement en volumes.
\item La dilution ne change pas le pH (rapport des concentrations constant) mais diminue le pouvoir tampon.
\end{itemize}
\end{aretenir}

\begin{savaistu}[Les tampons étalons des pH-mètres]
Les sachets de solutions étalons pH = 4, 7 et 10 que l'on utilise pour calibrer les pH-mètres — y compris dans les laboratoires de la SONARA (raffinerie), des brasseries (Isangui, Ngok) ou des laiteries camerounaises — sont préparés exactement selon les principes de cette leçon : un couple acide/base bien choisi, des concentrations égales des deux partenaires, et une grande pureté des réactifs pour garantir un pH stable et reproductible. Le tampon pH = 4 est souvent un tampon phtalate (pKa = 4,0), le tampon pH = 7 un tampon phosphate (pKa = 7,2) et le tampon pH = 10 un tampon borate (pKa = 9,2).
\end{savaistu}

\coursec{Zone tampon et pouvoir tampon sur une courbe de dosage}

Dans la section précédente, tu as appris à \emph{préparer} une solution tampon de pH = pKa, par exemple en mélangeant des quantités équimolaires d'acide éthanoïque \ce{CH3COOH} et d'ions éthanoate \ce{CH3COO-}. Une question naturelle se pose alors : comment \emph{voir} ce phénomène tampon à l'œuvre, sans rien préparer de spécial ? La réponse est magnifique : il suffit de doser un acide faible par une base forte et de regarder attentivement la courbe. La nature elle-même fabrique une solution tampon au milieu du dosage ! C'est ce que nous allons découvrir maintenant, et c'est une compétence directement exigible au Baccalauréat.

\courssub{Rappel : la courbe de dosage d'un acide faible par une base forte}

Reprenons la situation étudiée à la Leçon 9. On place dans un bécher un volume $V_a = \SI{20,0}{mL}$ d'acide éthanoïque \ce{CH3COOH} de concentration $C_a = \SI{0,10}{mol.L^{-1}}$, et on ajoute progressivement, à la burette, une solution d'hydroxyde de sodium (\ce{Na+} + \ce{HO-}) de concentration $C_b = \SI{0,10}{mol.L^{-1}}$. Un pH-mètre enregistre le pH après chaque ajout.

La réaction de dosage, quasi totale ($K \approx 10^{9,2} \gg 10^4$), s'écrit :

\[ \ce{CH3COOH + HO- -> CH3COO- + H2O} \]

Cette réaction est \textbf{rapide}, \textbf{unique} et \textbf{quasi totale} : elle transforme l'acide éthanoïque en sa base conjuguée, l'ion éthanoate, au fur et à mesure de l'ajout des ions hydroxyde.

L'équivalence est atteinte lorsque $C_a V_a = C_b V_e$, soit :

\[ V_e = \frac{C_a V_a}{C_b} = \frac{0,10 \times 20,0}{0,10} = \SI{20,0}{mL} \]

La courbe pH $= f(V_b)$ présente alors quatre régions que tu dois connaître par cœur :

\begin{itemize}
\item $V_b = 0$ : solution d'acide faible seul, pH $\approx 2,9$ ;
\item $0 < V_b < V_e$ : montée lente du pH, avec un point remarquable à la \cle{demi-équivalence} ;
\item $V_b = V_e = \SI{20,0}{mL}$ : \textbf{saut de pH} brutal (le pH à l'équivalence vaut environ $8,7$, supérieur à 7 car la solution contient alors la base faible \ce{CH3COO-}) ;
\item $V_b > V_e$ : excès de soude, le pH tend vers celui de la base forte diluée.
\end{itemize}

\courssub{La demi-équivalence : le point où pH = pKa}

\textbf{Intuition.} Imagine que tu transformes progressivement un stock de \ce{CH3COOH} en \ce{CH3COO-}, comme on vide un seau dans un autre. À mi-chemin, les deux seaux contiennent exactement la même quantité : il y a autant d'acide restant que de base conjuguée formée. Cet instant précis, c'est la demi-équivalence.

\begin{propriete}[Demi-équivalence]
À la \cle{demi-équivalence}, on a versé un volume $V_b = \dfrac{V_e}{2}$. La moitié de l'acide a été consommée, donc :
\[ n(\ce{CH3COOH})_{\text{restant}} = n(\ce{CH3COO-})_{\text{formé}} \quad \Longrightarrow \quad [\ce{CH3COOH}] = [\ce{CH3COO-}] \]
La relation de Henderson--Hasselbalch donne alors :
\[ \text{pH} = \text{p}K_a + \log\frac{[\ce{CH3COO-}]}{[\ce{CH3COOH}]} = \text{p}K_a + \log 1 = \boxed{\text{p}K_a} \]
\end{propriete}

C'est un résultat capital : \textbf{la lecture graphique du pH à la demi-équivalence fournit directement le p$K_a$ du couple}. Pour l'acide éthanoïque, on lit pH $= 4,8$ à $V_b = \SI{10}{mL}$, donc p$K_a(\ce{CH3COOH}/\ce{CH3COO-}) = 4,8$. C'est d'ailleurs la méthode expérimentale classique pour déterminer un p$K_a$ au laboratoire.

\begin{attention}[Demi-équivalence $\neq$ équivalence]
C'est \emph{la} confusion la plus fréquente au Baccalauréat. Retiens bien :
\begin{itemize}
\item à la \textbf{demi-équivalence} ($V_b = V_e/2$) : $[\ce{AH}] = [\ce{A-}]$, pH $=$ p$K_a$, la courbe est \emph{presque plate} (zone tampon) ;
\item à l'\textbf{équivalence} ($V_b = V_e$) : les réactifs sont en proportions stœchiométriques, la courbe présente un \emph{saut de pH} vertical.
\end{itemize}
Le point plat et le point vertical sont deux points différents de la courbe, séparés d'un volume $V_e/2$. Ne les intervertis jamais !
\end{attention}

\courssub{La zone tampon : là où la courbe s'aplatit}

Observons maintenant la courbe autour de la demi-équivalence, disons entre $V_b \approx \SI{4}{mL}$ et $V_b \approx \SI{16}{mL}$ dans notre exemple. Que constate-t-on ? Le pH ne monte presque plus : il passe péniblement d'environ $3,9$ à $5,7$ alors qu'on a versé \SI{12}{mL} de soude concentrée ! Par comparaison, près de l'équivalence, une seule goutte suffit à faire bondir le pH de plusieurs unités.

\begin{definition}[Zone tampon]
La \cle{zone tampon} d'une courbe de dosage est la région située \textbf{autour de la demi-équivalence}, où le pH varie très lentement lors de l'ajout de base (ou d'acide). Sur la courbe, c'est la portion \textbf{quasi horizontale} (pente faible) encadrant le point pH $=$ p$K_a$.
\end{definition}

\textbf{Interprétation microscopique.} Dans cette région, le mélange contient des quantités \emph{comparables} d'acide \ce{CH3COOH} non encore dosé et de base conjuguée \ce{CH3COO-} déjà formée. Or c'est précisément la définition d'une solution tampon ! Le système se comporte comme un tampon \emph{sans qu'on l'ait préparé} :

\begin{itemize}
\item les ions \ce{HO-} ajoutés sont consommés par l'acide restant : \ce{CH3COOH + HO- -> CH3COO- + H2O} ;
\item le rapport $\dfrac{[\ce{CH3COO-}]}{[\ce{CH3COOH}]}$ varie peu, donc le pH $= \text{p}K_a + \log\dfrac{[\ce{CH3COO-}]}{[\ce{CH3COOH}]}$ varie peu.
\end{itemize}

On considère en pratique que l'effet tampon est efficace tant que $\dfrac{1}{10} < \dfrac{[\ce{A-}]}{[\ce{AH}]} < 10$, c'est-à-dire dans la fourchette $\text{p}K_a - 1 < \text{pH} < \text{p}K_a + 1$.

\begin{methode}[Repérer la zone tampon sur une courbe de dosage]
\begin{enumerate}
\item Repérer le \textbf{saut de pH} et en déduire le volume équivalent $V_e$.
\item Calculer $V_e/2$ et placer le point de demi-équivalence sur la courbe.
\item Lire le pH en ce point : c'est le p$K_a$ du couple.
\item La zone tampon est la portion \textbf{quasi horizontale} de la courbe autour de ce point, approximativement entre p$K_a - 1$ et p$K_a + 1$.
\end{enumerate}
\end{methode}

\begin{popfigure}
\begin{center}
\begin{tikzpicture}
\begin{axis}[
width=0.92\linewidth, height=7.2cm,
xlabel={$V_b$ (mL)}, ylabel={pH},
xmin=0, xmax=40, ymin=0, ymax=14,
xtick={0,10,20,30,40}, ytick={0,2,4,6,8,10,12,14},
grid=major, grid style={popline!40, dashed},
legend style={draw=none, at={(0.03,0.97)}, anchor=north west, font=\small}
]
% Zone tampon grisée
\addplot[draw=none, fill=popblue!18, forget plot] coordinates {(4,0) (16,0) (16,14) (4,14)} -- cycle;
% Branche tampon : Henderson-Hasselbalch
\addplot[popcurve, domain=0.8:19.2, samples=150] {4.76 + ln(x/(20-x))/ln(10)};
% Branche après équivalence
\addplot[popcurve, domain=20.6:40, samples=80] {14 + log10(0.1*(x-20)/(x+20))};
% Point initial
\addplot[only marks, mark=*, popblue, forget plot] coordinates {(0,2.9)};
% Point demi-équivalence
\addplot[only marks, mark=*, poporange, mark size=2.5pt, forget plot] coordinates {(10,4.76)};
\draw[dashed, poporange] (axis cs:10,0) -- (axis cs:10,4.76) -- (axis cs:0,4.76);
\node[poporange, font=\small, anchor=west] at (axis cs:10.5,3.4) {$V_e/2 = \SI{10}{mL}$ ; pH $=$ p$K_a = 4,8$};
% Équivalence
\draw[dashed, poppurple] (axis cs:20,0) -- (axis cs:20,8.7);
\addplot[only marks, mark=*, poppurple, mark size=2.5pt, forget plot] coordinates {(20,8.7)};
\node[poppurple, font=\small, anchor=west] at (axis cs:20.5,7.2) {équivalence $V_e = \SI{20}{mL}$};
% Légende zone tampon
\node[popblue, font=\small\bfseries] at (axis cs:10,12.6) {zone tampon};
\end{axis}
\end{tikzpicture}
\end{center}
\legende{Dosage de \SI{20,0}{mL} de \ce{CH3COOH} à \SI{0,10}{mol.L^{-1}} par \ce{NaOH} à \SI{0,10}{mol.L^{-1}}. La zone grisée, quasi horizontale autour de la demi-équivalence, est la zone tampon : le mélange y contient des quantités comparables de \ce{CH3COOH} et \ce{CH3COO-}.}
\end{popfigure}

\begin{exemplebox}[Lecture de la zone tampon]
Pour le dosage de \SI{20,0}{mL} d'acide éthanoïque à \SI{0,10}{mol.L^{-1}} par la soude à \SI{0,10}{mol.L^{-1}} :
\begin{itemize}
\item volume équivalent : $V_e = \SI{20,0}{mL}$ ;
\item demi-équivalence : $V_b = \SI{10,0}{mL}$, où l'on lit pH $= 4,8 =$ p$K_a$ ;
\item zone tampon : approximativement pour $V_b$ compris entre \SI{4}{mL} et \SI{16}{mL}, soit un pH compris entre $3,8$ et $5,8$ (fourchette p$K_a \pm 1$).
\end{itemize}
Dans toute cette zone, ajouter \SI{1}{mL} de soude ne fait varier le pH que d'environ $0,1$ unité, alors que la même goutte versée près de l'équivalence le ferait varier de plusieurs unités !
\end{exemplebox}

\courssub{Le pouvoir tampon : mesurer l'efficacité du tampon}

Toutes les solutions tampons ne se valent pas. Un tampon dilué s'épuise vite ; un tampon concentré résiste longtemps. Pour comparer les tampons entre eux, les chimistes ont introduit une grandeur : le pouvoir tampon.

\begin{definition}[Pouvoir tampon]
Le \cle{pouvoir tampon} $\beta$ d'une solution mesure son aptitude à s'opposer aux variations de pH. Il est défini par :
\[ \beta = \left| \frac{\mathrm{d}n}{\mathrm{d}\text{pH}} \right| \]
où $\mathrm{d}n$ est la quantité (en mol) d'acide fort ou de base forte ajoutée, et $\mathrm{d}$pH la variation de pH qui en résulte. Plus $\beta$ est \textbf{grand}, plus il faut ajouter d'acide ou de base pour faire bouger le pH, et plus le tampon est \textbf{efficace}.
\end{definition}

\begin{propriete}[Facteurs influençant le pouvoir tampon]
\begin{enumerate}
\item $\beta$ est \textbf{maximal lorsque pH $=$ p$K_a$}, c'est-à-dire lorsque $[\ce{AH}] = [\ce{A-}]$ : les deux « réservoirs » (acide et base conjuguée) sont alors également remplis et prêts à neutraliser aussi bien un acide qu'une base ajoutés.
\item $\beta$ \textbf{augmente avec la concentration totale} du couple $C = [\ce{AH}] + [\ce{A-}]$ : un tampon dix fois plus concentré est environ dix fois plus résistant.
\item $\beta$ devient faible dès que le pH s'écarte de plus d'une unité du p$K_a$ : le tampon est « épuisé » d'un côté.
\end{enumerate}
\end{propriete}

\begin{popfigure}
\begin{center}
\begin{tikzpicture}
\begin{axis}[
width=0.85\linewidth, height=6cm,
xlabel={pH}, ylabel={$\beta$},
xmin=0.5, xmax=9.5, ymin=0, ymax=2.8,
xtick={1,2,3,4,4.76,5,6,7,8,9},
ytick=\empty,
axis lines=left,
]
\addplot[popcurve, domain=0.8:8.8, samples=200] {2.3*exp(-((x-4.76)^2)/1.1)};
\draw[dashed, poporange] (axis cs:4.76,0) -- (axis cs:4.76,2.3);
\addplot[only marks, mark=*, poporange, mark size=2.5pt] coordinates {(4.76,2.3)};
\node[poporange, font=\small, anchor=south] at (axis cs:4.76,2.32) {maximum à pH $=$ p$K_a$};
\draw[<->, popmuted] (axis cs:3.76,1.15) -- (axis cs:5.76,1.15);
\node[popmuted, font=\small, anchor=north] at (axis cs:4.76,1.05) {p$K_a \pm 1$};
\end{axis}
\end{tikzpicture}
\end{center}
\legende{Pouvoir tampon $\beta$ en fonction du pH pour le couple \ce{CH3COOH}/\ce{CH3COO-} : courbe en cloche centrée sur le p$K_a$. L'efficacité du tampon s'effondre dès que $|\text{pH} - \text{p}K_a| > 1$.}
\end{popfigure}

\begin{savaistu}[Pourquoi ton sang ne « saute » pas de pH ?]
Le sang humain est tamponné à pH $\approx 7,4$ principalement par le couple \ce{H2CO3}/\ce{HCO3-}. Son pouvoir tampon est remarquable : même après un effort musculaire intense (qui produit de l'acide lactique) ou l'ingestion d'aliments acides comme le jus de bissap bien sucré, le pH sanguin ne varie que de quelques centièmes d'unité de pH. Une variation de seulement $0,3$ unité serait mortelle ! C'est le même principe que la zone plate de notre courbe de dosage, à l'échelle de tout un organisme.
\end{savaistu}

\courssub{Complément pour le Bac : la courbe dérivée}

Une façon élégante de visualiser le phénomène consiste à tracer la dérivée $\dfrac{\mathrm{d}\text{pH}}{\mathrm{d}V_b}$ en fonction de $V_b$. Cette courbe dérivée présente :

\begin{itemize}
\item un \textbf{maximum très marqué à l'équivalence} : la pente de la courbe de dosage y est vertigineuse (saut de pH) — c'est d'ailleurs une méthode précise pour repérer $V_e$ ;
\item un \textbf{minimum à la demi-équivalence} : la pente y est la plus faible, ce qui traduit mathématiquement l'effet tampon. Comme $\beta$ est proportionnel à l'inverse de la pente ($\beta \propto \left|\dfrac{\mathrm{d}V_b}{\mathrm{d}\text{pH}}\right|$), le pouvoir tampon y est maximal.
\end{itemize}

\begin{aretenir}[L'essentiel de la section]
\begin{itemize}
\item À la demi-équivalence ($V_b = V_e/2$) : $[\ce{AH}] = [\ce{A-}]$ et pH $=$ p$K_a$ — lecture graphique directe du p$K_a$.
\item La zone tampon est la portion quasi horizontale de la courbe autour de la demi-équivalence (pH compris entre p$K_a - 1$ et p$K_a + 1$).
\item Le pouvoir tampon $\beta = \left|\dfrac{\mathrm{d}n}{\mathrm{d}\text{pH}}\right|$ est maximal à pH $=$ p$K_a$ et croît avec la concentration totale du couple.
\item La dérivée $\dfrac{\mathrm{d}\text{pH}}{\mathrm{d}V_b}$ est minimale à la demi-équivalence (tampon) et maximale à l'équivalence (saut de pH).
\end{itemize}
\end{aretenir}

\begin{exoresolu}[Exploiter une courbe de dosage]
On dose $V_a = \SI{25,0}{mL}$ d'une solution d'acide propanoïque \ce{C2H5COOH} de concentration $C_a = \SI{0,080}{mol.L^{-1}}$ par une solution de soude de concentration $C_b = \SI{0,100}{mol.L^{-1}}$. Le pH relevé à la demi-équivalence vaut $4,9$.
\begin{enumerate}
\item Calculer le volume équivalent $V_e$.
\item En déduire le volume pour lequel l'effet tampon est maximal, et le p$K_a$ du couple.
\item Pour quel volume $V_b$ a-t-on $[\ce{C2H5COO-}] = 3\,[\ce{C2H5COOH}]$ ? Quel est alors le pH ?
\end{enumerate}
\tcblower
\textbf{1. Volume équivalent.} À l'équivalence, $C_a V_a = C_b V_e$ :
\[ V_e = \frac{C_a V_a}{C_b} = \frac{0,080 \times 25,0}{0,100} = \SI{20,0}{mL} \]

\textbf{2. Effet tampon maximal.} Le pouvoir tampon est maximal à la demi-équivalence :
\[ V_b = \frac{V_e}{2} = \SI{10,0}{mL} \]
À ce volume, $[\ce{C2H5COOH}] = [\ce{C2H5COO-}]$, donc pH $=$ p$K_a$. Le relevé donne p$K_a(\ce{C2H5COOH}/\ce{C2H5COO-}) = 4,9$.

\textbf{3. Rapport base/acide égal à 3.} La réaction \ce{C2H5COOH + HO- -> C2H5COO- + H2O} étant quasi totale, un tableau d'avancement (en notant $n_0 = C_a V_a = \SI{2,0e-3}{mol}$ et $n = C_b V_b$ la quantité de soude versée) donne, avant l'équivalence :
\[ n(\ce{C2H5COOH}) = n_0 - n \qquad n(\ce{C2H5COO-}) = n \]
La condition $[\ce{C2H5COO-}] = 3\,[\ce{C2H5COOH}]$ s'écrit $n = 3(n_0 - n)$, soit $4n = 3n_0$, d'où $n = 0,75\,n_0$ et :
\[ V_b = 0,75\,V_e = \SI{15,0}{mL} \]
Le pH vaut alors :
\[ \text{pH} = \text{p}K_a + \log\frac{[\ce{C2H5COO-}]}{[\ce{C2H5COOH}]} = 4,9 + \log 3 \approx 5,4 \]
On vérifie que ce point ($V_b = \SI{15}{mL}$, pH $= 5,4$) se situe bien dans la zone tampon, comprise entre p$K_a - 1 = 3,9$ et p$K_a + 1 = 5,9$.
\end{exoresolu}

\begin{exercice}[À toi de jouer]
On dose \SI{10,0}{mL} d'une solution d'acide méthanoïque \ce{HCOOH} par de la soude ; l'équivalence est obtenue pour $V_e = \SI{12,5}{mL}$.
\begin{enumerate}
\item Pour quel volume versé le mélange constitue-t-il la meilleure solution tampon possible ?
\item Le pH mesuré à ce volume vaut $3,8$. Donner le p$K_a$ du couple \ce{HCOOH}/\ce{HCOO-}.
\item Tracer l'allure de la courbe en indiquant la zone tampon, la demi-équivalence et l'équivalence.
\end{enumerate}
\end{exercice}

\begin{corrige}[Exercice : À toi de jouer]
\textbf{1.} L'effet tampon est maximal à la demi-équivalence : $V_b = V_e/2 = \SI{6,25}{mL}$. À ce volume, $[\ce{HCOOH}] = [\ce{HCOO-}]$.

\textbf{2.} À la demi-équivalence, pH $=$ p$K_a$, donc p$K_a(\ce{HCOOH}/\ce{HCOO-}) = 3,8$.

\textbf{3.} La courbe part d'un pH voisin de $2,4$ (acide faible), s'aplatit autour de $V_b = \SI{6,25}{mL}$ où pH $= 3,8$ (zone tampon, entre pH $2,8$ et $4,8$ environ), présente un saut de pH à $V_b = \SI{12,5}{mL}$ (équivalence, pH $> 7$), puis se stabilise vers le pH de la soude en excès.
\end{corrige}

\coursec{Applications des solutions tampons}

Dans la section précédente, nous avons repéré la \cle{zone tampon} sur une courbe de dosage et quantifié le \cle{pouvoir tampon} : autour de la demi-équivalence, le pH varie très peu malgré l'ajout d'acide ou de base. Cette propriété remarquable n'est pas une simple curiosité de laboratoire : elle est mise à profit chaque jour, dans ton propre organisme, dans les hôpitaux de Yaoundé et de Douala, dans les usines pharmaceutiques, les conserveries et même dans ta salle de bains. Cette dernière section te montre \emph{pourquoi} les solutions tampons sont devenues un outil indispensable de la chimie appliquée.

\courssub{Les tampons biologiques : l'exemple vital du sang}

\textbf{Un problème de survie}\par
Les réactions biochimiques de notre organisme (respiration cellulaire, digestion, contraction musculaire) produisent en permanence des espèces acides : dioxyde de carbone \ce{CO2}, acide lactique \ce{CH3-CHOH-COOH} lors d'un effort musculaire intense, corps cétoniques, etc. Or les enzymes, ces catalyseurs biologiques, ne fonctionnent correctement que dans une plage de pH très étroite. Le sang doit donc maintenir son pH entre \num{7,35} et \num{7,45}, c'est-à-dire quasiment constant, malgré cet apport continu d'acidité. C'est exactement la mission d'une solution tampon.

\begin{definition}[Le tampon bicarbonate du sang]
Le principal système tampon du sang est le couple \ce{H2CO3}/\ce{HCO3-} (acide carbonique / ion hydrogénocarbonate), de $\mathrm{p}K_a = 6,1$ (valeur apparente à \SI{37}{\celsius}). Un couple secondaire, \ce{H2PO4-}/\ce{HPO4^2-} (phosphates), ainsi que les protéines plasmatiques et l'hémoglobine, renforcent ce pouvoir tampon.
\end{definition}

Le sang contient donc simultanément l'acide faible \ce{H2CO3} et sa base conjuguée \ce{HCO3-} en quantités comparables : c'est la définition même d'une solution tampon. Si un acide pénètre dans le sang, il est consommé par la base :

\[ \ce{HCO3- + H3O+ -> H2CO3 + H2O} \]

Si une base pénètre dans le sang, elle est consommée par l'acide :

\[ \ce{H2CO3 + HO- -> HCO3- + H2O} \]

Dans les deux cas, le rapport $\dfrac{[\ce{HCO3-}]}{[\ce{H2CO3}]}$ varie peu, donc le pH varie peu d'après la relation de Henderson-Hasselbalch :

\[ \mathrm{pH} = \mathrm{p}K_a + \log\frac{[\ce{HCO3-}]}{[\ce{H2CO3}]} \]

\begin{exemplebox}[Pourquoi le pH du sang vaut-il 7,4 alors que le pKa vaut 6,1 ?]
Dans le sang artériel, les concentrations sont approximativement $[\ce{HCO3-}] \approx \SI{24}{mmol.L^{-1}}$ et $[\ce{H2CO3}] \approx \SI{1,2}{mmol.L^{-1}}$, soit un rapport $\dfrac{[\ce{HCO3-}]}{[\ce{H2CO3}]} \approx 20$. On obtient :
\[ \mathrm{pH} = 6,1 + \log(20) = 6,1 + 1,3 \approx 7,4 \]
Le pH sanguin est donc \emph{supérieur} au $\mathrm{p}K_a$ du couple : le sang fonctionne dans la partie haute de sa zone tampon ($\mathrm{p}K_a \pm 1$), ce qui lui permet de neutraliser préférentiellement les \emph{acides} (le danger le plus fréquent pour l'organisme).
\end{exemplebox}

\textbf{Le rôle respiratoire : les poumons, partenaires du tampon}\par
L'acide carbonique \ce{H2CO3} est en équilibre avec le dioxyde de carbone dissous, lui-même en équilibre avec le \ce{CO2} gazeux de l'air alvéolaire :

\[ \ce{CO2(g) <=> CO2(aq) + H2O <=> H2CO3 <=> HCO3- + H+} \]

Cette chaîne d'équilibres donne à l'organisme un moyen d'action supplémentaire : en modifiant la ventilation, les poumons règlent la concentration en \ce{H2CO3}.

\begin{itemize}
\item Si le sang s'acidifie (excès de \ce{H+}), l'équilibre $\ce{H2CO3 <=> HCO3- + H+}$ est déplacé vers la gauche ; l'excès de \ce{H2CO3} se transforme en \ce{CO2} que la respiration \emph{accélérée} élimine : le pH remonte.
\item Si le sang devient trop basique, la ventilation ralentit, le \ce{CO2} s'accumule, forme \ce{H2CO3} qui libère des ions \ce{H+} : le pH redescend.
\end{itemize}

\begin{popfigure}
\begin{center}
\begin{tikzpicture}[font=\small, >=Stealth]
% Poumons
\draw[fill=poppinkL, draw=poppink, thick] (0,1.6) ellipse (1.5 and 0.9);
\node[popdark] at (0,1.6) {\textbf{POUMONS}};
\node[popmuted] at (0,1.15) {ventilation};
\node[popdark] at (0,-0.2) {$\ce{CO2(g)}$ éliminé};
\draw[->, thick, poppink] (0,0.75) -- (0,0.35);
% Chaîne d'équilibres
\node[draw=popblue, fill=popblueL, rounded corners, thick, minimum width=2.6cm, minimum height=0.9cm] (co2) at (3.2,0.4) {$\ce{CO2} + \ce{H2O}$};
\node[draw=popgreen, fill=popgreenL, rounded corners, thick, minimum width=2.6cm, minimum height=0.9cm] (h2co3) at (6.9,0.4) {$\ce{H2CO3}$};
\node[draw=poporange, fill=poporangeL, rounded corners, thick, minimum width=3.2cm, minimum height=0.9cm] (hco3) at (10.9,0.4) {$\ce{HCO3-} + \ce{H+}$};
\draw[<->, very thick, popblue] (co2) -- (h2co3);
\draw[<->, very thick, popgreen] (h2co3) -- (hco3);
\draw[<->, thick, poppink] (0.9,0.4) -- (co2);
% Sang
\draw[fill=poppink!25, draw=poppink, thick, rounded corners] (2.2,-2.3) rectangle (12.4,-1.2);
\node[popdark] at (7.3,-1.75) {\textbf{SANG} : pH maintenu entre \num{7,35} et \num{7,45} grâce au tampon \ce{H2CO3}/\ce{HCO3-}};
\draw[->, thick, poppurple] (6.9,-0.15) -- (6.9,-1.1);
\draw[->, thick, poppurple] (10.9,-0.15) -- (10.9,-1.1);
% Annotations
\node[align=center, popmuted] at (5.05,1.15) {hydratation /\\déshydratation};
\node[align=center, popmuted] at (9.0,1.15) {ionisation\\($\mathrm{p}K_a = 6,1$)};
\end{tikzpicture}
\end{center}
\legende{Le système tampon du sang : la chaîne d'équilibres reliant le \ce{CO2} éliminé par les poumons au couple \ce{H2CO3}/\ce{HCO3-} permet de réguler finement le pH sanguin.}
\end{popfigure}

\begin{attention}[Acidose et alcalose : quand le tampon est débordé]
Le tampon sanguin a un pouvoir tampon \emph{fini}. Si l'apport d'acide ou de base dépasse sa capacité, le pH sort de la plage \num{7,35}--\num{7,45} :
\begin{itemize}
\item \cle{acidose} ($\mathrm{pH} < 7,35$) : diabète non traité (cétose), effort musculaire extrême (acide lactique), insuffisance respiratoire (rétention de \ce{CO2}) ;
\item \cle{alcalose} ($\mathrm{pH} > 7,45$) : hyperventilation prolongée (élimination excessive de \ce{CO2}), vomissements répétés (perte d'acide chlorhydrique gastrique).
\end{itemize}
Ces deux états sont des urgences médicales. Ne confonds pas : l'acidose n'est pas une « acidité normale », c'est la \emph{rupture} de la régulation tampon.
\end{attention}

\begin{savaistu}[0,4 unité de pH : la frontière entre la vie et la mort]
Une variation du pH sanguin de seulement \num{0,4} unité (par exemple de 7,4 à 7,0 ou à 7,8) peut entraîner le coma, voire la mort. En dessous de 6,8 ou au-dessus de 7,8, la survie est impossible. C'est dire l'importance \emph{vitale} du tampon bicarbonate : chaque seconde, il neutralise l'acidité produite par ton métabolisme sans que tu t'en rendes compte. Les sportifs de haut niveau exploitent d'ailleurs ce système : l'acide lactique produit pendant un sprint est d'abord neutralisé par \ce{HCO3-}, puis le \ce{CO2} formé est évacué par l'essoufflement qui suit l'effort.
\end{savaistu}

\courssub{L'étalonnage des pH-mètres}

Un pH-mètre ne donne une mesure fiable que s'il a été \cle{étalonné} au préalable : la tension électrique délivrée par l'électrode de verre doit être convertie en pH grâce à une relation qui dépend de l'état de l'électrode (vieillissement, température). L'étalonnage consiste à plonger l'électrode dans des \cle{solutions tampons étalons}, dont le pH est connu avec une grande précision et, surtout, \emph{stable dans le temps} — c'est là que la propriété tampon est exploitée : le pH de la solution étalon ne bouge pratiquement pas malgré la dilution due aux traces d'eau de rinçage ou la contamination par le \ce{CO2} de l'air.

\begin{methode}[Étalonner un pH-mètre (méthode des deux points)]
\begin{enumerate}
\item Rincer l'électrode à l'eau distillée et la sécher délicatement (sans frotter le bulbe).
\item Plonger l'électrode dans la solution tampon \textbf{pH = 7,0} (étalon neutre) ; régler le bouton « \emph{offset} » (ou « pH 7 ») jusqu'à affichage de 7,00.
\item Rincer à nouveau, puis plonger dans une seconde solution tampon : \textbf{pH = 4,0} si on mesurera des milieux acides, \textbf{pH = 10,0} (ou 9,2) si on mesurera des milieux basiques ; régler la « pente » (slope) jusqu'à affichage correct.
\item Rincer une dernière fois : l'appareil est prêt pour les mesures.
\end{enumerate}
\end{methode}

\begin{attention}[Pourquoi ne pas étalonner avec de l'eau distillée ou une solution d'acide fort ?]
L'eau distillée n'est \emph{pas} une solution tampon : son pH, proche de 5,7 à cause du \ce{CO2} dissous, varie énormément à la moindre contamination. Une solution d'acide fort (ex. \ce{HCl} \SI{0,01}{mol.L^{-1}}) a un pH théorique connu (2,00), mais elle n'est pas tamponnée : son pH change dès qu'une trace d'eau ou de \ce{CO2} la dilue ou la modifie. Une solution d'étalonnage doit avoir un pH \emph{connu avec certitude} et \emph{stable} malgré ces perturbations : seule une solution tampon certifiée garantit les deux. C'est une question classique au Baccalauréat.
\end{attention}

\courssub{L'industrie pharmaceutique}

Les médicaments administrés par voie injectable (perfusions, vaccins, sérums) pénètrent directement dans le sang ou les tissus : leur pH doit être proche du pH physiologique (\num{7,4}) pour éviter douleurs, brûlures ou destruction des tissus (hémolyse). On y incorpore donc des tampons (phosphate, citrate, bicarbonate) qui maintiennent le pH du produit constant pendant toute sa durée de conservation. De même, les \cle{collyres} (gouttes pour les yeux) sont tamponnés au pH des larmes ($\approx 7,4$) : une simple goutte de solution non tamponnée, même faiblement acide, provoquerait une vive irritation. Les collyres contiennent souvent un tampon phosphate ou borique. Enfin, de nombreux principes actifs ne sont stables ou solubles que dans une étroite plage de pH : le tampon garantit l'efficacité du médicament jusqu'à sa date de péremption.

\courssub{L'agroalimentaire}

Dans l'industrie alimentaire, le pH est un paramètre de \emph{conservation} et de \emph{goût} :

\begin{itemize}
\item \textbf{Conserves} : un pH inférieur à 4,5 empêche le développement de bactéries dangereuses (notamment \emph{Clostridium botulinum}, agent du botulisme). L'acidité des conserves de tomates, de haricots ou de fruits est contrôlée et stabilisée par des tampons (citrate, lactate).
\item \textbf{Boissons} : les sodas et jus de fruits industriels (jus de bissap, jus de gingembre embouteillés) contiennent de l'acide citrique et du citrate de sodium — un couple tampon — qui fixe un pH stable, garant d'un goût constant d'un lot à l'autre et d'une saveur équilibrée (ni trop acide, ni plat).
\item \textbf{Produits laitiers} : la fermentation du lait en yaourt fait chuter le pH ; maîtriser le pouvoir tampon du lait (dû aux phosphates \ce{H2PO4-}/\ce{HPO4^2-} et aux protéines comme la caséine) permet de contrôler la cinétique de fermentation et la texture finale (ferme ou brassé).
\end{itemize}

\courssub{Cosmétiques et détergents}

La peau et les cheveux ont un pH naturellement légèrement acide, autour de \num{5,5} (le « film hydrolipidique » ou manteau acide protège contre les bactéries et les agressions extérieures). Les shampoings, gels douche et crèmes sont donc tamponnés pour respecter ce pH et éviter irritations, dessèchement ou perturbation de la flore cutanée. Les savons classiques, fabriqués par saponification (comme dans les savonneries artisanales à partir d'huile de palme au Cameroun), sont naturellement basiques (pH $\approx 9$ à $10$) du fait de l'hydrolyse du carbonate formé ; les « pains dermatologiques » ou « savons sans savon » (syndets) modernes sont tamponnés pour se rapprocher du pH cutané.

\begin{attention}[Le piège du « pH neutre »]
Sur un shampoing ou un gel douche, la mention « \emph{pH neutre} » signifie un pH \emph{proche de celui de la peau}, soit environ \num{5,5}, et \textbf{pas} $\mathrm{pH} = 7$ ! Il s'agit d'un usage commercial du mot « neutre » (synonyme de « non agressif »). En chimie, la neutralité à \SI{25}{\celsius} reste rigoureusement $\mathrm{pH} = 7$ ($[\ce{H3O+}] = [\ce{HO-}]$). Au Baccalauréat, ne laisse jamais passer cette confusion : un produit « pH neutre » cosmétique est \emph{acide} au sens chimique.
\end{attention}

\begin{popfigure}
\begin{center}
\begin{tikzpicture}[font=\small, >=Stealth]
\draw[->, thick, popline] (0,0) -- (13.6,0);
\foreach \x/\titre/\icone/\couleur/\teinte in {
  1.3/{Sang}/{pH $\approx 7,4$\\Couple \ce{H2CO3}/\ce{HCO3-}}/poppink/poppinkL,
  4.0/{pH-mètre}/{Étalons pH 4, 7, 10}/popblue/popblueL,
  6.7/{Pharmacie}/{Injectables, collyres\\Stabilité principes actifs}/popgreen/popgreenL,
  9.4/{Agroalimentaire}/{Conserves, boissons,\\laitages}/poporange/poporangeL,
  12.1/{Cosmétiques}/{Shampoings, syndets\\pH $\approx 5,5$}/poppurple/poppurpleL}{
  \draw[fill=\teinte, draw=\couleur, thick] (\x,0) circle (0.14);
  \node[draw=\couleur, fill=\teinte, rounded corners, thick, align=center, minimum width=2.4cm] at (\x,1.3) {\textbf{\titre}\\ \icone};
  \draw[thick, \couleur] (\x,0.14) -- (\x,0.7);
}
\node[popmuted, align=center] at (6.8,-0.6) {Les solutions tampons : un même principe chimique, des applications partout};
\end{tikzpicture}
\end{center}
\legende{Frise des grandes applications des solutions tampons : du vivant à l'industrie, la même propriété — un pH stable — est exploitée.}
\end{popfigure}

\begin{exoresolu}[Choisir un tampon pour un collyre]
Un laboratoire pharmaceutique veut préparer un collyre de $\mathrm{pH} = 7,4$. Il dispose des couples suivants : \ce{CH3COOH}/\ce{CH3COO-} ($\mathrm{p}K_a = 4,8$), \ce{H2PO4-}/\ce{HPO4^2-} ($\mathrm{p}K_a = 7,2$), \ce{NH4+}/\ce{NH3} ($\mathrm{p}K_a = 9,2$).
\begin{enumerate}
\item Quel couple choisir ? Justifier.
\item Calculer le rapport $\dfrac{[\text{base}]}{[\text{acide}]}$ nécessaire.
\end{enumerate}
\tcblower
\textbf{1. Choix du couple.} Une solution tampon est d'autant plus efficace (pouvoir tampon maximal) que le pH visé est proche du $\mathrm{p}K_a$ du couple (zone tampon efficace : $\mathrm{p}K_a - 1 \leqslant \mathrm{pH} \leqslant \mathrm{p}K_a + 1$). Pour $\mathrm{pH} = 7,4$ :
\begin{itemize}
\item $\mathrm{p}K_a = 4,8$ (acide acétique) : $7,4 \notin [3,8\,;\,5,8]$ : pouvoir tampon nul $\rightarrow$ rejeté ;
\item $\mathrm{p}K_a = 7,2$ (phosphate) : $7,4 \in [6,2\,;\,8,2]$ : \textbf{retenu} (pouvoir tampon optimal) ;
\item $\mathrm{p}K_a = 9,2$ (ammonium) : $7,4 \notin [8,2\,;\,10,2]$ : pouvoir tampon nul $\rightarrow$ rejeté.
\end{itemize}
On choisit le couple phosphate \ce{H2PO4-}/\ce{HPO4^2-} (c'est d'ailleurs l'un des tampons réellement utilisés dans les collyres et les perfusions).

\textbf{2. Rapport des concentrations.} D'après la relation de Henderson-Hasselbalch :
\[ \mathrm{pH} = \mathrm{p}K_a + \log\frac{[\ce{HPO4^2-}]}{[\ce{H2PO4-}]} \implies \log\frac{[\ce{HPO4^2-}]}{[\ce{H2PO4-}]} = 7,4 - 7,2 = 0,2 \]
\[ \frac{[\ce{HPO4^2-}]}{[\ce{H2PO4-}]} = 10^{0,2} \approx 1,58 \approx 1,6 \]
Il faut donc environ \num{1,6} fois plus de base conjuguée \ce{HPO4^2-} que d'acide \ce{H2PO4-} (par exemple \SI{0,16}{mol} de \ce{Na2HPO4} pour \SI{0,10}{mol} de \ce{NaH2PO4} dans un litre de solution).
\end{exoresolu}

\begin{exercice}[Pour t'entraîner]
\begin{enumerate}
\item Expliquer, à l'aide des équations d'équilibre, pourquoi une hyperventilation (respiration rapide et profonde) fait augmenter le pH du sang.
\item Un technicien de laboratoire veut vérifier le bon fonctionnement d'un pH-mètre. Pourquoi utilise-t-il des solutions tampons pH 4 et pH 7 plutôt qu'une solution d'acide chlorhydrique de concentration connue ?
\item Un fabricant de shampoing annonce « pH neutre ». Le pH mesuré du produit est 5,6. L'annonce est-elle mensongère ? Justifier ta réponse avec le vocabulaire du cours.
\end{enumerate}
\end{exercice}

\begin{corrige}[Exercice]
\begin{enumerate}
\item L'hyperventilation élimine massivement le \ce{CO2} gazeux. Par loi de Le Chatelier, l'équilibre $\ce{CO2(g) <=> CO2(aq)}$ est déplacé vers la gauche (consommation de \ce{CO2(aq)}), ce qui entraîne le déplacement vers la gauche de $\ce{CO2(aq) + H2O <=> H2CO3}$ (consommation de \ce{H2CO3}), puis de $\ce{H2CO3 <=> HCO3- + H+}$ : la concentration en ions \ce{H+} diminue, donc le pH du sang augmente (risque d'alcalose respiratoire).
\item Une solution de \ce{HCl} n'est pas tamponnée : son pH peut être modifié par la dilution (traces d'eau sur l'électrode) ou par absorption du \ce{CO2} atmosphérique. Une solution tampon pH 4 ou 7 possède un pH \emph{connu avec précision} (valeur certifiée) et \emph{stable} malgré ces perturbations : c'est la condition indispensable pour un étalonnage fiable (offset et pente corrects).
\item Non, l'annonce n'est pas mensongère au sens commercial : dans l'industrie cosmétique, « pH neutre » désigne un pH proche de celui de la peau ($\approx 5,5$), ce qui est le cas ici (5,6). En revanche, au sens \emph{chimique} strict, la neutralité correspond à $\mathrm{pH} = 7$ à \SI{25}{\celsius} : le shampoing est en réalité légèrement acide, ce qui est voulu pour respecter le film hydrolipidique protecteur de la peau (manteau acide).
\end{enumerate}
\end{corrige}

\begin{aretenir}[L'essentiel de la section]
\begin{itemize}
\item Le sang est tamponné par le couple \ce{H2CO3}/\ce{HCO3-} ($\mathrm{p}K_a = 6,1$ à \SI{37}{\celsius}) : avec un rapport $\dfrac{[\ce{HCO3-}]}{[\ce{H2CO3}]} \approx 20$, on obtient $\mathrm{pH} \approx 7,4$.
\item Les poumons régulent le pH sanguin en éliminant plus ou moins de \ce{CO2}, ce qui déplace les équilibres acido-basiques (régulation respiratoire).
\item L'étalonnage des pH-mètres exige des solutions tampons étalons (pH 4, 7, 10), seules à offrir un pH connu avec certitude et stable.
\item Pharmacie (injectables, collyres), agroalimentaire (conserves, boissons, laitages) et cosmétiques (shampoings, syndets) exploitent les tampons pour stabiliser le pH de leurs produits.
\item Piège : « pH neutre » d'un shampoing $\approx 5,5$ (pH de la peau), pas 7 !
\end{itemize}
\end{aretenir}

\begin{savaistu}[Lien avec le Baccalauréat]
Les questions appliquant les solutions tampons au \emph{sang} (calcul de pH à partir du rapport $[\ce{HCO3-}]/[\ce{H2CO3}]$, rôle des poumons, acidose/alcalose) et à l'\emph{étalonnage des pH-mètres} (pourquoi des solutions tampons ? quelles valeurs d'étalonnage ?) reviennent très fréquemment dans les sujets de Terminale C et D au Cameroun. Maîtrise la relation $\mathrm{pH} = \mathrm{p}K_a + \log\frac{[\text{base}]}{[\text{acide}]}$ et sache rédiger la justification du choix d'un couple tampon (proximité pH/pKa) : ce sont des points faciles à gagner le jour de l'examen.
\end{savaistu}

\newpage

\coursec{Travaux pratiques}

\courssub{Objectifs du TP}

À l'issue de cette séance de travaux pratiques, tu dois être capable de :

\begin{itemize}
\item préparer une solution tampon éthanoïque de $\mathrm{pH} = \mathrm{p}K_a(\ce{CH3COOH}/\ce{CH3COO-}) = 4{,}8$ par mélange équimolaire d'acide éthanoïque et d'éthanoate de sodium ;
\item mesurer le pH d'une solution à l'aide d'un pH-mètre étalonné (ou, à défaut, de papier pH de précision) ;
\item vérifier expérimentalement la propriété fondamentale d'une solution tampon : son pH varie \cle{peu} lors de l'ajout modéré d'un acide fort, d'une base forte ou par dilution ;
\item comparer le comportement du tampon à celui de l'eau distillée soumise aux mêmes ajouts ;
\item exploiter la relation de Henderson-Hasselbalch pour calculer les pH théoriques et les confronter aux mesures ;
\item faire le lien entre cette préparation (mélange équimolaire) et la demi-équivalence d'un dosage acide-base.
\end{itemize}

\courssub{Matériel et produits}

\begin{center}
\begin{tabularx}{\linewidth}{|l|X|}
\hline
\rowcolor{popblueL} \textbf{Produits} & \textbf{Matériel} \\
\hline
Solution d'acide éthanoïque \ce{CH3COOH} à \SI{0,10}{mol.L^{-1}} & 5 béchers de \SI{150}{mL} propres et secs \\
\hline
Solution d'éthanoate de sodium \ce{CH3COONa} à \SI{0,10}{mol.L^{-1}} & 2 pipettes jaugées de \SI{50}{mL} + poire aspirante \\
\hline
Solution d'acide chlorhydrique \ce{HCl} à \SI{0,10}{mol.L^{-1}} & 1 pipette jaugée de \SI{1}{mL} (ou pipette graduée) \\
\hline
Solution d'hydroxyde de sodium \ce{NaOH} à \SI{0,10}{mol.L^{-1}} & 1 fiole jaugée de \SI{100}{mL} (facultatif) \\
\hline
Eau distillée & pH-mètre étalonné (solutions étalons pH 4 et pH 7) \\
\hline
Solutions tampons étalons pH 4 et pH 7 & Agitateur en verre, éprouvette de \SI{100}{mL}, pissette \\
\hline
\end{tabularx}
\end{center}

\begin{attention}[Si un produit manque dans ton laboratoire]
\begin{itemize}
\item \textbf{Pas d'éthanoate de sodium ?} Prépare le tampon par \cle{demi-neutralisation} : verse \SI{100}{mL} d'acide éthanoïque à \SI{0,10}{mol.L^{-1}} et ajoute \SI{50}{mL} de soude à \SI{0,10}{mol.L^{-1}}. La moitié de l'acide est transformée en ions éthanoate : tu obtiens bien un mélange équimolaire \ce{CH3COOH}/\ce{CH3COO-}.
\item \textbf{Pas de pH-mètre ?} Utilise du papier pH de précision (gamme 3--6, puis gamme large 1--14 pour l'eau distillée). Les mesures seront moins fines mais la conclusion restera spectaculaire.
\item \textbf{Pas de solutions étalons commerciales ?} Le jus de citron filtré (pH $\approx 2{,}3$) et une solution de bicarbonate de soude (pH $\approx 8{,}3$) peuvent servir de repères grossiers pour vérifier qu'un pH-mètre n'est pas déréglé.
\end{itemize}
\end{attention}

\courssub{Sécurité}

\begin{attention}[Consignes de sécurité]
Les solutions utilisées sont diluées, mais la prudence reste de mise :
\begin{itemize}
\item \textbf{Pictogramme « corrosif »} (gouttes attaquant une main et une surface) : l'acide chlorhydrique et la soude, même à \SI{0,10}{mol.L^{-1}}, sont irritants pour la peau et les yeux. Porte la blouse, des lunettes de protection et, si possible, des gants.
\item \textbf{Jamais de pipetage à la bouche} : utilise toujours la poire aspirante (propipette).
\item L'acide éthanoïque dégage des vapeurs âcres (odeur de vinaigre) : travaille dans un espace aéré, ne porte pas le bécher au nez.
\item En cas de contact avec la peau, rince abondamment à l'eau claire pendant plusieurs minutes et préviens l'enseignant.
\item Ne jette les solutions dans l'évier qu'après accord de l'enseignant, en faisant couler l'eau.
\end{itemize}
\end{attention}

\courssub{Schéma du dispositif}

\begin{popfigure}
\begin{center}
\begin{tikzpicture}[scale=0.92, every node/.style={font=\small}]
% Bécher tampon (gauche)
\draw[thick, popblue] (0,0) -- (0,2.4) (2.2,0) -- (2.2,2.4);
\draw[thick, popblue] (0,0) -- (2.2,0);
\draw[popblue!50, fill=popblue!15] (0.05,0.05) rectangle (2.15,1.3);
\node[popdark] at (1.1,0.65) {tampon};
\node[popdark, align=center] at (1.1,-0.5) {Bécher A$_1$\\ \SI{50}{mL} tampon};
% sonde pH
\draw[thick, popdark, fill=popcream] (1.5,1.1) rectangle (1.75,3.2);
\node[popdark, rotate=90] at (1.62,2.6) {\footnotesize sonde};
\draw[thick, popdark] (1.62,3.2) -- (1.62,3.7) -- (3.1,3.7);
\draw[thick, popdark, fill=poporangeL] (3.1,3.3) rectangle (4.9,4.1);
\node[popdark] at (4.0,3.7) {\footnotesize pH-mètre};
% pipette ajout HCl
\draw[thick, poppurple] (-0.9,3.6) -- (-0.4,2.2);
\draw[thick, poppurple, -{Stealth}] (-0.4,2.2) -- (-0.2,1.5);
\node[poppurple, align=center] at (-1.6,3.9) {\footnotesize \SI{1}{mL} \ce{HCl}\\ \footnotesize \SI{0,10}{mol.L^{-1}}};

% Bécher eau (droite)
\begin{scope}[xshift=7.2cm]
\draw[thick, popgreen] (0,0) -- (0,2.4) (2.2,0) -- (2.2,2.4);
\draw[thick, popgreen] (0,0) -- (2.2,0);
\draw[popgreen!60, fill=popgreen!12] (0.05,0.05) rectangle (2.15,1.3);
\node[popdark] at (1.1,0.65) {eau distillée};
\node[popdark, align=center] at (1.1,-0.5) {Bécher B$_1$\\ \SI{100}{mL} eau};
\draw[thick, popdark, fill=popcream] (1.5,1.1) rectangle (1.75,3.2);
\draw[thick, popdark] (1.62,3.2) -- (1.62,3.7) -- (3.1,3.7);
\draw[thick, popdark, fill=poporangeL] (3.1,3.3) rectangle (4.9,4.1);
\node[popdark] at (4.0,3.7) {\footnotesize pH-mètre};
\draw[thick, poppurple] (-0.9,3.6) -- (-0.4,2.2);
\draw[thick, poppurple, -{Stealth}] (-0.4,2.2) -- (-0.2,1.5);
\node[poppurple, align=center] at (-1.6,3.9) {\footnotesize \SI{1}{mL} \ce{HCl}\\ \footnotesize \SI{0,10}{mol.L^{-1}}};
\end{scope}
\end{tikzpicture}
\end{center}
\legende{Dispositif de vérification de l'effet tampon : on ajoute le même volume d'acide chlorhydrique dans le tampon éthanoïque (bécher A$_1$) et dans l'eau distillée (bécher B$_1$), puis on compare les variations de pH lues au pH-mètre.}
\end{popfigure}

\courssub{Protocole}

\begin{methode}[Préparation et vérification du tampon]
\begin{enumerate}
\item \textbf{Étalonnage.} Étale le pH-mètre avec les solutions étalons pH 7 puis pH 4. Rince la sonde à l'eau distillée et essuie-la délicatement entre deux mesures.
\item \textbf{Préparation du tampon.} À l'aide des deux pipettes jaugées de \SI{50}{mL}, introduis dans un bécher propre \SI{50,0}{mL} de solution d'acide éthanoïque à \SI{0,10}{mol.L^{-1}} et \SI{50,0}{mL} de solution d'éthanoate de sodium à \SI{0,10}{mol.L^{-1}}. Homogénéise à l'agitateur. Tu obtiens \SI{100}{mL} de solution tampon.
\item \textbf{Première mesure.} Plonge la sonde dans le tampon, attends la stabilisation et note le pH. Compare à la valeur attendue $\mathrm{pH} = 4{,}8$.
\item \textbf{Partage.} Répartis le tampon en deux béchers A$_1$ et A$_2$ contenant chacun environ \SI{50}{mL}.
\item \textbf{Ajout d'acide.} Dans A$_1$, ajoute \SI{1,0}{mL} d'acide chlorhydrique à \SI{0,10}{mol.L^{-1}}, homogénéise, mesure le pH.
\item \textbf{Ajout de base.} Dans A$_2$, ajoute \SI{1,0}{mL} de soude à \SI{0,10}{mol.L^{-1}}, homogénéise, mesure le pH.
\item \textbf{Témoin eau distillée.} Verse \SI{100}{mL} d'eau distillée dans deux béchers B$_1$ et B$_2$. Mesure le pH initial, puis ajoute \SI{1,0}{mL} de \ce{HCl} \SI{0,10}{mol.L^{-1}} dans B$_1$ et \SI{1,0}{mL} de \ce{NaOH} \SI{0,10}{mol.L^{-1}} dans B$_2$. Mesure les pH après chaque ajout.
\item \textbf{Test de dilution (facultatif).} Prélève \SI{10}{mL} du tampon restant, dilue avec \SI{90}{mL} d'eau distillée et mesure le pH.
\end{enumerate}
\end{methode}

\courssub{Observations et résultats attendus}

Voici un tableau de mesures type, obtenues avec un pH-mètre correctement étalonné à \SI{25}{\celsius} :

\begin{center}
\begin{tabularx}{\linewidth}{|X|c|c|c|}
\hline
\rowcolor{popblueL} \textbf{Solution} & \textbf{pH initial} & \textbf{pH après ajout} & \textbf{$\Delta$pH} \\
\hline
Tampon + \SI{1}{mL} \ce{HCl} \SI{0,10}{mol.L^{-1}} & 4{,}8 & 4{,}7 & $-0{,}1$ \\
\hline
Tampon + \SI{1}{mL} \ce{NaOH} \SI{0,10}{mol.L^{-1}} & 4{,}8 & 4{,}9 & $+0{,}1$ \\
\hline
Tampon dilué 10 fois & 4{,}8 & 4{,}8 & $\approx 0$ \\
\hline
Eau distillée + \SI{1}{mL} \ce{HCl} \SI{0,10}{mol.L^{-1}} & 7{,}0 & 3{,}0 & $-4{,}0$ \\
\hline
Eau distillée + \SI{1}{mL} \ce{NaOH} \SI{0,10}{mol.L^{-1}} & 7{,}0 & 11{,}0 & $+4{,}0$ \\
\hline
\end{tabularx}
\end{center}

Le contraste est frappant : la même quantité d'acide fait chuter le pH de l'eau de 4 unités, alors qu'il ne fait varier celui du tampon que d'environ un dixième d'unité.

\begin{attention}[Écarts possibles entre mesures et valeurs attendues]
\begin{itemize}
\item Un pH initial du tampon compris entre 4{,}6 et 5{,}0 est tout à fait acceptable : un pH-mètre mal étalonné, une température différente de \SI{25}{\celsius} ou des concentrations légèrement fausses expliquent l'écart.
\item L'eau distillée exposée à l'air dissout le dioxyde de carbone : son pH initial est souvent voisin de 5{,}5--6{,}5 et non 7{,}0. Ce n'est pas une erreur de manipulation !
\item Rince toujours la sonde entre deux mesures pour éviter les contaminations croisées.
\end{itemize}
\end{attention}

\courssub{Exploitation}

\begin{exoresolu}[Questions guidées d'exploitation]
\textbf{Question 1.} Montrer que le mélange préparé contient des quantités égales d'acide éthanoïque et d'ions éthanoate, puis calculer son pH théorique. On donne $\mathrm{p}K_a(\ce{CH3COOH}/\ce{CH3COO-}) = 4{,}8$.

\textbf{Question 2.} Calculer le pH théorique du tampon après ajout de \SI{1,0}{mL} de \ce{HCl} à \SI{0,10}{mol.L^{-1}} dans \SI{50}{mL} de tampon. Comparer à la mesure.

\textbf{Question 3.} Calculer le pH de l'eau distillée après ajout de \SI{1,0}{mL} de \ce{HCl} à \SI{0,10}{mol.L^{-1}} dans \SI{100}{mL}. Conclure.

\textbf{Question 4.} Expliquer, à l'aide des réactions mises en jeu, pourquoi le pH du tampon varie peu.

\textbf{Question 5.} Relier la préparation effectuée (mélange équimolaire) à la notion de demi-équivalence lors d'un dosage de l'acide éthanoïque par la soude.
\tcblower
\textbf{Question 1.} Quantités introduites :
\[ n(\ce{CH3COOH}) = \SI{0,10}{mol.L^{-1}} \times \SI{50,0e-3}{L} = \SI{5,0e-3}{mol} \]
\[ n(\ce{CH3COO-}) = \SI{0,10}{mol.L^{-1}} \times \SI{50,0e-3}{L} = \SI{5,0e-3}{mol} \]
Le mélange est bien équimolaire. La relation de Henderson-Hasselbalch donne :
\[ \mathrm{pH} = \mathrm{p}K_a + \log\frac{[\ce{CH3COO-}]}{[\ce{CH3COOH}]} = 4{,}8 + \log 1 = 4{,}8 \]

\textbf{Question 2.} Les ions \ce{H3O+} apportés réagissent quantitativement sur la base du couple :
\[ \ce{CH3COO- + H3O+ -> CH3COOH + H2O} \]
Quantité d'acide ajoutée : $n(\ce{H3O+}) = 0{,}10 \times 1{,}0\times10^{-3} = \SI{1,0e-4}{mol}$.

Dans les \SI{50}{mL} de tampon : $n(\ce{CH3COOH}) = n(\ce{CH3COO-}) = \SI{2,5e-3}{mol}$.

Après réaction : $n(\ce{CH3COO-}) = 2{,}5\times10^{-3} - 1{,}0\times10^{-4} = \SI{2,4e-3}{mol}$ et $n(\ce{CH3COOH}) = 2{,}5\times10^{-3} + 1{,}0\times10^{-4} = \SI{2,6e-3}{mol}$.
\[ \mathrm{pH} = 4{,}8 + \log\frac{2{,}4}{2{,}6} = 4{,}8 - 0{,}03 \approx 4{,}77 \]
La variation théorique ($-0{,}03$) est encore plus faible que la variation mesurée ($-0{,}1$) : l'accord est excellent compte tenu de la précision du pH-mètre.

\textbf{Question 3.} Dans l'eau, l'acide fort est totalement dissocié :
\[ [\ce{H3O+}] = \frac{\SI{1,0e-4}{mol}}{\SI{101e-3}{L}} \approx \SI{9,9e-4}{mol.L^{-1}} \]
\[ \mathrm{pH} = -\log(9{,}9\times10^{-4}) \approx 3{,}0 \]
Le pH chute de 7 à 3, soit une variation de 4 unités, contre moins de 0{,}1 pour le tampon : l'effet tampon est mis en évidence.

\textbf{Question 4.} Le tampon contient à la fois un acide faible et sa base conjuguée en quantités comparables. L'acide ajouté est « consommé » par les ions \ce{CH3COO-} ; une base ajoutée serait consommée par \ce{CH3COOH} selon $\ce{CH3COOH + HO- -> CH3COO- + H2O}$. Les espèces agressives sont donc neutralisées et le rapport $\dfrac{[\ce{CH3COO-}]}{[\ce{CH3COOH}]}$ ne varie que très peu : le pH reste quasi constant.

\textbf{Question 5.} Lors du dosage de l'acide éthanoïque par la soude, la demi-équivalence correspond au volume de soude pour lequel la moitié de l'acide initial a été neutralisée. On a alors $n(\ce{CH3COOH}) = n(\ce{CH3COO-})$. Le pH au point de demi-équivalence vaut exactement $\mathrm{p}K_a$. Notre préparation par mélange de \SI{50}{mL} d'acide \SI{0,10}{M} et \SI{50}{mL} de sel \SI{0,10}{M} reproduit exactement cette situation : c'est une réalisation « statique » de la demi-équivalence.
\end{exoresolu}

\begin{savaistu}[Zone tampon, pouvoir tampon et applications au Cameroun]
\begin{itemize}
\item \textbf{Zone tampon :} Sur une courbe de dosage pH = f(V), la zone tampon est l'intervalle de volumes autour de la demi-équivalence où le pH varie faiblement (généralement $\mathrm{p}K_a \pm 1$). C'est la région où les deux formes du couple coexistent en proportions significatives.
\item \textbf{Pouvoir tampon :} Il mesure la résistance au changement de pH. Il est maximal pour un mélange équimolaire (pH = pKa) et augmente avec la concentration totale du couple ($C_{\text{tot}} = [\ce{AH}] + [\ce{A-}]$). Diluer le tampon 10 fois ne change pas le pH (rapport constant) mais divise le pouvoir tampon par 10 : il « encaissera » moins d'acide/base avant de céder.
\item \textbf{Sang humain :} Le couple principal est \ce{H2CO3}/\ce{HCO3-} (pKa $\approx$ 6,1 à 37 °C). Le pH sanguin est maintenu à 7,40 $\pm$ 0,05 grâce à la régulation respiratoire (\ce{CO2}) et rénale (\ce{HCO3-}). Une acidose (pH < 7,35) ou alcalose (pH > 7,45) met en jeu le pronostic vital.
\item \textbf{Industrie agroalimentaire (Cameroun) :} La conservation du \emph{bili-bili} (bière de mil), du vin de palme ou des jus (bissap, ananas) repose sur le contrôle du pH (souvent acide, pH 3--4) pour inhiber les bactéries de contamination. Les tampons alimentaires (citrates, phosphates) stabilisent la saveur et la texture.
\item \textbf{Pharmacie et cosmétiques :} Les sirops, collyres et crèmes sont formulés à pH constant (tampons phosphate, citrate, lactate) pour garantir la stabilité chimique du principe actif, son absorption cutanée ou muqueuse, et l'absence d'irritation.
\item \textbf{Étalonnage pH-mètre :} Les solutions tampons étalons (pH 4,01 ; 7,00 ; 9,21 à 25 °C) sont des références primaires ou secondaires traçables au SI.
\end{itemize}
\end{savaistu}

\courssub{Conclusion}

Nous avons préparé une solution tampon éthanoïque en mélangeant des volumes égaux d'acide éthanoïque et d'éthanoate de sodium de même concentration : le pH mesuré ($\approx 4{,}8$) est bien égal au $\mathrm{p}K_a$ du couple, conformément à la relation de Henderson-Hasselbalch. Les ajouts modérés d'acide chlorhydrique ou de soude ne font varier le pH du tampon que d'environ $0{,}1$ unité, alors qu'ils provoquent des sauts de 4 unités dans l'eau distillée ; la dilution laisse également le pH pratiquement inchangé. Nous avons ainsi vérifié expérimentalement la définition d'une solution tampon : \emph{son pH varie peu par addition modérée d'acide ou de base, ou par dilution modérée}. C'est précisément ce mécanisme qui maintient le pH de ton sang autour de 7{,}4 grâce au couple \ce{H2CO3}/\ce{HCO3-}, et qui garantit la stabilité de nombreux médicaments et produits alimentaires fabriqués au Cameroun.

\newpage

\coursec{Méthodes et exercices résolus}

\courssub{Savoir-faire 1 : reconnaître et définir une solution tampon}

\begin{definition}[Solution tampon]
Une \cle{solution tampon} (ou solution tamponnée) est une solution dont le pH varie \cle{peu} lors de l'addition d'une quantité modérée d'acide fort, de base forte, ou lors d'une dilution modérée. Elle est constituée d'un \cle{couple acide faible / base conjuguée} \ce{AH}/\ce{A-} (ou \ce{BH+}/\ce{B}) dont les deux espèces sont présentes en \cle{quantités comparables} (ni l'une ni l'autre n'est négligeable).
\end{definition}

\begin{popfigure}
\begin{tikzpicture}[scale=0.9, every node/.style={font=\small}]
% Beaker left: Acid addition
\begin{scope}[xshift=-4cm]
\draw[thick, popblue] (0,0) -- (0,3.5) -- (0.2,3.7) -- (2.2,3.7) -- (2.4,3.5) -- (2.4,0) -- cycle;
\draw[thick, popblue] (0.2,3.7) -- (2.2,3.7);
\fill[popblue!20] (0.2,0.2) rectangle (2.2,3.0);
\node[popblue, align=center] at (1.2,3.3) {Avant ajout};
\node[align=center] at (1.2,2.5) {\ce{AH} (majoritaire)\\\ce{A-} (majoritaire)};
\draw[->, thick, poporange] (1.2,3.8) -- (1.2,4.5) node[above, align=center] {Ajout \ce{H3O+}};
\draw[thick, popblue] (0,5) -- (0,8.5) -- (0.2,8.7) -- (2.2,8.7) -- (2.4,8.5) -- (2.4,5) -- cycle;
\fill[popblue!20] (0.2,5.2) rectangle (2.2,8.0);
\node[align=center] at (1.2,7.5) {\ce{A- + H3O+ -> AH + H2O}};
\node[align=center] at (1.2,6.5) {\ce{AH} \textbf{augmente}\\\ce{A-} \textbf{diminue}};
\node[popdark, align=center] at (1.2,4.8) {pH \textbf{quasi stable}};
\end{scope}

% Beaker right: Base addition
\begin{scope}[xshift=4cm]
\draw[thick, popgreen] (0,0) -- (0,3.5) -- (0.2,3.7) -- (2.2,3.7) -- (2.4,3.5) -- (2.4,0) -- cycle;
\fill[popgreen!20] (0.2,0.2) rectangle (2.2,3.0);
\node[popgreen, align=center] at (1.2,3.3) {Avant ajout};
\node[align=center] at (1.2,2.5) {\ce{AH} (majoritaire)\\\ce{A-} (majoritaire)};
\draw[->, thick, poporange] (1.2,3.8) -- (1.2,4.5) node[above, align=center] {Ajout \ce{HO-}};
\draw[thick, popgreen] (0,5) -- (0,8.5) -- (0.2,8.7) -- (2.2,8.7) -- (2.4,8.5) -- (2.4,5) -- cycle;
\fill[popgreen!20] (0.2,5.2) rectangle (2.2,8.0);
\node[align=center] at (1.2,7.5) {\ce{AH + HO- -> A- + H2O}};
\node[align=center] at (1.2,6.5) {\ce{AH} \textbf{diminue}\\\ce{A-} \textbf{augmente}};
\node[popdark, align=center] at (1.2,4.8) {pH \textbf{quasi stable}};
\end{scope}

% Central equilibrium
\node[align=center, popdark] at (0,-0.5) {\textbf{Équilibre acido-basique :} \ce{AH + H2O <=> A- + H3O+}};
\node[align=center, poppurple] at (0,-1.2) {Le principe de Le Chatelier assure la stabilité du pH};
\end{tikzpicture}
\legende{Mécanisme d'action d'une solution tampon \ce{AH}/\ce{A-} : consommation des ions \ce{H3O+} ou \ce{HO-} ajoutés par réaction totale avec l'espèce conjuguée appropriée.}
\end{popfigure}

\begin{methode}[Identifier une solution tampon]
\begin{enumerate}
\item \textbf{Réflexe de composition} : une solution tampon contient un couple acide faible/base faible conjugués \ce{AH}/\ce{A-} en quantités \cle{comparables}. Exemples types : 
\chemfig{CH_3-C(=[1]O)-[7]OH}/\chemfig{CH_3-C(=[1]O)-[7]{O^{-}}} (\ce{CH3COOH}/\ce{CH3COO-}), 
\chemfig{H_3C-CH_2-C(=[1]O)-[7]OH}/\chemfig{H_3C-CH_2-C(=[1]O)-[7]{O^{-}}} (propanoïque), 
\ce{NH4+}/\chemfig{H-N(-[2]H)(-[6]H)} (\ce{NH4+}/\ce{NH3}), \ce{H2CO3}/\ce{HCO3-}.
\item \textbf{Vérifier l'absence de réaction totale} : l'acide et la base présents doivent être \emph{conjugués} (sinon ils réagissent entre eux jusqu'à disparition de l'un des deux).
\item \textbf{Énoncer les trois propriétés} : le pH d'une solution tampon varie \cle{peu} par (a) addition modérée d'acide fort, (b) addition modérée de base forte, (c) dilution modérée.
\item \textbf{Justifier} : l'acide ajouté est consommé par \ce{A-} (\ce{A- + H3O+ -> AH + H2O}) et la base ajoutée est consommée par \ce{AH} (\ce{AH + HO- -> A- + H2O}) : les ions \ce{H3O+} ou \ce{HO-} introduits ne s'accumulent pas.
\end{enumerate}
\end{methode}

\begin{exoresolu}[QCM de reconnaissance — type Bac]
On dispose des solutions suivantes, toutes à \SI{0,10}{mol.L^{-1}} : (a) \ce{HCl} ; (b) \ce{CH3COOH} seul ; (c) mélange équimolaire \ce{CH3COOH} + \ce{CH3COONa} ; (d) mélange \ce{HCl} + \ce{NaCl} ; (e) mélange équimolaire \ce{NH4Cl} + \ce{NH3}. Identifier la (ou les) solution(s) tampon(s) et justifier.
\tcblower
\textbf{Analyse de chaque solution :}
\begin{itemize}
\item (a) \ce{HCl} : acide fort seul, pas de base conjuguée en quantité comparable $\Rightarrow$ \textbf{pas tampon}.
\item (b) \ce{CH3COOH} seul : acide faible, mais sa base conjuguée \ce{CH3COO-} n'existe qu'en quantité infime (ionisation partielle, $\tau \ll 1$) $\Rightarrow$ \textbf{pas tampon}.
\item (c) \ce{CH3COOH} + \ce{CH3COONa} : l'acétate de sodium est totalement dissocié (\ce{CH3COONa -> Na+ + CH3COO-}). On a donc le couple \ce{CH3COOH}/\ce{CH3COO-} en quantités égales et comparables $\Rightarrow$ \textbf{solution tampon}, de $\text{pH} = \text{p}K_a = 4,8$.
\item (d) \ce{HCl} + \ce{NaCl} : \ce{Cl-} est une base inerte (spectatrice), sans propriété basique (conjuguée d'un acide fort) $\Rightarrow$ \textbf{pas tampon}.
\item (e) \ce{NH4Cl} + \ce{NH3} : couple \ce{NH4+}/\ce{NH3} en quantités équimolaires $\Rightarrow$ \textbf{solution tampon}, de $\text{pH} = \text{p}K_a = 9,2$.
\end{itemize}
\textbf{Conclusion :} les solutions (c) et (e) sont des solutions tampons. Le critère décisif est la présence simultanée, en quantités comparables, d'un acide faible et de sa base conjuguée.
\end{exoresolu}

\courssub{Savoir-faire 2 : calculer le pH d'un mélange acide faible / base conjuguée}

\begin{propriete}[Relation de Henderson-Hasselbalch]
Pour un mélange contenant un acide faible \ce{AH} et sa base conjuguée \ce{A-} en quantités non négligeables, le pH est donné par :
\[ \text{pH} = \text{p}K_a(\ce{AH}/\ce{A-}) + \log\frac{[\ce{A-}]}{[\ce{AH}]} = \text{p}K_a + \log\frac{n(\ce{A-})}{n(\ce{AH})} \]
Le volume total de la solution se simplifie dans le rapport des concentrations : on peut utiliser directement les quantités de matière $n$ (en mol ou \si{mmol}).
\end{propriete}

\begin{methode}[Appliquer la relation de Henderson-Hasselbalch]
\begin{enumerate}
\item \textbf{Identifier le couple} \ce{AH}/\ce{A-} et relever son $\text{p}K_a$ (donné ou déduit de $K_a$).
\item \textbf{Déterminer les quantités} $n(\ce{AH})$ et $n(\ce{A-})$ \emph{après mélange} (attention : si une réaction acido-basique totale a lieu lors du mélange — ex. ajout d'acide fort sur base faible —, faire d'abord un tableau d'avancement de cette réaction totale !).
\item \textbf{Appliquer la relation} : $\text{pH} = \text{p}K_a + \log\dfrac{n(\ce{A-})}{n(\ce{AH})}$.
\item \textbf{Cas particulier} : si $n(\ce{AH}) = n(\ce{A-})$, alors $\text{pH} = \text{p}K_a$ (mélange équimolaire).
\item \textbf{Vérifier la cohérence} : le pH doit rester proche du $\text{p}K_a$ (en pratique dans l'intervalle $\text{p}K_a \pm 1$).
\end{enumerate}
\end{methode}

\begin{exoresolu}[pH d'un mélange acide éthanoïque / éthanoate]
On mélange $V_a = \SI{50,0}{mL}$ d'acide éthanoïque \ce{CH3COOH} à $C_a = \SI{0,20}{mol.L^{-1}}$ avec $V_b = \SI{100,0}{mL}$ d'éthanoate de sodium \ce{CH3COONa} à $C_b = \SI{0,10}{mol.L^{-1}}$. Donnée : $\text{p}K_a(\ce{CH3COOH}/\ce{CH3COO-}) = 4,8$.
\begin{enumerate}
\item Calculer le pH du mélange.
\item On ajoute \SI{1,0}{mmol} d'ions \ce{H3O+} (sans variation de volume). Calculer la variation de pH et conclure.
\end{enumerate}
\tcblower
\textbf{1. pH du mélange.}

Quantités apportées (l'acétate de sodium est totalement dissocié ; \ce{Na+} est spectateur) :
\begin{align*}
n(\ce{CH3COOH}) &= C_a V_a = 0,20 \times 50,0\times 10^{-3} = \SI{1,0e-2}{mol} = \SI{10,0}{mmol}\\
n(\ce{CH3COO-}) &= C_b V_b = 0,10 \times 100,0\times 10^{-3} = \SI{1,0e-2}{mol} = \SI{10,0}{mmol}
\end{align*}
Les deux espèces conjuguées coexistent sans réaction totale entre elles. La relation de Henderson-Hasselbalch s'applique :
\[ \text{pH} = \text{p}K_a + \log\frac{n(\ce{CH3COO-})}{n(\ce{CH3COOH})} = 4,8 + \log\frac{10,0}{10,0} = 4,8 + \log 1 \]
\[ \boxed{\text{pH} = 4,8} \]
Le mélange est équimolaire : c'est une solution tampon de $\text{pH} = \text{p}K_a$.

\textbf{2. Ajout de \SI{1,0}{mmol} de \ce{H3O+}.}

Les ions \ce{H3O+} réagissent de façon totale sur la base du couple :
\[ \ce{CH3COO- + H3O+ -> CH3COOH + H2O} \]
Tableau d'avancement (quantités en \si{mmol}) :
\begin{center}
\begin{tabularx}{\linewidth}{|l|X|X|X|}
\hline
\rowcolor{popblueL} & \ce{CH3COO-} & \ce{H3O+} & \ce{CH3COOH} \\
\hline
État initial & $10,0$ & $1,0$ & $10,0$ \\
\hline
État final & $9,0$ & $0$ & $11,0$ \\
\hline
\end{tabularx}
\end{center}
Nouveau pH :
\[ \text{pH}' = 4,8 + \log\frac{9,0}{11,0} = 4,8 + \log(0,818) = 4,8 - 0,087 = 4,71 \]
Variation : $\Delta\text{pH} = 4,71 - 4,80 = -0,09$ unité de pH.

\textbf{Conclusion :} l'ajout de \SI{1,0}{mmol} d'acide fort ne fait varier le pH que de $0,09$ unité. À titre de comparaison, la même quantité d'acide ajoutée à \SI{150}{mL} d'eau pure ferait chuter le pH de $7,0$ à environ $2,2$ ($\Delta\text{pH} \approx -4,8$). Le mélange joue bien son rôle de \cle{tampon}.
\end{exoresolu}

\courssub{Savoir-faire 3 : préparer une solution tampon de pH = pKa}

\begin{methode}[Trois stratégies de préparation pour $\text{pH} = \text{p}K_a$]
\begin{enumerate}
\item \textbf{Mélange direct équimolaire} : mélanger des volumes d'acide faible et de sel de la base conjuguée tels que $C_a V_a = C_b V_b$. Alors $\text{pH} = \text{p}K_a$.
\item \textbf{Demi-neutralisation d'un acide faible par une base forte} : ajouter à $n_0$ mol d'acide faible exactement $n_0/2$ mol de soude (ou potasse). La réaction totale \ce{AH + HO- -> A- + H2O} transforme la moitié de l'acide en base conjuguée : $n(\ce{AH}) = n(\ce{A-}) = n_0/2$, donc $\text{pH} = \text{p}K_a$. C'est le point de \cle{demi-équivalence}.
\item \textbf{Demi-neutralisation d'une base faible par un acide fort} : ajouter à $n_0$ mol de base faible exactement $n_0/2$ mol d'acide fort : $n(\ce{B}) = n(\ce{BH+}) = n_0/2$, donc $\text{pH} = \text{p}K_a$.
\item \textbf{Pour un pH différent du pKa} : imposer le rapport $\dfrac{n(\ce{A-})}{n(\ce{AH})} = 10^{\text{pH}-\text{p}K_a}$ et en déduire les volumes.
\end{enumerate}
\end{methode}

\begin{exoresolu}[Préparation d'un tampon ammoniacal — type Bac]
Un technicien du laboratoire d'une brasserie doit préparer \SI{200,0}{mL} de solution tampon de $\text{pH} = 9,2$. Il dispose d'une solution d'ammoniac \ce{NH3} à $C_b = \SI{0,40}{mol.L^{-1}}$ et d'une solution d'acide chlorhydrique à $C_a = \SI{0,20}{mol.L^{-1}}$. Donnée : $\text{p}K_a(\ce{NH4+}/\ce{NH3}) = 9,2$.
\begin{enumerate}
\item Expliquer la stratégie de préparation par demi-neutralisation.
\item Calculer les volumes $V_b$ d'ammoniac et $V_a$ d'acide chlorhydrique à mélanger pour obtenir \SI{200,0}{mL} de tampon.
\end{enumerate}
\tcblower
\textbf{1. Stratégie.}

On veut $\text{pH} = \text{p}K_a = 9,2$, donc $n(\ce{NH4+}) = n(\ce{NH3})$ dans le mélange final. En versant de l'acide chlorhydrique (acide fort) sur l'ammoniac (base faible), la réaction totale suivante a lieu :
\[ \ce{NH3 + H3O+ -> NH4+ + H2O} \]
Si l'on neutralise exactement la \textbf{moitié} de l'ammoniac initial, il reste $n_0/2$ de \ce{NH3} et il s'est formé $n_0/2$ de \ce{NH4+} : le mélange est équimolaire, c'est un tampon de $\text{pH} = \text{p}K_a = 9,2$.

\textbf{2. Calcul des volumes.}

Soit $n_0 = C_b V_b$ la quantité initiale d'ammoniac. La demi-neutralisation exige :
\[ n(\ce{H3O+}) = \frac{n_0}{2} \quad\Longrightarrow\quad C_a V_a = \frac{C_b V_b}{2} \]
Contrainte de volume final : $V_a + V_b = \SI{200,0}{mL}$.

De la première relation : $0,20\,V_a = \dfrac{0,40\,V_b}{2} = 0,20\,V_b$, donc $V_a = V_b$.

Avec $V_a + V_b = \SI{200,0}{mL}$ : $V_a = V_b = \SI{100,0}{mL}$.

\textbf{Vérification :} $n_0 = 0,40 \times 0,1000 = \SI{4,0e-2}{mol}$ de \ce{NH3} ; acide versé : $0,20 \times 0,1000 = \SI{2,0e-2}{mol} = n_0/2$. Il reste \SI{2,0e-2}{mol} de \ce{NH3} et il s'est formé \SI{2,0e-2}{mol} de \ce{NH4+} : rapport égal à 1, $\text{pH} = 9,2$. \checkmark

\textbf{Protocole :} prélever \SI{100,0}{mL} de solution d'ammoniac à la fiole jaugée ou à l'éprouvette graduée, ajouter lentement \SI{100,0}{mL} d'acide chlorhydrique en agitant, puis vérifier le pH au pH-mètre étalonné.

\textbf{Conclusion :} il faut mélanger \SI{100,0}{mL} d'ammoniac à \SI{0,40}{mol.L^{-1}} et \SI{100,0}{mL} d'acide chlorhydrique à \SI{0,20}{mol.L^{-1}}.
\end{exoresolu}

\begin{exoresolu}[Préparer un tampon de pH imposé différent du pKa]
On veut préparer \SI{250,0}{mL} d'une solution tampon de $\text{pH} = 5,0$ à partir d'acide éthanoïque ($\text{p}K_a = 4,8$) à $C_a = \SI{0,10}{mol.L^{-1}}$ et d'éthanoate de sodium à $C_b = \SI{0,10}{mol.L^{-1}}$, de mêmes concentrations. Calculer les volumes $V_a$ et $V_b$ à mélanger.
\tcblower
\textbf{Condition sur le rapport :}
\[ \text{pH} = \text{p}K_a + \log\frac{n(\ce{CH3COO-})}{n(\ce{CH3COOH})} \Longrightarrow \log\frac{n_b}{n_a} = 5,0 - 4,8 = 0,2 \]
\[ \frac{n_b}{n_a} = 10^{0,2} \approx 1,585 \]
Les concentrations étant égales ($C_a = C_b$), le rapport des quantités est égal au rapport des volumes :
\[ \frac{V_b}{V_a} = 1,585 \quad\text{et}\quad V_a + V_b = \SI{250,0}{mL} \]
D'où $V_b = 1,585\,V_a$, puis $V_a(1 + 1,585) = 250,0$, soit :
\[ V_a = \frac{250,0}{2,585} \approx \SI{96,7}{mL} \qquad V_b = 250,0 - 96,7 \approx \SI{153,3}{mL} \]
\textbf{Vérification :} $\text{pH} = 4,8 + \log(153,3/96,7) = 4,8 + \log(1,585) = 4,8 + 0,20 = 5,0$. \checkmark

\textbf{Conclusion :} mélanger environ \SI{96,7}{mL} d'acide éthanoïque et \SI{153,3}{mL} d'éthanoate de sodium. Remarque : pour obtenir un pH \emph{supérieur} au $\text{p}K_a$, il faut davantage de base conjuguée que d'acide ($V_b > V_a$).
\end{exoresolu}

\courssub{Savoir-faire 4 : interpréter la zone tampon d'une courbe de dosage}

\begin{popfigure}
\begin{tikzpicture}
\begin{axis}[
    width=\linewidth, height=8cm,
    xlabel={Volume de base forte versé $V_b$ (\si{mL})},
    ylabel={pH},
    xmin=0, xmax=25,
    ymin=2, ymax=12,
    grid=major, grid style={popline},
    axis lines=middle,
    enlargelimits={abs=0.5},
    tick label style={font=\small},
    label style={font=\small},
    ]
% Sigmoid curve approximation for weak acid titration
\addplot[popcurve, domain=0:24, samples=300, smooth] 
    ({x}, {4.8 + 2.5*atan(3*(x-16))/90 + 1.5*atan(10*(x-8))/90}); % rough shape
% Better analytical approximation for weak acid titration curve
% Using Henderson-Hasselbalch before equivalence, then excess base
% Let's draw a schematic curve using coordinates for clarity
\addplot[popcurve, thick, smooth] coordinates {
    (0, 2.9) (2, 3.8) (4, 4.3) (6, 4.6) (8, 4.8) (10, 5.0) (12, 5.3) (14, 6.0) (15.5, 8.5) (16, 8.8) (16.5, 10.0) (18, 10.5) (20, 10.8) (24, 11.0)
};
% Equivalence point
\draw[poporange, dashed, thick] (16,2) -- (16,11) node[above, pos=0.9, poporange, font=\small] {$V_E = \SI{16,0}{mL}$};
% Half-equivalence
\draw[popgreen, dashed, thick] (8,2) -- (8,11) node[above, pos=0.5, popgreen, font=\small] {$V_E/2 = \SI{8,0}{mL}$};
% Buffer zone highlight
\fill[popgreen!15] (4,2) rectangle (12,11);
\draw[popgreen, thick, <->] (4, 10.5) -- (12, 10.5) node[midway, above, font=\small] {Zone tampon};
% pKa label
\node[poppurple, font=\small, fill=popcream, rounded corners, inner sep=2pt] at (8, 5.5) {$\text{pH} = \text{p}K_a = 4,8$};
\end{axis}
\end{tikzpicture}
\legende{Courbe de dosage d'un acide faible par une base forte : zone tampon (verte) centrée sur la demi-équivalence ($V_E/2$), où $\text{pH} = \text{p}K_a$. Le saut de pH (orange) marque l'équivalence.}
\end{popfigure}

\begin{methode}[Exploiter la zone tampon sur pH = f(V)]
\begin{enumerate}
\item \textbf{Repérer la zone tampon} : région de la courbe où le pH varie \cle{lentement} (pente faible, courbe presque horizontale), située \emph{avant} le saut de pH, autour de la demi-équivalence $V_{E}/2$.
\item \textbf{Lire le pKa} : à la demi-équivalence, $n(\ce{AH}) = n(\ce{A-})$, donc $\text{pH} = \text{p}K_a$. Méthode : déterminer $V_E$ (milieu du saut de pH), puis lire le pH à $V_E/2$.
\item \textbf{Délimiter le domaine d'efficacité} : le tampon est efficace tant que $0,1 < \dfrac{[\ce{A-}]}{[\ce{AH}]} < 10$, c'est-à-dire pour $\text{pH} \in [\text{p}K_a - 1\,;\,\text{p}K_a + 1]$. Sur la courbe, cela correspond à un intervalle de volumes centré sur $V_E/2$.
\item \textbf{Qualifier le pouvoir tampon} : il est \cle{maximal} à la demi-équivalence ($\text{pH} = \text{p}K_a$) et d'autant plus grand que la solution est \cle{concentrée} (grande réserve d'acide et de base).
\end{enumerate}
\end{methode}

\begin{exoresolu}[Lecture d'une courbe de dosage — type Bac]
On dose $V_a = \SI{20,0}{mL}$ d'une solution d'acide éthanoïque de concentration $C_a$ inconnue par de la soude à $C_b = \SI{0,10}{mol.L^{-1}}$. La courbe $\text{pH} = f(V_b)$ présente un saut de pH pour $V_E = \SI{16,0}{mL}$, et le pH à la demi-équivalence vaut $4,8$.
\begin{enumerate}
\item Déterminer le $\text{p}K_a$ du couple et la concentration $C_a$.
\item Délimiter la zone tampon sur l'axe des volumes et expliquer pourquoi le pH y varie peu.
\item Où le pouvoir tampon est-il maximal ?
\end{enumerate}
\tcblower
\textbf{1. pKa et concentration.}

À la demi-équivalence, $V_b = V_E/2 = \SI{8,0}{mL}$, la moitié de l'acide a été transformée en ion éthanoate : $n(\ce{CH3COOH}) = n(\ce{CH3COO-})$, donc :
\[ \boxed{\text{p}K_a = \text{pH}_{V_E/2} = 4,8} \]
À l'équivalence : $C_a V_a = C_b V_E$, d'où :
\[ C_a = \frac{C_b V_E}{V_a} = \frac{0,10 \times 16,0}{20,0} = \SI{8,0e-2}{mol.L^{-1}} \]

\textbf{2. Zone tampon.}

Le tampon est efficace pour $\text{pH} \in [3,8\,;\,5,8]$, ce qui correspond sur la courbe à la région autour de $V_b = \SI{8,0}{mL}$ (en pratique de quelques mL après le début du dosage jusqu'aux abords du saut, soit environ $V_b \in [4\,;\,12]\,\si{mL}$).

Dans cette zone, la solution contient simultanément \ce{CH3COOH} et \ce{CH3COO-} en quantités comparables : les ions \ce{HO-} versés sont consommés par la réaction totale \ce{CH3COOH + HO- -> CH3COO- + H2O} et ne s'accumulent pas. Le rapport $\dfrac{[\ce{CH3COO-}]}{[\ce{CH3COOH}]}$ varie lentement, donc le pH aussi (faible pente de la courbe).

\textbf{3. Pouvoir tampon maximal.}

Le pouvoir tampon est maximal à la \textbf{demi-équivalence} ($V_b = \SI{8,0}{mL}$, $\text{pH} = \text{p}K_a = 4,8$), car c'est là que les deux espèces du couple sont en quantités égales : la solution peut absorber aussi bien un ajout d'acide qu'un ajout de base avec la même efficacité.

\textbf{Conclusion :} $\text{p}K_a = 4,8$ ; $C_a = \SI{8,0e-2}{mol.L^{-1}}$ ; zone tampon centrée sur $V_b = \SI{8,0}{mL}$, où le pouvoir tampon est maximal.
\end{exoresolu}

\courssub{Savoir-faire 5 : comparer l'effet tampon et choisir un tampon adapté}

\begin{methode}[Choisir et évaluer un tampon]
\begin{enumerate}
\item \textbf{Règle de choix} : pour tamponner à un pH donné, choisir un couple dont le $\text{p}K_a$ est le \cle{plus proche} du pH visé (idéalement $|\text{pH} - \text{p}K_a| \leq 1$).
\item \textbf{Comparer deux solutions} : ajouter la même petite quantité d'acide (ou de base) à chacune et calculer $\Delta\text{pH}$ ; le tampon est celui dont la variation est la plus faible.
\item \textbf{Effet de la dilution} : la dilution modifie peu le pH (le rapport $[\ce{A-}]/[\ce{AH}]$ est inchangé) mais \cle{diminue le pouvoir tampon} (moins de « réserve » d'acide et de base pour contrer les ajouts).
\end{enumerate}
\end{methode}

\begin{exoresolu}[Quel tampon pour une fermentation contrôlée ?]
Une unité de production de yaourt à Yaoundé doit maintenir le milieu à $\text{pH} \approx 4,6$. On hésite entre deux couples : \ce{CH3COOH}/\ce{CH3COO-} ($\text{p}K_a = 4,8$) et \ce{NH4+}/\ce{NH3} ($\text{p}K_a = 9,2$).
\begin{enumerate}
\item Quel couple choisir ? Justifier.
\item On ajoute \SI{2,0}{mmol} de soude à \SI{100}{mL} d'un tampon équimolaire \ce{CH3COOH}/\ce{CH3COO-} contenant \SI{10}{mmol} de chaque espèce. Calculer le pH final.
\end{enumerate}
\tcblower
\textbf{1. Choix du couple.}

Un tampon est efficace dans la zone $\text{p}K_a \pm 1$ :
\begin{itemize}
\item \ce{CH3COOH}/\ce{CH3COO-} : zone $[3,8\,;\,5,8]$ $\rightarrow$ contient $\text{pH} = 4,6$. \checkmark
\item \ce{NH4+}/\ce{NH3} : zone $[8,2\,;\,10,2]$ $\rightarrow$ ne convient pas.
\end{itemize}
On choisit le couple \textbf{acide éthanoïque / ion éthanoate}, dont le $\text{p}K_a = 4,8$ est proche de $4,6$.

\textbf{2. pH après ajout de soude.}

La base forte réagit totalement sur l'acide du couple :
\[ \ce{CH3COOH + HO- -> CH3COO- + H2O} \]
Tableau d'avancement (en \si{mmol}) :
\begin{center}
\begin{tabularx}{\linewidth}{|l|X|X|X|}
\hline
\rowcolor{popblueL} & \ce{CH3COOH} & \ce{HO-} & \ce{CH3COO-} \\
\hline
État initial & $10,0$ & $2,0$ & $10,0$ \\
\hline
État final & $8,0$ & $0$ & $12,0$ \\
\hline
\end{tabularx}
\end{center}
\[ \text{pH} = 4,8 + \log\frac{12,0}{8,0} = 4,8 + \log(1,5) = 4,8 + 0,176 = 4,98 \]
\textbf{Conclusion :} le pH passe de $4,80$ à $4,98$, soit $\Delta\text{pH} = +0,18$ seulement : le tampon a bien amorti l'ajout de base. Le milieu reste dans la zone utile pour la fermentation.
\end{exoresolu}

\begin{experience}[TP : Préparation et test d'une solution tampon acide éthanoïque / éthanoate]
\textbf{Objectif :} Préparer \SI{100,0}{mL} d'une solution tampon de $\text{pH} = 4,8$ et mettre en évidence son pouvoir tampon par addition d'acide et de base forts.
\textbf{Matériel :} Acide éthanoïque \SI{0,20}{mol.L^{-1}}, éthanoate de sodium \SI{0,20}{mol.L^{-1}}, acide chlorhydrique \SI{0,10}{mol.L^{-1}}, soude \SI{0,10}{mol.L^{-1}}, pH-mètre étalonné, fioles jaugées \SI{100}{mL}, béchers, pipettes, agitateur magnétique.
\textbf{Protocole :}
\begin{enumerate}
\item \textbf{Préparation :} Mélanger \SI{50,0}{mL} d'acide éthanoïque \SI{0,20}{mol.L^{-1}} et \SI{50,0}{mL} d'éthanoate de sodium \SI{0,20}{mol.L^{-1}} dans une fiole jaugée \SI{100}{mL}. Compléter à l'eau distillée. Mesurer le pH initial (théorique 4,8).
\item \textbf{Test acide :} Prélever \SI{25,0}{mL} de tampon dans un bécher. Ajouter \SI{1,0}{mL} de \ce{HCl} \SI{0,10}{mol.L^{-1}} (\SI{0,10}{mmol}). Mesurer le pH final. Calculer $\Delta\text{pH}$.
\item \textbf{Test base :} Prélever \SI{25,0}{mL} de tampon. Ajouter \SI{1,0}{mL} de \ce{NaOH} \SI{0,10}{mol.L^{-1}} (\SI{0,10}{mmol}). Mesurer le pH final. Calculer $\Delta\text{pH}$.
\item \textbf{Témoin (eau) :} Faire les mêmes ajouts sur \SI{25,0}{mL} d'eau distillée (pH initial $\approx 7$) et mesurer les variations.
\end{enumerate}
\textbf{Interprétation attendue :} $\Delta\text{pH}_{\text{tampon}} \ll \Delta\text{pH}_{\text{eau}}$. Le tampon résiste ; l'eau non.
\end{experience}

\begin{savaistu}[Le sang : un tampon vital au quotidien]
Le pH du sang humain est maintenu à \SI{7,40}{\pm 0,05} grâce à plusieurs systèmes tampons couplés :
\begin{itemize}
\item \textbf{Le couple carbonique/bicarbonate} \ce{H2CO3}/\ce{HCO3-} ($\text{p}K_a' \approx 6,1$ à \SI{37}{\celsius}, concentration \SI{24}{mmol.L^{-1}} pour \ce{HCO3-}, \SI{1,2}{mmol.L^{-1}} pour \ce{CO2} dissous). C'est le tampon principal du plasma. L'élimination du \ce{CO2} par les poumons (respiration) et la réabsorption rénale de \ce{HCO3-} en font un tampon \emph{ouvert} (non isolé), d'une efficacité remarquable.
\item \textbf{L'hémoglobine} (Hb) dans les globules rouges : ses groupes histidine agissent comme couples \ce{HbH+}/\ce{Hb} ($\text{p}K_a \approx 6,8-7,8$).
\item \textbf{Les phosphates} \ce{H2PO4-}/\ce{HPO4^2-} ($\text{p}K_a = 7,2$) : tampon intracellulaire et urinaire.
\end{itemize}
\textbf{Application camerounaise :} Dans les laboratoires d'analyses médicales (Centre Pasteur du Cameroun, hôpitaux de référence), le dosage des gaz du sang (pH, $p\ce{CO2}$, $[\ce{HCO3-}]$) permet de diagnostiquer les déséquilibres acido-basiques (acidose/cétoacidose diabétique, alcalose respiratoire) et d'ajuster les perfusions (bicarbonate de sodium, lactate de Ringer).
\end{savaistu}

\begin{attention}[Les 5 erreurs qui coûtent des points au Bac]
\begin{enumerate}
\item \textbf{Croire qu'un acide faible seul est un tampon.} Sans sa base conjuguée en quantité comparable, il n'y a pas d'effet tampon. Il faut \emph{les deux} espèces du couple.
\item \textbf{Oublier la réaction préalable lors du mélange.} Si l'on mélange \ce{AH} avec une base forte (ou \ce{A-} avec un acide fort), il faut \emph{d'abord} faire un tableau d'avancement de la réaction totale, \emph{puis} appliquer Henderson-Hasselbalch sur les quantités restantes. Appliquer la formule sur les quantités initiales est faux.
\item \textbf{Confondre équivalence et demi-équivalence.} C'est à la \textbf{demi}-équivalence ($V_E/2$) que $\text{pH} = \text{p}K_a$ ; à l'équivalence d'un dosage d'acide faible, le pH est basique ($> 7$) et n'est \emph{pas} égal au $\text{p}K_a$.
\item \textbf{Écrire le rapport à l'envers.} $\text{pH} = \text{p}K_a + \log\dfrac{[\text{base}]}{[\text{acide}]}$ : la base conjuguée est au \emph{numérateur}. Un pH supérieur au $\text{p}K_a$ signifie plus de base que d'acide — vérifie toujours la cohérence.
\item \textbf{Négliger les unités et les chiffres significatifs.} Convertir les volumes en litres dans $n = CV$ (ou rester cohérent en \si{mmol} avec $V$ en \si{mL}), et donner le pH avec deux décimales. Une variation de pH se note $\Delta\text{pH}$ avec son signe.
\end{enumerate}
\end{attention}

\begin{exercice}[Entraînement Bac — Solutions tampons]
\begin{enumerate}
\item On mélange \SI{30,0}{mL} d'une solution d'acide éthanoïque \SI{0,10}{mol.L^{-1}} avec \SI{20,0}{mL} d'une solution de soude \SI{0,10}{mol.L^{-1}}. $\text{p}K_a = 4,8$. Calculer le pH du mélange obtenu. Ce mélange est-il une solution tampon ? Justifier.
\item On souhaite préparer \SI{500,0}{mL} d'une solution tampon \ce{NH4+}/\ce{NH3} de $\text{pH} = 9,0$ ($\text{p}K_a = 9,2$). On dispose de \ce{NH3} \SI{0,20}{mol.L^{-1}} et de \ce{NH4Cl} solide ($M = \SI{53,5}{g.mol^{-1}}$). Calculer le volume de solution d'ammoniac et la masse de chlorure d'ammonium à dissoudre.
\item \textbf{Épreuve pratique :} On dose \SI{20,0}{mL} d'une solution d'acide faible AH (\SI{0,10}{mol.L^{-1}}) par une solution de soude \SI{0,10}{mol.L^{-1}}. La courbe $\text{pH} = f(V_b)$ montre un saut de pH pour $V_E = \SI{20,0}{mL}$. Le pH à $V_b = \SI{10,0}{mL}$ est $4,5$.
\begin{enumerate}
\item Déterminer $\text{p}K_a$ et la nature possible de l'acide AH (données : acide méthanoïque 3,7 ; acide éthanoïque 4,8 ; acide benzoïque 4,2 ; acide hypochloreux 7,5).
\item Tracer schématiquement la courbe et hachurer la zone tampon. Préciser les bornes de pH de cette zone.
\end{enumerate}
\end{enumerate}
\end{exercice}

\begin{corrige}[Exercice 1]
\textbf{1.} Réaction totale : \ce{CH3COOH + HO- -> CH3COO- + H2O}.
$n_0(\ce{CH3COOH}) = 3,0\,\si{mmol}$ ; $n_0(\ce{HO-}) = 2,0\,\si{mmol}$.
Avancement max $x_{\max} = 2,0\,\si{mmol}$ (réactif limitant : \ce{HO-}).
État final : $n(\ce{CH3COOH}) = 1,0\,\si{mmol}$ ; $n(\ce{CH3COO-}) = 2,0\,\si{mmol}$.
Présence des deux espèces du couple en quantités comparables $\Rightarrow$ \textbf{solution tampon}.
$\text{pH} = 4,8 + \log(2,0/1,0) = 4,8 + 0,30 = \boxed{5,10}$.
\end{corrige}

\begin{corrige}[Exercice 2]
$\text{pH} = \text{p}K_a + \log\frac{n(\ce{NH3})}{n(\ce{NH4+})} \Rightarrow 9,0 = 9,2 + \log\frac{n_b}{n_a} \Rightarrow \log\frac{n_b}{n_a} = -0,2 \Rightarrow \frac{n_b}{n_a} = 10^{-0,2} \approx 0,631$.
Soit $n_a = n(\ce{NH4+})$, $n_b = n(\ce{NH3}) = 0,631\,n_a$.
Volume final $V = \SI{0,500}{L}$. Concentration totale $C_{\text{tot}} = \frac{n_a + n_b}{0,500}$.
Choisissons une concentration totale raisonnable, ex. $C_{\text{tot}} = \SI{0,10}{mol.L^{-1}}$ (pouvoir tampon correct).
$n_a + n_b = 0,050\,\si{mol} \Rightarrow n_a(1+0,631) = 0,050 \Rightarrow n_a \approx \SI{3,07e-2}{mol}$, $n_b \approx \SI{1,93e-2}{mol}$.
Volume \ce{NH3} : $V_b = n_b / 0,20 = \SI{96,5}{mL}$.
Masse \ce{NH4Cl} : $m = n_a \times 53,5 \approx \SI{1,64}{g}$.
\textit{Note : d'autres couples $(V_b, m)$ sont possibles si on change $C_{\text{tot}}$ ; le rapport $n_b/n_a$ est imposé.}
\end{corrige}

\begin{corrige}[Exercice 3]
\textbf{a.} $V_E = \SI{20,0}{mL} \Rightarrow V_E/2 = \SI{10,0}{mL}$. À ce volume, $\text{pH} = 4,5 = \text{p}K_a$.
L'acide est l'\textbf{acide éthanoïque} ($\text{p}K_a = 4,8 \approx 4,5$ ; l'acide benzoïque a $\text{p}K_a = 4,2$, l'acide méthanoïque 3,7, l'hypochloreux 7,5). La valeur expérimentale 4,5 est compatible avec l'acide éthanoïque (erreur de mesure ou force ionique).
\textbf{b.} Zone tampon : $\text{pH} \in [\text{p}K_a - 1\,;\, \text{p}K_a + 1] = [3,5\,;\,5,5]$. Centrée sur $V_b = \SI{10,0}{mL}$. Courbe sigmoïde avec plateau autour de 4,5 entre $\approx \SI{4}{mL}$ et $\approx \SI{16}{mL}$.
\end{corrige}

\newpage

\coursec{Exercices}

\coursec{Les solutions tampons}

\courssub{1. Définition et propriétés d'une solution tampon}

\begin{definition}[Solution tampon]
Une \textbf{solution tampon} (ou solution tamponnée) est une solution dont le pH varie très peu :
\begin{itemize}
  \item lors de l'addition d'une \textbf{quantité modérée} d'acide (ions \ce{H3O+}) ou de base (ions \ce{HO-}) ;
  \item lors d'une \textbf{dilution modérée} (ajout d'eau).
\end{itemize}
Elle est constituée d'un \textbf{couple acide/base conjugué} \ce{AH}/\ce{A-} (acide faible / base conjuguée) ou \ce{BH+}/\ce{B} (acide conjugué / base faible), les deux espèces étant présentes en \textbf{concentrations voisines} (généralement rapport entre 0,1 et 10).
\end{definition}

\begin{propriete}[Propriétés fondamentales]
Une solution tampon possède trois propriétés caractéristiques :
\begin{enumerate}
  \item \textbf{Stabilité du pH par addition d'acide :} les ions \ce{H3O+} ajoutés sont consommés par la base \ce{A-} (ou \ce{B}) : \ce{A- + H3O+ -> AH + H2O}.
  \item \textbf{Stabilité du pH par addition de base :} les ions \ce{HO-} ajoutés sont consommés par l'acide \ce{AH} (ou \ce{BH+}) : \ce{AH + HO- -> A- + H2O}.
  \item \textbf{Stabilité du pH par dilution :} le pH dépend du \textbf{rapport} des concentrations (ou quantités de matière) $[\ce{A-}]/[\ce{AH}]$, qui ne change pas lors d'une dilution (les deux concentrations sont divisées par le même facteur).
\end{enumerate}
\end{propriete}

\begin{aretenir}[Condition d'existence]
Pour qu'un mélange soit tampon, il faut :
\begin{itemize}
  \item un couple acide/base \textbf{faible} ($K_a$ entre $10^{-4}$ et $10^{-10}$ environ) ;
  \item les deux espèces du couple présentes en \textbf{quantités comparables} (ni l'une ni l'autre en excès majeur) ;
  \item des concentrations \textbf{suffisamment élevées} (généralement $> \SI{10^{-2}}{mol.L^{-1}}$) pour avoir un pouvoir tampon significatif.
\end{itemize}
\end{aretenir}

\begin{savaistu}[Le tampon sanguin : une question de vie ou de mort]
Le sang humain est l'exemple type de solution tampon. Son pH doit rester strictement entre \textbf{7,35 et 7,45}. Le couple principal est \ce{H2CO3}/\ce{HCO3-} (acide carbonique / hydrogénocarbonate), avec un $\mathrm{p}K_a \approx 6,4$ (valeur apparente à \SI{37}{\celsius}). Le rein et les poumons régulent les concentrations de \ce{HCO3-} et \ce{CO2} dissous pour maintenir ce pH. Une déviation (acidose ou alcalose) met en péril les fonctions enzymatiques vitales.
\end{savaistu}

\courssub{2. pH d'un mélange acide faible / base conjuguée}

\begin{propriete}[Équation de Henderson-Hasselbalch]
Pour un mélange contenant un acide faible \ce{AH} et sa base conjuguée \ce{A-} (issus par exemple d'un sel \ce{Na+A-}), le pH est donné par :
\[
\mathrm{pH} = \mathrm{p}K_a(\ce{AH}/\ce{A-}) + \log\frac{[\ce{A-}]}{[\ce{AH}]}
\]
Comme les deux espèces sont dans le même volume de solution, on peut utiliser les quantités de matière :
\[
\mathrm{pH} = \mathrm{p}K_a + \log\frac{n_{\ce{A-}}}{n_{\ce{AH}}}
\]
Cette relation est valable si les approximations usuelles tiennent (acide faible, concentrations non nulles, force ionique modérée).
\end{propriete}

\begin{methode}[Calculer le pH d'un mélange tampon]
\begin{enumerate}
  \item Identifier le couple acide/base \ce{AH}/\ce{A-} et son $\mathrm{p}K_a$.
  \item Calculer les quantités de matière $n_{\ce{AH}}$ et $n_{\ce{A-}}$ introduites (ou restantes après réaction).
  \item Appliquer la formule $\mathrm{pH} = \mathrm{p}K_a + \log\dfrac{n_{\ce{A-}}}{n_{\ce{AH}}}$.
  \item Vérifier que le rapport $n_{\ce{A-}}/n_{\ce{AH}}$ est compris entre 0,1 et 10 (zone tampon efficace).
\end{enumerate}
\end{methode}

\begin{exemplebox}[Calcul de pH : mélange acétique]
On mélange \SI{50,0}{mL} d'acide éthanoïque \ce{CH3COOH} \SI{0,20}{mol.L^{-1}} ($n_{\ce{AH}} = \SI{0,010}{mol}$) et \SI{50,0}{mL} d'éthanoate de sodium \ce{CH3COONa} \SI{0,10}{mol.L^{-1}} ($n_{\ce{A-}} = \SI{0,0050}{mol}$). $\mathrm{p}K_a = 4,8$.
\[
\mathrm{pH} = 4,8 + \log\frac{0,0050}{0,010} = 4,8 + \log(0,5) = 4,8 - 0,30 = 4,5
\]
Le rapport est 0,5 (entre 0,1 et 10) : c'est une solution tampon de pH 4,5.
\end{exemplebox}

\courssub{3. Préparation d'une solution tampon de pH = pKa}

\begin{propriete}[Condition pH = pKa]
D'après Henderson-Hasselbalch, $\mathrm{pH} = \mathrm{p}K_a \iff \log\frac{[\ce{A-}]}{[\ce{AH}]} = 0 \iff [\ce{A-}] = [\ce{AH}]$ (mélange \textbf{équimolaire}).
\end{propriete}

\begin{methode}[Trois méthodes classiques pour préparer un tampon de pH = pKa]
\begin{enumerate}
  \item \textbf{Mélange direct équimolaire :} mélanger un volume $V_a$ d'une solution d'acide faible \ce{AH} (concentration $C$) avec un volume $V_b = V_a$ d'une solution de sa base conjuguée \ce{A-} (même concentration $C$).
  \item \textbf{Demi-neutralisation d'un acide faible par une base forte :} ajouter à un volume $V_a$ d'acide \ce{AH} (concentration $C_a$) un volume $V_b$ de base forte \ce{HO-} (concentration $C_b$) tel que $n_{\ce{HO-}} = \frac{1}{2}n_{\ce{AH},\text{initial}}$. La réaction \ce{AH + HO- -> A- + H2O} laisse $n_{\ce{AH}} = n_{\ce{A-}}$.
  \item \textbf{Demi-neutralisation d'une base faible par un acide fort :} ajouter à une base \ce{B} un acide fort \ce{H3O+} en quantité égale à la moitié de la quantité initiale de \ce{B}. La réaction \ce{B + H3O+ -> BH+} donne $n_{\ce{B}} = n_{\ce{BH+}}$.
\end{enumerate}
\end{methode}

\begin{experience}[TP : Préparation d'un tampon acétique pH = 4,8 (Méthode 2)]
\textbf{Matériel :} Acide éthanoïque \SI{0,20}{mol.L^{-1}}, Soude \SI{0,20}{mol.L^{-1}}, Bécher \SI{250}{mL}, pH-mètre étalonné, Agitateur magnétique.
\textbf{Protocole :}
\begin{enumerate}
  \item Verser \SI{100,0}{mL} d'acide éthanoïque \SI{0,20}{mol.L^{-1}} dans le bécher.
  \item Plonger l'électrode, agiter, noter le pH initial (acide, $\approx$ 2,7).
  \item Ajouter \textbf{lentement} \SI{50,0}{mL} de soude \SI{0,20}{mol.L^{-1}} (demi-neutralisation : $n_{\ce{HO-}} = \SI{0,010}{mol} = \frac{1}{2}n_{\ce{AH}}$).
  \item Noter le pH final : il doit être \textbf{4,8}.
\end{enumerate}
\textbf{Interprétation :} À la demi-équivalence, $n_{\ce{CH3COOH}} = n_{\ce{CH3COO-}}$, donc $\mathrm{pH} = \mathrm{p}K_a$.
\end{experience}

\courssub{4. Zone tampon et pouvoir tampon sur une courbe de dosage}

\begin{definition}[Zone tampon]
Lors du dosage d'un acide faible par une base forte (ou réciproque), la \textbf{zone tampon} est la portion de la courbe $\mathrm{pH} = f(V_b)$ située \textbf{autour de la demi-équivalence} (généralement pour $0,5 \le V_b/V_e \le 1,5$ ou pH = $\mathrm{p}K_a \pm 1$). Dans cette zone, la pente $\mathrm{dpH}/\mathrm{d}V_b$ est minimale.
\end{definition}

\begin{propriete}[Pouvoir tampon (capacité tampon)]
Le \textbf{pouvoir tampon} $\beta$ quantifie la résistance au changement de pH :
\[
\beta = \frac{\mathrm{d}n_{\ce{HO-}}}{\mathrm{dpH}} \quad (\text{ou } \frac{\mathrm{d}n_{\ce{H3O+}}}{\mathrm{dpH}})
\]
Il est \textbf{maximal} lorsque $[\ce{AH}] = [\ce{A-}]$ (pH = $\mathrm{p}K_a$) et \textbf{proportionnel à la concentration totale} $C_{\text{tot}} = [\ce{AH}] + [\ce{A-}]$. Plus le tampon est concentré, plus il résiste aux additions d'acide/base.
\end{propriete}

\begin{attention}[Limites du pouvoir tampon]
\begin{itemize}
  \item Le tampon \textbf{sature} : si on ajoute plus d'acide (ou base) qu'il n'y a de base (ou d'acide) dans le tampon, le pH chute (ou monte) brutalement.
  \item La dilution \textbf{extrême} diminue le pouvoir tampon (concentrations trop faibles, effets de force ionique, autoprotolyse de l'eau non négligeable).
  \item Un tampon n'est efficace que dans une gamme de pH d'environ $\mathrm{p}K_a \pm 1$.
\end{itemize}
\end{attention}

\begin{popfigure}
\begin{tikzpicture}[scale=0.8, domain=0:20]
\begin{axis}[
    width=\linewidth, height=0.6\linewidth,
    xlabel={Volume de soude $V_b$ (mL)}, ylabel={pH},
    xmin=0, xmax=20, ymin=0, ymax=14,
    xtick={0,5,10,15,20}, ytick={0,2,4,6,8,10,12,14},
    grid=major, grid style={popline},
    axis lines=middle, enlargelimits=0.05,
    clip=false
]
% Courbe sigmoïde approximative pour acide faible/base forte (pKa=4.8, Ca=0.1, Cb=0.1, Va=25 -> Ve=25)
% Utilisation d'une fonction tangente hyperbolique décalée pour la forme
%\addplot[popblue, thick, samples=300] ({10+12.5*tanh(1.5*(x-10))}, {7+3.5*tanh(1.5*(x-10))}); % Fake curve, better to use analytical approx
% Let's use a more realistic analytical shape: pH = pKa + log((Cb*x)/(Ca*Va - Cb*x)) before eq
% For simplicity in TikZ, we draw a schematic curve.
\draw[popblue, very thick] plot[smooth, tension=0.8] coordinates {
    (0,2.9) (2,3.8) (4,4.3) (5,4.8) (6,5.2) (8,5.8) (9,6.5) (9.5,7.5) (10,8.7) (10.5,10.5) (11,11.2) (12,11.8) (14,12.2) (20,12.5)
};
% Zone tampon highlight
\fill[popblue!15] (4,0) rectangle (8,14);
\draw[popblue, dashed, thick] (4,0) -- (4,14) node[above, pos=0.9] {$V_e/2$};
\draw[popblue, dashed, thick] (8,0) -- (8,14);
\node[popblue, align=center] at (6, 13) {\textbf{Zone tampon}\\$\mathrm{p}K_a \pm 1$};
% Equivalence
\draw[poporange, dashed, thick] (10,0) -- (10,14) node[above, pos=0.9] {$V_e$};
% pKa line
\draw[poppurple, dashed] (0,4.8) -- (20,4.8) node[right] {$\mathrm{pH}=\mathrm{p}K_a$};
\end{axis}
\end{tikzpicture}
\legende{Courbe de dosage d'un acide faible par une base forte : zone tampon (bleu) autour de la demi-équivalence ($V_e/2$), $\mathrm{pH}=\mathrm{p}K_a$ (violet), saut à l'équivalence $V_e$ (orange).}
\end{popfigure}

\courssub{5. Applications des solutions tampons}

\begin{savaistu}[Applications au Cameroun et dans le monde]
\begin{itemize}
  \item \textbf{Biologie/Médecine :} Sang (\ce{H2CO3}/\ce{HCO3-}, pH 7,4), liquide céphalo-rachidien, urines. Solutions pour perfusion (Ringer lactate, sérum physiologique tamponné phosphate).
  \item \textbf{Industrie agroalimentaire :} Stabilisation du pH des jus (bissap, ananas), boissons gazeuses, bière (bili-bili, vin de palme), conserves. Contrôle de la fermentation.
  \item \textbf{Industrie pharmaceutique/cosmétique :} Crèmes, shampoings, collyres, médicaments injectables (stabilité du principe actif, non-irritation). pH peau $\approx$ 5,5 (film hydrolipidique).
  \item \textbf{Industrie chimique (SONARA, CIMENCAM, Savonneries) :} Traitement des eaux, contrôle des bains galvaniques, fabrication de savons (hydrolyse des esters, contrôle pH pâte à savon), lessives.
  \item \textbf{Laboratoire :} Étalonnage des pH-mètres (solutions tampons étalons pH 4,01 ; 7,00 ; 9,21 ; 10,00), milieux de réaction enzymatiques, dosages.
\end{itemize}
\end{savaistu}

\courssub{Je vérifie mes connaissances}

\begin{exercice}[QCM : les solutions tampons]
Pour chaque question, choisir la (ou les) bonne(s) réponse(s) et justifier brièvement.
\begin{enumerate}
  \item Une solution tampon est une solution dont le pH :
  \begin{enumerate}
    \item est toujours égal à 7 ;
    \item varie peu par addition modérée d'un acide ou d'une base, ou par dilution modérée ;
    \item ne varie jamais, quelle que soit la quantité d'acide ajoutée ;
    \item est toujours égal au $\mathrm{p}K_a$ du couple utilisé.
  \end{enumerate}
  \item Une solution tampon peut être préparée en mélangeant :
  \begin{enumerate}
    \item un acide fort et une base forte ;
    \item un acide faible et sa base conjuguée en quantités voisines ;
    \item un acide faible et un excès de base forte ;
    \item une base faible et un acide fort en quantité deux fois moindre (demi-neutralisation).
  \end{enumerate}
  \item Sur la courbe de dosage d'un acide faible par une base forte, la zone tampon se situe :
  \begin{enumerate}
    \item autour de l'équivalence ;
    \item autour de la demi-équivalence ;
    \item au début du dosage ;
    \item après l'équivalence.
  \end{enumerate}
  \item Le pouvoir tampon est maximal lorsque :
  \begin{enumerate}
    \item $[\text{base}] \gg [\text{acide}]$ ;
    \item $[\text{acide}] = [\text{base}]$, c'est-à-dire pour $\mathrm{pH} = \mathrm{p}K_a$ ;
    \item la solution est très diluée ;
    \item le pH est égal à 7.
  \end{enumerate}
\end{enumerate}
\end{exercice}

\begin{exercice}[Vrai ou faux]
Répondre par vrai ou faux en justifiant chaque réponse.
\begin{enumerate}
  \item Le pH d'une solution tampon ne dépend pas de sa dilution, tant que la dilution reste modérée.
  \item Une solution d'acide chlorhydrique de concentration \SI{0,1}{mol.L^{-1}} constitue une bonne solution tampon.
  \item À la demi-équivalence du dosage d'un acide faible par la soude, on a $\mathrm{pH} = \mathrm{p}K_a$.
  \item Plus les concentrations de l'acide faible et de sa base conjuguée sont élevées, plus le pouvoir tampon est grand.
  \item Le sang humain est tamponné par le couple \ce{H2CO3}/\ce{HCO3-}.
\end{enumerate}
\end{exercice}

\begin{exercice}[Texte à trous]
Compléter le texte suivant avec les mots ou expressions : \emph{demi-équivalence} ; \emph{pKa} ; \emph{base conjuguée} ; \emph{modérée} ; \emph{pouvoir tampon} ; \emph{équimolaire}.

Une solution tampon contient un acide faible et sa \ldots\ldots\ldots\ldots\ldots en concentrations voisines. Son pH varie peu par addition \ldots\ldots\ldots\ldots\ldots d'acide, de base ou par dilution. Le mélange est dit \ldots\ldots\ldots\ldots\ldots lorsque les deux espèces sont en quantités égales : le pH est alors égal au \ldots\ldots\ldots\ldots\ldots du couple. Cette situation correspond, sur une courbe de dosage, à la \ldots\ldots\ldots\ldots\ldots. La capacité de la solution à s'opposer aux variations de pH est appelée \ldots\ldots\ldots\ldots\ldots.
\end{exercice}

\begin{exercice}[Définitions]
\begin{enumerate}
  \item Définir une solution tampon et citer ses trois propriétés fondamentales.
  \item Définir le pouvoir tampon d'une solution. De quels facteurs dépend-il ?
  \item Citer deux méthodes de préparation d'une solution tampon de $\mathrm{pH} = \mathrm{p}K_a$.
  \item Donner deux exemples d'utilisation des solutions tampons, l'un en biologie, l'autre au laboratoire.
\end{enumerate}
\end{exercice}

\courssub{Je m'entraîne}

\begin{exercice}[pH d'un mélange acide / base conjuguée]
On mélange $V_a = \SI{50,0}{mL}$ d'une solution d'acide éthanoïque \ce{CH3COOH} de concentration $C_a = \SI{0,20}{mol.L^{-1}}$ avec $V_b = \SI{50,0}{mL}$ d'une solution d'éthanoate de sodium \ce{CH3COO- + Na+} de concentration $C_b = \SI{0,10}{mol.L^{-1}}$.
\par Donnée : $\mathrm{p}K_a(\ce{CH3COOH}/\ce{CH3COO-}) = 4,8$.
\begin{enumerate}
  \item Calculer les quantités de matière d'acide éthanoïque et d'ion éthanoate dans le mélange.
  \item En déduire le pH du mélange.
  \item Ce mélange est-il une solution tampon ? Justifier.
\end{enumerate}
\end{exercice}

\begin{exercice}[Préparation d'un tampon par demi-neutralisation]
Un élève de Terminale C du lycée de Nkongsamba veut préparer un tampon de $\mathrm{pH} = 4,8$ à partir de $V_a = \SI{100}{mL}$ d'une solution d'acide éthanoïque de concentration $C_a = \SI{0,20}{mol.L^{-1}}$ ($\mathrm{p}K_a = 4,8$), en y ajoutant un volume $V_b$ de soude de concentration $C_b = \SI{0,20}{mol.L^{-1}}$.
\begin{enumerate}
  \item Écrire l'équation de la réaction entre l'acide éthanoïque et les ions hydroxyde.
  \item Montrer que pour obtenir $\mathrm{pH} = \mathrm{p}K_a$, il faut neutraliser exactement la moitié de l'acide.
  \item Calculer le volume $V_b$ de soude à verser.
\end{enumerate}
\end{exercice}

\begin{exercice}[Effet d'un ajout d'acide sur un tampon]
On dispose de $\SI{100}{mL}$ d'une solution tampon contenant $n_a = \SI{0,010}{mol}$ d'acide éthanoïque et $n_b = \SI{0,010}{mol}$ d'ions éthanoate ($\mathrm{p}K_a = 4,8$). On y ajoute $V = \SI{2,0}{mL}$ d'une solution d'acide chlorhydrique de concentration $C = \SI{0,10}{mol.L^{-1}}$ (on négligera la variation de volume).
\begin{enumerate}
  \item Calculer le pH du tampon avant l'ajout.
  \item Écrire l'équation de la réaction qui se produit lors de l'ajout de \ce{H3O+}.
  \item Dresser un bilan de matière et calculer le pH après l'ajout.
  \item Calculer le pH qu'aurait obtenu le même ajout d'acide dans $\SI{100}{mL}$ d'eau pure. Conclure sur l'efficacité du tampon.
\end{enumerate}
\end{exercice}

\begin{exercice}[Choix d'un couple pour un tampon]
Un technicien du laboratoire d'une savonnerie de Douala doit préparer une solution tampon de $\mathrm{pH} = 9,2$. Il dispose des couples acide/base suivants :
\begin{center}
\begin{tabularx}{\linewidth}{|l|X|}
\hline
\rowcolor{popblueL} Couple acide/base & $\mathrm{p}K_a$ \\
\hline
\ce{CH3COOH}/\ce{CH3COO-} & 4,8 \\
\hline
\ce{H2PO4-}/\ce{HPO4^2-} & 7,2 \\
\hline
\ce{NH4+}/\ce{NH3} & 9,2 \\
\hline
\ce{HCO3-}/\ce{CO3^2-} & 10,3 \\
\hline
\end{tabularx}
\end{center}
\begin{enumerate}
  \item Quel couple faut-il choisir ? Justifier.
  \item Il décide de mélanger un volume $V_a$ d'une solution de chlorure d'ammonium \ce{NH4+ + Cl-} à \SI{0,10}{mol.L^{-1}} avec un volume $V_b$ d'une solution d'ammoniac \ce{NH3} à \SI{0,10}{mol.L^{-1}}, de façon à obtenir $\SI{200}{mL}$ de tampon. Calculer $V_a$ et $V_b$.
  \item Le tampon obtenu résistera-t-il mieux à l'ajout d'acide ou à l'ajout de base ? Justifier.
\end{enumerate}
\end{exercice}

\courssub{Je me prépare au Bac}

\begin{exercice}[Type Bac — Préparation et étude d'un tampon acétique]
Au laboratoire du Lycée Classique de Dschang, on prépare une solution tampon en mélangeant $V_a = \SI{40,0}{mL}$ d'une solution d'acide éthanoïque \ce{CH3COOH} de concentration $C_a = \SI{0,15}{mol.L^{-1}}$ et $V_b = \SI{60,0}{mL}$ d'une solution d'éthanoate de sodium de concentration $C_b = \SI{0,10}{mol.L^{-1}}$.
\par Données : $\mathrm{p}K_a(\ce{CH3COOH}/\ce{CH3COO-}) = 4,8$ ; $M(\ce{CH3COOH}) = \SI{60}{g.mol^{-1}}$ ; $M(\ce{CH3COONa}) = \SI{82}{g.mol^{-1}}$.
\begin{enumerate}
  \item Rappeler la définition d'une solution tampon.
  \item Calculer les quantités de matière $n_a$ d'acide éthanoïque et $n_b$ d'ions éthanoate introduites.
  \item Montrer que le pH du mélange vaut $\mathrm{pH} = 4,8$.
  \item Justifier que ce mélange constitue une solution tampon et préciser dans quelle gamme de pH son pouvoir tampon est efficace.
  \item On ajoute à ce mélange $n = \SI{1,0e-3}{mol}$ d'ions hydroxyde \ce{HO-} (sans variation de volume).
  \begin{enumerate}
    \item Écrire l'équation de la réaction qui se produit.
    \item Calculer le nouveau pH du mélange et la variation de pH correspondante. Conclure.
  \end{enumerate}
  \item Quelle masse d'éthanoate de sodium solide aurait-il fallu peser pour préparer les $\SI{60,0}{mL}$ de solution utilisés ?
\end{enumerate}
\end{exercice}

\begin{exercice}[Type Bac — Dosage de l'acide d'un vinaigre et zone tampon]
Une coopérative de la région de l'Ouest produit du vinaigre artisanal. Pour contrôler sa qualité, on dose $V_a = \SI{25,0}{mL}$ d'une solution de vinaigre dilué (acide éthanoïque \ce{CH3COOH}) par une solution de soude de concentration $C_b = \SI{0,10}{mol.L^{-1}}$. Le suivi pH-métrique donne les résultats suivants :
\begin{center}
\begin{tabularx}{\linewidth}{|l|X|X|X|X|X|X|X|X|X|X|X|X|}
\hline
\rowcolor{popblueL} $V_b$ (mL) & 0 & 2 & 4 & 5 & 6 & 8 & 9 & 9,5 & 10 & 10,5 & 11 & 12 \\
\hline pH & 2,9 & 3,8 & 4,3 & 4,8 & 5,2 & 5,8 & 6,5 & 7,5 & 8,7 & 10,5 & 11,2 & 11,7 \\
\hline
\end{tabularx}
\end{center}
\begin{center}
\begin{tabularx}{\linewidth}{|l|X|X|X|X|}
\hline
\rowcolor{popblueL} $V_b$ (mL) & 14 & 16 & 18 & 20 \\
\hline pH & 12,0 & 12,2 & 12,3 & 12,4 \\
\hline
\end{tabularx}
\end{center}
\begin{enumerate}
  \item Écrire l'équation de la réaction de dosage.
  \item Tracer l'allure de la courbe $\mathrm{pH} = f(V_b)$ sur une feuille (échelle conseillée : $\SI{1}{cm}$ pour $\SI{2}{mL}$ et $\SI{1}{cm}$ pour 1 unité de pH).
  \item Déterminer graphiquement le volume équivalent $V_e$ et en déduire la concentration $C_a$ de l'acide éthanoïque dans le vinaigre dilué.
  \item Définir la demi-équivalence. Déterminer graphiquement le pH à la demi-équivalence et en déduire le $\mathrm{p}K_a$ du couple \ce{CH3COOH}/\ce{CH3COO-}.
  \item Délimiter sur la courbe la zone tampon. Pourquoi le pH varie-t-il peu dans cette zone ?
  \item Pourquoi la phénolphtaléine (zone de virage : 8,2 -- 10,0) convient-elle comme indicateur coloré pour ce dosage ?
\end{enumerate}
\end{exercice}

\begin{exercice}[Type Bac — Le tampon sanguin]
Le sang humain doit maintenir son pH entre 7,35 et 7,45 pour que l'organisme fonctionne correctement ; c'est le rôle du couple \ce{H2CO3}/\ce{HCO3-} (dioxyde de carbone dissous / ion hydrogénocarbonate).
\par Données : $\mathrm{p}K_a(\ce{H2CO3}/\ce{HCO3-}) = 6,4$ ; pH normal du sang : 7,4 ; concentration totale $[\ce{H2CO3}] + [\ce{HCO3-}] = \SI{2,5e-2}{mol.L^{-1}}$.
\begin{enumerate}
  \item Écrire l'équation de la réaction de l'acide \ce{H2CO3} avec l'eau et donner l'expression de la constante d'acidité $K_a$ du couple.
  \item Montrer que le rapport des concentrations dans le sang vaut $\dfrac{[\ce{HCO3-}]}{[\ce{H2CO3}]} = 10$.
  \item En déduire les valeurs de $[\ce{HCO3-}]$ et $[\ce{H2CO3}]$ dans le sang.
  \item Expliquer pourquoi ce mélange constitue une solution tampon efficace alors que le pH du sang (7,4) est éloigné du $\mathrm{p}K_a$ (6,4). On s'appuiera sur la notion de pouvoir tampon.
  \item Lors d'un effort intense, les muscles produisent de l'acide lactique \ce{CH3-CHOH-COOH}, qui libère des ions \ce{H3O+} dans le sang.
  \begin{enumerate}
    \item Écrire l'équation de la réaction entre les ions \ce{H3O+} et l'espèce du tampon qui les consomme.
    \item Expliquer qualitativement pourquoi le pH sanguin varie alors très peu.
  \end{enumerate}
  \item L'hyperventilation (respiration rapide et profonde) élimine le \ce{CO2} dissous, c'est-à-dire fait diminuer $[\ce{H2CO3}]$. Dans quel sens varie alors le pH du sang ? Justifier à l'aide de la relation $\mathrm{pH} = \mathrm{p}K_a + \log\dfrac{[\ce{HCO3-}]}{[\ce{H2CO3}]}$.
\end{enumerate}
\end{exercice}



\newpage

\coursec{Corrigés des exercices}

\begin{corrige}[Exercice 1 — QCM : les solutions tampons]
\begin{enumerate}
  \item \textbf{Bonne réponse : b.} C'est la définition même d'une solution tampon. \\
  a : faux, un tampon peut avoir n'importe quel pH (4,8 ; 7,4 ; 9,2\ldots), pas nécessairement 7. \\
  c : faux, le tampon \emph{sature} si l'ajout d'acide ou de base dépasse sa capacité. \\
  d : faux, on a $\mathrm{pH} = \mathrm{p}K_a$ seulement pour un mélange équimolaire.
  \item \textbf{Bonnes réponses : b et d.} \\
  b : méthode directe, le couple faible \ce{AH}/\ce{A-} est présent en quantités voisines. \\
  d : demi-neutralisation d'une base faible par un acide fort : $n_{\ce{H3O+}} = \tfrac{1}{2}n_{\ce{B},0}$ donne $n_{\ce{B}} = n_{\ce{BH+}}$, donc $\mathrm{pH} = \mathrm{p}K_a$. \\
  a : faux, un acide fort et une base forte donnent un sel neutre sans couple faible. \\
  c : faux, l'excès de base forte impose le pH ; il ne reste plus d'acide faible pour jouer le rôle de réserve.
  \item \textbf{Bonne réponse : b.} La zone tampon se situe autour de la demi-équivalence ($V_b = V_e/2$), où $\mathrm{pH} = \mathrm{p}K_a$ et où la pente de la courbe est minimale.
  \item \textbf{Bonne réponse : b.} Le pouvoir tampon $\beta$ est maximal pour $[\ce{AH}] = [\ce{A-}]$, c'est-à-dire $\mathrm{pH} = \mathrm{p}K_a$, et il croît avec la concentration totale.
\end{enumerate}
\end{corrige}

\begin{corrige}[Exercice 2 — Vrai ou faux]
\begin{enumerate}
  \item \textbf{Vrai.} $\mathrm{pH} = \mathrm{p}K_a + \log\dfrac{n_{\ce{A-}}}{n_{\ce{AH}}}$ : la dilution divise les deux concentrations par le même facteur, le rapport reste inchangé, donc le pH aussi (tant que la dilution reste modérée).
  \item \textbf{Faux.} \ce{HCl} est un acide \emph{fort}, totalement dissocié : il n'existe pas de couple faible \ce{AH}/\ce{A-} en quantités voisines. Tout ajout d'acide ou de base fait varier brutalement le pH.
  \item \textbf{Vrai.} À la demi-équivalence, la moitié de l'acide a été consommée : $n_{\ce{AH}} = n_{\ce{A-}}$, donc $\mathrm{pH} = \mathrm{p}K_a + \log 1 = \mathrm{p}K_a$.
  \item \textbf{Vrai.} Le pouvoir tampon est proportionnel à la concentration totale $C_{\text{tot}} = [\ce{AH}] + [\ce{A-}]$ : plus les réserves d'acide et de base sont grandes, plus le tampon absorbe d'ajouts sans variation notable de pH.
  \item \textbf{Vrai.} Le couple \ce{H2CO3}/\ce{HCO3-} (système bicarbonate) est le principal tampon du sang, couplé à l'action des poumons (élimination de \ce{CO2}) et des reins (régulation de \ce{HCO3-}).
\end{enumerate}
\end{corrige}

\begin{corrige}[Exercice 3 — Texte à trous]
Une solution tampon contient un acide faible et sa \textbf{base conjuguée} en concentrations voisines. Son pH varie peu par addition \textbf{modérée} d'acide, de base ou par dilution. Le mélange est dit \textbf{équimolaire} lorsque les deux espèces sont en quantités égales : le pH est alors égal au \textbf{pKa} du couple. Cette situation correspond, sur une courbe de dosage, à la \textbf{demi-équivalence}. La capacité de la solution à s'opposer aux variations de pH est appelée \textbf{pouvoir tampon}.
\end{corrige}

\begin{corrige}[Exercice 4 — Définitions]
\begin{enumerate}
  \item Une \textbf{solution tampon} est une solution dont le pH varie très peu lors de l'addition modérée d'un acide ou d'une base, ou lors d'une dilution modérée. Elle contient un couple acide faible / base conjuguée en concentrations voisines. Ses trois propriétés : (1) stabilité du pH par ajout d'acide (les \ce{H3O+} sont consommés par la base du couple) ; (2) stabilité par ajout de base (les \ce{HO-} sont consommés par l'acide du couple) ; (3) stabilité par dilution (le rapport $[\ce{A-}]/[\ce{AH}]$ est invariant).
  \item Le \textbf{pouvoir tampon} $\beta$ est la quantité de matière d'acide ou de base forte qu'il faut ajouter à \SI{1}{L} de solution pour faire varier son pH d'une unité. Il dépend : de la concentration totale du couple ($\beta$ croît avec $C_{\text{tot}}$) ; du rapport $[\ce{A-}]/[\ce{AH}]$ (maximal pour 1, c'est-à-dire $\mathrm{pH} = \mathrm{p}K_a$) ; de la nature du couple (gamme d'efficacité $\mathrm{p}K_a \pm 1$).
  \item (1) Mélanger des volumes égaux de solutions d'acide faible et de sel de sa base conjuguée de même concentration (mélange équimolaire). (2) Neutraliser exactement la moitié d'un acide faible par une base forte (demi-équivalence). On peut aussi citer la demi-neutralisation d'une base faible par un acide fort.
  \item Biologie : le sang, tamponné par le couple \ce{H2CO3}/\ce{HCO3-} (pH $\approx$ 7,4). Laboratoire : l'étalonnage des pH-mètres avec des solutions tampons étalons (pH 4,01 ; 7,00 ; 9,21).
\end{enumerate}
\end{corrige}

\begin{corrige}[Exercice 5 — pH d'un mélange acide éthanoïque / ion éthanoate]
\begin{enumerate}
  \item Quantités de matière introduites :
  \[
  n_{\ce{CH3COOH}} = C_a V_a = \SI{0,20}{mol.L^{-1}} \times \SI{0,0500}{L} = \SI{1,00e-2}{mol}
  \]
  \[
  n_{\ce{CH3COO-}} = C_b V_b = \SI{0,10}{mol.L^{-1}} \times \SI{0,0500}{L} = \SI{5,0e-3}{mol}
  \]
  \item Le mélange contient le couple \ce{CH3COOH}/\ce{CH3COO-} (acide éthanoïque / ion éthanoate) ; on applique Henderson-Hasselbalch avec les quantités de matière (même volume total) :
  \[
  \mathrm{pH} = \mathrm{p}K_a + \log\frac{n_{\ce{CH3COO-}}}{n_{\ce{CH3COOH}}} = 4,8 + \log\frac{5,0\times10^{-3}}{1,00\times10^{-2}} = 4,8 + \log(0,50)
  \]
  \[
  \boxed{\mathrm{pH} = 4,8 - 0,30 = 4,50}
  \]
  \item \textbf{Oui}, ce mélange est une solution tampon : il contient un acide faible et sa base conjuguée en quantités comparables, le rapport $n_{\ce{A-}}/n_{\ce{AH}} = 0,50$ étant compris entre 0,1 et 10. Il résistera aux ajouts modérés d'acide, de base et à la dilution.
\end{enumerate}
\end{corrige}

\begin{corrige}[Exercice 6 — Préparation d'un tampon éthanoïque par demi-neutralisation]
\begin{enumerate}
  \item Réaction (totale) entre l'acide éthanoïque et la soude :
  \[ \ce{CH3COOH + HO- -> CH3COO- + H2O} \]
  \item Pour obtenir $\mathrm{pH} = \mathrm{p}K_a$, il faut $\log\dfrac{n_{\ce{CH3COO-}}}{n_{\ce{CH3COOH}}} = 0$, soit $n_{\ce{CH3COO-}} = n_{\ce{CH3COOH}}$. \\
  Quantité initiale d'acide : $n_0 = C_a V_a = \SI{0,20}{mol.L^{-1}} \times \SI{0,100}{L} = \SI{2,00e-2}{mol}$. \\
  Après ajout de $n$ mol de \ce{HO-} : $n_{\ce{CH3COOH}} = n_0 - n$ et $n_{\ce{CH3COO-}} = n$. \\
  La condition $n_0 - n = n$ donne $n = \dfrac{n_0}{2}$ : il faut neutraliser exactement la \textbf{moitié} de l'acide.
  \item $n_{\ce{HO-}} = \dfrac{2,00\times10^{-2}}{2} = \SI{1,00e-2}{mol}$.
  \[
  V_b = \frac{n_{\ce{HO-}}}{C_b} = \frac{\SI{1,00e-2}{mol}}{\SI{0,20}{mol.L^{-1}}} = \SI{5,00e-2}{L} = \boxed{\SI{50,0}{mL}}
  \]
  Vérification : $n_{\ce{CH3COOH}} = n_{\ce{CH3COO-}} = \SI{1,00e-2}{mol}$, donc $\mathrm{pH} = 4,8$. Le tampon est prêt.
\end{enumerate}
\end{corrige}

\begin{corrige}[Exercice 7 — Effet d'un ajout d'acide sur un tampon éthanoïque]
\begin{enumerate}
  \item Avant l'ajout : $n_{\ce{AH}} = n_{\ce{A-}} = \SI{0,010}{mol}$, mélange équimolaire :
  \[
  \mathrm{pH} = 4,8 + \log(1) = \boxed{4,80}
  \]
  \item Les ions \ce{H3O+} apportés par \ce{HCl} sont consommés par la base du tampon (réaction totale) :
  \[ \ce{CH3COO- + H3O+ -> CH3COOH + H2O} \]
  \item Quantité d'acide ajoutée : $n_{\ce{H3O+}} = C V = \SI{0,10}{mol.L^{-1}} \times \SI{2,0e-3}{L} = \SI{2,0e-4}{mol}$ (réactif limitant).
  \begin{center}
  \begin{tabular}{|l|c|c|c|}
  \hline
  \rowcolor{popblueL} & \ce{CH3COO-} & \ce{H3O+} & \ce{CH3COOH} \\
  \hline État initial (mol) & $1,00\times10^{-2}$ & $2,0\times10^{-4}$ & $1,00\times10^{-2}$ \\
  \hline Variation (mol) & $-2,0\times10^{-4}$ & $-2,0\times10^{-4}$ & $+2,0\times10^{-4}$ \\
  \hline État final (mol) & $9,8\times10^{-3}$ & $\approx 0$ & $1,02\times10^{-2}$ \\
  \hline
  \end{tabular}
  \end{center}
  \[
  \mathrm{pH}_{\text{final}} = 4,8 + \log\frac{9,8\times10^{-3}}{1,02\times10^{-2}} = 4,8 + \log(0,961) = \boxed{4,78}
  \]
  La variation n'est que de $\Delta\mathrm{pH} = -0,02$ unité.
  \item Dans \SI{100}{mL} d'eau pure (volume final $\approx \SI{0,102}{L}$) :
  \[
  [\ce{H3O+}] = \frac{2,0\times10^{-4}}{0,102} = \SI{1,96e-3}{mol.L^{-1}} \Rightarrow \mathrm{pH} = -\log(1,96\times10^{-3}) \approx \boxed{2,71}
  \]
  \textbf{Conclusion :} dans l'eau pure, le pH chuterait de 7,0 à 2,7 ($\Delta\mathrm{pH} \approx -4,3$) ; dans le tampon, il ne baisse que de 0,02. Le tampon amortit spectaculairement l'ajout d'acide.
\end{enumerate}
\end{corrige}

\begin{corrige}[Exercice 8 — Choix d'un couple pour un tampon de pH 9,2]
\begin{enumerate}
  \item Un tampon est optimal lorsque $\mathrm{pH} = \mathrm{p}K_a$. Pour un tampon de pH 9,2, il faut donc choisir le couple \textbf{\ce{NH4+}/\ce{NH3}} (ion ammonium / ammoniac, $\mathrm{p}K_a = 9,2$). Les autres couples donneraient un tampon de pH très différent du $\mathrm{p}K_a$, donc de pouvoir tampon faible.
  \item Pour $\mathrm{pH} = \mathrm{p}K_a$, il faut un mélange équimolaire : $n_{\ce{NH4+}} = n_{\ce{NH3}}$. Les deux solutions mères ayant la même concentration (\SI{0,10}{mol.L^{-1}}), il faut des volumes égaux :
  \[
  V_a = V_b \quad \text{et} \quad V_a + V_b = \SI{200}{mL} \Rightarrow \boxed{V_a = V_b = \SI{100}{mL}}
  \]
  \item Le mélange étant équimolaire ($[\ce{NH4+}] = [\ce{NH3}]$), le pouvoir tampon est \textbf{symétrique} : la solution résiste \emph{également bien} aux deux types d'ajout. L'acide ajouté est consommé par \ce{NH3} (\ce{NH3 + H3O+ -> NH4+ + H2O}) et la base ajoutée par \ce{NH4+} (\ce{NH4+ + HO- -> NH3 + H2O}), les deux réserves étant identiques.
\end{enumerate}
\end{corrige}

\begin{corrige}[Exercice 9 — Type Bac : préparation et étude d'un tampon éthanoïque]
\begin{enumerate}
  \item Une solution tampon est une solution dont le pH varie très peu lors de l'addition modérée d'un acide ou d'une base, ou lors d'une dilution modérée. Elle contient un couple acide faible / base conjuguée en concentrations voisines.
  \item Quantités introduites :
  \[
  n_a = C_a V_a = \SI{0,15}{mol.L^{-1}} \times \SI{0,0400}{L} = \SI{6,0e-3}{mol}
  \]
  \[
  n_b = C_b V_b = \SI{0,10}{mol.L^{-1}} \times \SI{0,0600}{L} = \SI{6,0e-3}{mol}
  \]
  \item Le mélange est équimolaire ($n_a = n_b$) :
  \[
  \mathrm{pH} = \mathrm{p}K_a + \log\frac{n_b}{n_a} = 4,8 + \log(1) = \boxed{4,8}
  \]
  \item C'est bien un tampon : acide faible \ce{CH3COOH} (acide éthanoïque) et base conjuguée \ce{CH3COO-} (ion éthanoate) en quantités égales et non nulles. Son pouvoir tampon est efficace dans la gamme $\mathrm{p}K_a \pm 1$, soit \textbf{pH compris entre 3,8 et 5,8}.
  \item \begin{enumerate}
      \item \ce{CH3COOH + HO- -> CH3COO- + H2O} (réaction totale, \ce{HO-} limitant).
      \item Bilan de matière :
      \begin{center}
      \begin{tabular}{|l|c|c|}
      \hline
      \rowcolor{popblueL} & \ce{CH3COOH} & \ce{CH3COO-} \\
      \hline État initial (mol) & $6,0\times10^{-3}$ & $6,0\times10^{-3}$ \\
      \hline Variation (mol) & $-1,0\times10^{-3}$ & $+1,0\times10^{-3}$ \\
      \hline État final (mol) & $5,0\times10^{-3}$ & $7,0\times10^{-3}$ \\
      \hline
      \end{tabular}
      \end{center}
      \[
      \mathrm{pH}_{\text{final}} = 4,8 + \log\frac{7,0\times10^{-3}}{5,0\times10^{-3}} = 4,8 + \log(1,4) = \boxed{4,95}
      \]
      $\Delta\mathrm{pH} = +0,15$ seulement. \textbf{Conclusion :} l'ajout de \SI{1,0e-3}{mol} de base forte (qui aurait porté \SI{100}{mL} d'eau pure à $\mathrm{pH} \approx 12$) ne fait varier le pH du tampon que de 0,15 unité : le tampon joue pleinement son rôle.
    \end{enumerate}
  \item Masse d'éthanoate de sodium :
  \[
  m = n_b \times M = \SI{6,0e-3}{mol} \times \SI{82}{g.mol^{-1}} = \boxed{\SI{0,49}{g}}
  \]
\end{enumerate}
\textbf{Barème indicatif (sur 10) :} définition 1 ; calculs $n_a$, $n_b$ 2 ; pH 1,5 ; justification + gamme 1,5 ; équation 0,5 ; bilan + pH final + conclusion 2,5 ; masse 1. \\
\textbf{Points de vigilance :} convertir les volumes en litres ; ne pas oublier que le rapport des concentrations égale le rapport des quantités (même volume total) ; vérifier que \ce{HO-} est bien le réactif limitant ; arrondir à deux chiffres significatifs.
\end{corrige}

\begin{corrige}[Exercice 10 — Type Bac : dosage du vinaigre et zone tampon]
\begin{enumerate}
  \item Réaction de dosage (totale) :
  \[ \ce{CH3COOH + HO- -> CH3COO- + H2O} \]
  \item La courbe $\mathrm{pH} = f(V_b)$ est croissante, en forme de « S » : départ à pH 2,9, zone de faible pente autour de pH 4,8, saut vertical centré sur $V_b = \SI{10}{mL}$, puis plateau basique vers pH 12,4.
  \item Le saut de pH maximal se produit entre $V_b = 9,5$ et $\SI{10,5}{mL}$ : le volume équivalent est $\boxed{V_e = \SI{10,0}{mL}}$. \\
  À l'équivalence : $C_a V_a = C_b V_e$, donc :
  \[
  C_a = \frac{C_b V_e}{V_a} = \frac{\SI{0,10}{mol.L^{-1}} \times \SI{10,0}{mL}}{\SI{25,0}{mL}} = \boxed{\SI{0,040}{mol.L^{-1}}}
  \]
  \item La \textbf{demi-équivalence} est le point où $V_b = V_e/2 = \SI{5,0}{mL}$ : la moitié de l'acide éthanoïque a été neutralisée et $n_{\ce{CH3COOH}} = n_{\ce{CH3COO-}}$. \\
  Lecture graphique (confirmée par le tableau) : à $V_b = \SI{5,0}{mL}$, $\mathrm{pH} = 4,8$. Donc :
  \[
  \boxed{\mathrm{p}K_a(\ce{CH3COOH}/\ce{CH3COO-}) = 4,8}
  \]
  \item La zone tampon s'étend approximativement de $V_b = \SI{2}{mL}$ à $V_b = \SI{8}{mL}$ (pH entre 3,8 et 5,8, soit $\mathrm{p}K_a \pm 1$). Dans cette zone, le mélange contient des quantités comparables d'acide \ce{CH3COOH} et de base \ce{CH3COO-} : la soude ajoutée est consommée par l'acide restant, et le rapport $[\ce{CH3COO-}]/[\ce{CH3COOH}]$ varie lentement ; le pH, qui dépend du \emph{logarithme} de ce rapport, varie donc très peu.
  \item À l'équivalence, la solution contient l'ion éthanoate (base faible) : le pH vaut \textbf{8,7}, valeur basique. La zone de virage de la phénolphtaléine (8,2 -- 10,0) \textbf{contient ce pH} et se situe dans le saut de pH : le virage (incolore $\to$ rose) signale donc l'équivalence avec précision. L'indicateur convient parfaitement.
\end{enumerate}
\textbf{Barème indicatif (sur 10) :} équation 1 ; courbe soignée 2 ; $V_e$ et $C_a$ 2 ; demi-équivalence et $\mathrm{p}K_a$ 2 ; zone tampon et explication 2 ; choix de l'indicateur 1. \\
\textbf{Points de vigilance :} ne pas confondre demi-équivalence ($\mathrm{pH} = \mathrm{p}K_a$) et équivalence (saut de pH, pH basique ici) ; à l'équivalence d'un dosage acide faible / base forte, le pH n'est \emph{pas} 7 ; respecter l'échelle imposée pour le tracé.
\end{corrige}

\begin{corrige}[Exercice 11 — Type Bac : le tampon sanguin (couple acide carbonique / hydrogénocarbonate)]
\begin{enumerate}
  \item Réaction avec l'eau :
  \[ \ce{H2CO3 + H2O <=> H3O+ + HCO3-} \]
  \[
  K_a = \frac{[\ce{H3O+}]_{\text{eq}}\,[\ce{HCO3-}]_{\text{eq}}}{[\ce{H2CO3}]_{\text{eq}}}
  \]
  \item À partir de la relation de Henderson-Hasselbalch :
  \[
  \mathrm{pH} = \mathrm{p}K_a + \log\frac{[\ce{HCO3-}]}{[\ce{H2CO3}]}
  \Rightarrow \log\frac{[\ce{HCO3-}]}{[\ce{H2CO3}]} = 7,4 - 6,4 = 1
  \]
  \[
  \boxed{\frac{[\ce{HCO3-}]}{[\ce{H2CO3}]} = 10^{1} = 10}
  \]
  \item On résout le système : $[\ce{HCO3-}] = 10\,[\ce{H2CO3}]$ et $[\ce{HCO3-}] + [\ce{H2CO3}] = \SI{2,5e-2}{mol.L^{-1}}$.
  \[
  11\,[\ce{H2CO3}] = \SI{2,5e-2}{mol.L^{-1}} \Rightarrow \boxed{[\ce{H2CO3}] \approx \SI{2,3e-3}{mol.L^{-1}}}
  \]
  \[
  \boxed{[\ce{HCO3-}] \approx \SI{2,3e-2}{mol.L^{-1}}}
  \]
  \item Bien que $\mathrm{pH} \neq \mathrm{p}K_a$, le tampon reste efficace pour trois raisons : (1) le rapport 10 est encore dans la gamme d'efficacité (0,1 à 10) ; (2) la concentration totale est élevée (\SI{2,5e-2}{mol.L^{-1}}), ce qui garantit un pouvoir tampon important ; (3) c'est un tampon \emph{ouvert} : les poumons éliminent le \ce{CO2} et les reins régulent \ce{HCO3-}, ce qui restaure en permanence les concentrations. Le sang dispose ainsi d'une réserve de base (\ce{HCO3-}) dix fois supérieure, adaptée à la charge acide permanente du métabolisme.
  \item \begin{enumerate}
      \item Les ions \ce{H3O+} libérés par l'acide lactique sont consommés par la base du tampon :
      \[ \ce{H3O+ + HCO3- -> H2CO3 + H2O} \]
      (puis \ce{H2CO3} se décompose en \ce{CO2}, éliminé par la respiration).
      \item Les \ce{H3O+} ajoutés sont quasiment tous consommés : le rapport $[\ce{HCO3-}]/[\ce{H2CO3}]$ diminue légèrement, mais les quantités en présence sont grandes devant l'apport d'acide, donc $\log(\text{rapport})$ — et par suite le pH — varie très peu.
    \end{enumerate}
  \item L'hyperventilation élimine le \ce{CO2} dissous, donc fait \textbf{diminuer} $[\ce{H2CO3}]$. Le rapport $\dfrac{[\ce{HCO3-}]}{[\ce{H2CO3}]}$ \textbf{augmente}, son logarithme aussi, et par conséquent :
  \[
  \boxed{\text{le pH du sang augmente (alcalose respiratoire).}}
  \]
\end{enumerate}
\textbf{Barème indicatif (sur 10) :} équation + $K_a$ 1,5 ; rapport 10 1,5 ; concentrations 2 ; efficacité du tampon 2 ; réaction avec \ce{H3O+} + explication 2 ; sens de variation du pH 1. \\
\textbf{Points de vigilance :} ne pas conclure hâtivement « pH $\neq$ pKa donc mauvais tampon » — raisonner sur le rapport (0,1 à 10) et la concentration totale ; bien identifier que c'est \ce{HCO3-} (la base) qui piège les ions \ce{H3O+} ; dans la question 6, c'est le \emph{dénominateur} qui diminue, donc le rapport et le pH augmentent.
\end{corrige}

\newpage

\coursec{Fiche bilan}

\coursec{Les solutions tampons}

\noindent
Dans un laboratoire de contrôle qualité de la SONARA (Société Nationale de Raffinage) ou dans une brasserie artisanale de bili-bili à Douala, il est crucial de maintenir le pH d'un milieu stable malgré l'ajout d'acides ou de bases. Comment le sang humain maintient-il son pH à 7,4 alors que le métabolisme produit constamment du \ce{CO2} acide ? La réponse réside dans les \cle{solutions tampons}. Cette leçon en dévoile le fonctionnement, le calcul et les applications.

\begin{popfigure}
\centering
\begin{tikzpicture}[
  central/.style={circle, draw=popdark, fill=popblue!25, text=popdark, font=\bfseries\small, align=center, minimum size=2.6cm},
  branche/.style={rounded corners=4pt, draw=#1, fill=#1!15, text=popdark, font=\footnotesize, align=center, inner sep=4pt, minimum width=3.1cm},
  lien/.style={thick, draw=#1, -{Stealth[length=2mm]}}
]
\node[central] (C) {Solutions\\tampons};
\node[branche=popblue] (def) at (-5.4,2.6) {\textbf{Définition}\\pH variant peu par\\ajout d'acide, de base\\ou par dilution};
\node[branche=popgreen] (compo) at (-5.6,-0.2) {\textbf{Composition}\\couple \ce{AH}/\ce{A-}\\en quantités\\voisines};
\node[branche=poporange] (ph) at (-4.6,-3.0) {\textbf{pH du mélange}\\$\mathrm{pH}=\mathrm{p}K_a+\log\dfrac{[\ce{A-}]}{[\ce{AH}]}$};
\node[branche=poppurple] (prep) at (0,-3.9) {\textbf{Préparation}\\$\mathrm{pH}=\mathrm{p}K_a$ : mélange\\équimolaire ou\\demi-neutralisation};
\node[branche=popgold] (pouvoir) at (4.6,-3.0) {\textbf{Pouvoir tampon}\\$\beta=\left|\dfrac{\Delta n}{\Delta\mathrm{pH}}\right|$\\maximal à la\\demi-équivalence};
\node[branche=poppink] (zone) at (5.6,-0.2) {\textbf{Zone tampon}\\autour de $V_{E}/2$\\sur la courbe\\$\mathrm{pH}=f(V)$};
\node[branche=popgreen!60!black] (appli) at (5.4,2.6) {\textbf{Applications}\\sang ($\mathrm{pH}\approx7{,}4$),\\étalonnage pH-mètre,\\pharmacie, agroalimentaire};
\draw[lien=popblue] (C) -- (def);
\draw[lien=popgreen] (C) -- (compo);
\draw[lien=poporange] (C) -- (ph);
\draw[lien=poppurple] (C) -- (prep);
\draw[lien=popgold] (C) -- (pouvoir);
\draw[lien=poppink] (C) -- (zone);
\draw[lien=popgreen!60!black] (C) -- (appli);
\end{tikzpicture}
\legende{Carte mentale de la leçon : les solutions tampons.}
\end{popfigure}

\courssub{Définition et propriétés d'une solution tampon}

\begin{definition}[Solution tampon]
Une \cle{solution tampon} (ou solution tamponnée) est une solution dont le \cle{pH varie très peu} :
\begin{itemize}
  \item lors de l'addition d'une \emph{faible quantité} d'acide (ions \ce{H3O+}) ;
  \item lors de l'addition d'une \emph{faible quantité} de base (ions \ce{HO-}) ;
  \item lors d'une \emph{dilution modérée} (ajout d'eau).
\end{itemize}
\end{definition}

\begin{propriete}[Composition d'une solution tampon]
Une solution tampon contient obligatoirement un \cle{couple acide faible / base conjuguée} \ce{AH}/\ce{A-} (ou \ce{BH+}/\ce{B}) dont les deux espèces sont présentes en \cle{quantités de matière voisines} (ni l'une ni l'autre n'est négligeable devant l'autre).
\begin{itemize}
  \item L'acide faible \ce{AH} neutralise la base ajoutée.
  \item La base conjuguée \ce{A-} neutralise l'acide ajouté.
\end{itemize}
\end{propriete}

\begin{experience}[Mise en évidence qualitative du pouvoir tampon]
\textbf{Matériel :} 2 béchers, eau distillée, solution tampon acétate (\ce{CH3COOH}/\ce{CH3COO-}, pH $\approx$ 4,7), indicateur coloré (rouge de méthyle ou bleu de bromothymol), solution d'acide chlorhydrique \SI{0,1}{mol.L^{-1}}, solution de soude \SI{0,1}{mol.L^{-1}}, pipettes, baguette de verre.
\textbf{Protocole :}
\begin{enumerate}
  \item Dans le bécher 1 : verser \SI{50}{mL} d'eau + quelques gouttes d'indicateur. Noter la couleur.
  \item Dans le bécher 2 : verser \SI{50}{mL} de solution tampon acétate + mêmes gouttes d'indicateur. Noter la couleur (identique au bécher 1).
  \item Ajouter \SI{1}{mL} de \ce{HCl} \SI{0,1}{mol.L^{-1}} dans chaque bécher, agiter, observer.
  \item Recommencer avec \SI{1}{mL} de \ce{NaOH} \SI{0,1}{mol.L^{-1}} dans deux nouveaux béchers identiques.
\end{enumerate}
\textbf{Observations :}
\begin{itemize}
  \item Eau + \ce{HCl} : virage brutal de l'indicateur (pH passe de 7 à $\approx$ 3).
  \item Tampon + \ce{HCl} : couleur quasi inchangée (pH varie de $\approx$ 0,1 unité).
  \item Eau + \ce{NaOH} : virage brutal (pH passe de 7 à $\approx$ 11).
  \item Tampon + \ce{NaOH} : couleur quasi inchangée.
\end{itemize}
\textbf{Interprétation :} Le couple \ce{CH3COOH}/\ce{CH3COO-} réagit avec les ions ajoutés :
\[ \ce{CH3COO- + H3O+ -> CH3COOH + H2O} \quad \text{(acide ajouté)} \]
\[ \ce{CH3COOH + HO- -> CH3COO- + H2O} \quad \text{(base ajoutée)} \]
Les réactions sont totales, les quantités de \ce{CH3COOH} et \ce{CH3COO-} varient peu \emph{relativement} à leurs valeurs initiales, donc le rapport $[\ce{A-}]/[\ce{AH}]$ varie peu $\Rightarrow$ pH stable.
\end{experience}

\begin{exemplebox}[Reconnaître une solution tampon]
Parmi les mélanges suivants, lesquels constituent une solution tampon ?
\begin{enumerate}
  \item \SI{100}{mL} \ce{CH3COOH} \SI{0,1}{M} + \SI{100}{mL} \ce{CH3COONa} \SI{0,1}{M}
  \item \SI{100}{mL} \ce{HCl} \SI{0,1}{M} + \SI{100}{mL} \ce{NaCl} \SI{0,1}{M}
  \item \SI{100}{mL} \ce{NH3} \SI{0,1}{M} + \SI{50}{mL} \ce{HCl} \SI{0,1}{M}
  \item \SI{100}{mL} \ce{CH3COOH} \SI{0,1}{M} + \SI{100}{mL} \ce{NaOH} \SI{0,1}{M}
\end{enumerate}
\textbf{Analyse :}
\begin{enumerate}
  \item \textbf{Oui.} Couple \ce{CH3COOH}/\ce{CH3COO-}, $n_{\text{acide}} = n_{\text{base}} = \SI{10}{mmol}$. Quantités voisines.
  \item \textbf{Non.} \ce{HCl} est un acide fort, \ce{Cl-} est une base négligeable. Pas de couple faible/faible.
  \item \textbf{Oui.} Réaction \ce{NH3 + H3O+ -> NH4+}. $n(\ce{NH3})_{\text{init}} = \SI{10}{mmol}$, $n(\ce{H3O+}) = \SI{5}{mmol}$. Après réaction : $n(\ce{NH3}) = \SI{5}{mmol}$, $n(\ce{NH4+}) = \SI{5}{mmol}$. Couple \ce{NH4+}/\ce{NH3} en quantités égales.
  \item \textbf{Non.} Neutralisation totale. $n(\ce{CH3COOH}) = n(\ce{NaOH}) = \SI{10}{mmol}$. Produit : \ce{CH3COONa} seul (solution de base conjuguée, pas de couple).
\end{enumerate}
\end{exemplebox}

\courssub{pH d'un mélange acide faible / base conjuguée}

\noindent
Soit un mélange contenant un acide faible \ce{AH} et sa base conjuguée \ce{A-}. L'équilibre acido-basique est :
\[ \ce{AH + H2O <=> A- + H3O+} \qquad K_a = \frac{[\ce{A-}][\ce{H3O+}]}{[\ce{AH}]} \]
En prenant le logarithme décimal (en base 10) et en changeant de signe :
\[ -\log[\ce{H3O+}] = -\log K_a + \log\frac{[\ce{A-}]}{[\ce{AH}]} \]
On obtient la \cle{relation de Henderson-Hasselbalch} :

\begin{propriete}[Relation de Henderson-Hasselbalch]
Pour une solution tampon constituée du couple \ce{AH}/\ce{A-} :
\[ \boxed{\mathrm{pH} = \mathrm{p}K_a + \log\frac{[\ce{A-}]}{[\ce{AH}]} = \mathrm{p}K_a + \log\frac{n(\ce{A-})}{n(\ce{AH})}} \]
La seconde égalité tient au fait que les deux espèces sont dans le même volume de solution (le volume se simplifie).
\end{propriete}

\begin{methode}[Calculer le pH d'un mélange tampon]
\begin{enumerate}
  \item Identifier le couple \ce{AH}/\ce{A-} et sa constante \cle{p$K_a$} (donnée ou déduite de $K_a$).
  \item Calculer les \textbf{quantités de matière initiales} $n(\ce{AH})$ et $n(\ce{A-})$ apportées par chaque solution mère.
  \item S'il y a réaction entre les espèces apportées (ex: acide faible + base forte), dresser le \textbf{tableau d'avancement} pour déterminer les quantités \textbf{à l'équilibre} (ou après réaction totale) $n_{\text{eq}}(\ce{AH})$ et $n_{\text{eq}}(\ce{A-})$.
  \item Vérifier que les deux quantités sont \textbf{non nulles} (sinon ce n'est pas un tampon).
  \item Appliquer la formule : $\mathrm{pH} = \mathrm{p}K_a + \log\dfrac{n_{\text{eq}}(\ce{A-})}{n_{\text{eq}}(\ce{AH})}$.
\end{enumerate}
\end{methode}

\begin{exoresolu}[Calcul du pH d'un tampon acétate]
\textbf{Énoncé :} On mélange \SI{50,0}{mL} d'acide éthanoïque \ce{CH3COOH} \SI{0,20}{mol.L^{-1}} ($pK_a = 4,74$) et \SI{50,0}{mL} d'éthanoate de sodium \ce{CH3COONa} \SI{0,10}{mol.L^{-1}}. Calculer le pH de la solution obtenue.
\textbf{Solution détaillée :}
\begin{enumerate}
  \item Couple : \ce{CH3COOH}/\ce{CH3COO-}, $\mathrm{p}K_a = 4,74$.
  \item Quantités initiales :
  $n(\ce{CH3COOH}) = 0,050 \times 0,20 = \SI{10,0}{mmol}$
  $n(\ce{CH3COO-}) = 0,050 \times 0,10 = \SI{5,0}{mmol}$ (apporté par le sel, dissocié totalement).
  \item Pas de réaction entre les espèces apportées (acide faible + sel de sa base conjuguée). Quantités à l'équilibre = quantités initiales.
  \item Application de la formule :
  \[ \mathrm{pH} = 4,74 + \log\frac{5,0}{10,0} = 4,74 + \log(0,5) = 4,74 - 0,30 = \mathbf{4,44} \]
\end{enumerate}
\textbf{Résultat :} Le pH de la solution tampon est \textbf{4,44}.
\end{exoresolu}

\begin{attention}[Pièges fréquents]
\begin{itemize}
  \item \textbf{Oublier le tableau d'avancement} si on mélange acide faible et base forte (ou base faible et acide fort). Il faut calculer ce qui \emph{reste} après réaction.
  \item \textbf{Utiliser les concentrations initiales} au lieu des concentrations \emph{finales} (après mélange, volume total $V_{\text{tot}} = V_1 + V_2$). \emph{Astuce :} Utiliser directement les quantités de matière $n$ (le volume se simplifie).
  \item \textbf{Confondre $K_a$ et $pK_a$} : $\mathrm{pH} = \mathrm{p}K_a + \log(...)$, pas $K_a$.
  \item \textbf{Négliger l'autoprotolyse de l'eau} : La formule est valide si $[\ce{AH}]$ et $[\ce{A-}]$ sont $\gg 10^{-6}$ mol/L (généralement $> 10^{-3}$ M).
\end{itemize}
\end{attention}

\courssub{Préparation d'une solution tampon de pH = pKa}

\noindent
Pour obtenir un tampon de pH \emph{exactement égal} au $\mathrm{p}K_a$ du couple choisi, il faut que le rapport $[\ce{A-}]/[\ce{AH}] = 1$, c'est-à-dire $n(\ce{A-}) = n(\ce{AH})$. Deux méthodes classiques existent.

\begin{methode}[Préparer un tampon de pH = pKa (Méthode 1 : Mélange équimolaire)]
On dispose d'une solution d'acide faible \ce{AH} ($C_a$) et d'une solution de sa base conjuguée \ce{A-} (souvent un sel, $C_b$).
\begin{enumerate}
  \item Choisir le couple dont le $\mathrm{p}K_a$ est le plus proche du pH visé.
  \item Prendre des volumes $V_a$ et $V_b$ tels que : $C_a V_a = C_b V_b$ (quantités de matière égales).
  \item Mélanger dans une fiole jaugée, compléter au trait de jauge, homogénéiser.
  \item Vérifier le pH au pH-mètre (éventuel ajustement goutte à goutte avec acide/base fort).
\end{enumerate}
\end{methode}

\begin{methode}[Préparer un tampon de pH = pKa (Méthode 2 : Demi-neutralisation)]
On dispose d'une solution d'acide faible \ce{AH} ($C_a, V_a$) et d'une base forte \ce{HO-} ($C_b$).
\begin{enumerate}
  \item Calculer le volume $V_b$ de base forte pour neutraliser \textbf{la moitié} de l'acide :
  \[ n(\ce{HO-}) = \frac{1}{2} n(\ce{AH}) \quad \Rightarrow \quad C_b V_b = \frac{C_a V_a}{2} \]
  \item Verser $V_a$ d'acide faible dans un bécher, ajouter $V_b$ de base forte (sous agitation).
  \item Transférer en fiole jaugée, compléter, homogénéiser.
  \item Vérifier au pH-mètre.
\end{enumerate}
\emph{Remarque :} Même principe pour une base faible \ce{B} demi-neutralisée par un acide fort.
\end{methode}

\begin{exoresolu}[Préparation d'un tampon acétate pH = 4,74]
\textbf{Énoncé :} On souhaite préparer \SI{250,0}{mL} de solution tampon \ce{CH3COOH}/\ce{CH3COO-} de pH = 4,74 ($pK_a = 4,74$) à partir d'acide éthanoïque \SI{0,200}{mol.L^{-1}} et de soude \SI{0,100}{mol.L^{-1}}.
\textbf{Solution détaillée :}
\begin{enumerate}
  \item pH visé = $\mathrm{p}K_a$ $\Rightarrow$ il faut $n(\ce{CH3COOH}) = n(\ce{CH3COO-})$ dans le mélange final.
  \item On choisit la méthode 2 (demi-neutralisation). Soit $V_a$ le volume d'acide pris. Quantité d'acide initiale : $n_a = 0,200 \times V_a$.
  \item Quantité de soude à ajouter : $n_b = \frac{n_a}{2} = 0,100 \times V_a$.
  \item Volume de soude : $V_b = \frac{n_b}{C_b} = \frac{0,100 \times V_a}{0,100} = V_a$.
  \item Volume total final approximatif : $V_a + V_b = 2 V_a = \SI{250,0}{mL} \Rightarrow V_a = \SI{125,0}{mL}$.
  \item \textbf{Protocole :} Mesurer \SI{125,0}{mL} d'acide \SI{0,200}{M} (fiole jaugée ou pipette), verser dans un bécher. Ajouter \SI{125,0}{mL} de soude \SI{0,100}{M} (burette ou fiole). Transférer en fiole jaugée \SI{250}{mL}, compléter à l'eau distillée, homogénéiser.
  \item \textbf{Vérification :} $n(\ce{CH3COOH})_{\text{final}} = \frac{0,200 \times 125}{2} = \SI{12,5}{mmol}$ ; $n(\ce{CH3COO-})_{\text{final}} = \SI{12,5}{mmol}$. Rapport = 1 $\Rightarrow$ pH = 4,74.
\end{enumerate}
\end{exoresolu}

\begin{experience}[TP : Préparation et contrôle d'une solution tampon phosphate pH = 7,2]
\textbf{Objectif :} Préparer \SI{100}{mL} de tampon phosphate (\ce{H2PO4-}/\ce{HPO4^2-}, $\mathrm{p}K_{a2} = 7,21$) par mélange équimolaire.
\textbf{Matériel :} Solutions mères \ce{NaH2PO4} \SI{0,1}{M} et \ce{Na2HPO4} \SI{0,1}{M}, fiole jaugée \SI{100}{mL}, pipettes \SI{50}{mL}, pH-mètre étalonné.
\textbf{Protocole :}
\begin{enumerate}
  \item Pipeter \SI{50,0}{mL} de \ce{NaH2PO4} \SI{0,1}{M} (acide) dans la fiole.
  \item Pipeter \SI{50,0}{mL} de \ce{Na2HPO4} \SI{0,1}{M} (base) dans la fiole.
  \item Compléter à l'eau distillée au trait de jauge, boucher, agiter.
  \item Mesurer le pH au pH-mètre.
\end{enumerate}
\textbf{Résultats attendus :} pH mesuré $\approx 7,21$. Écart possible $\pm 0,05$ dû à l'activité ionique, température, pureté réactifs.
\textbf{Compétence :} Manipuler le matériel volumétrique, étalonner et utiliser un pH-mètre, rendre compte.
\end{experience}

\courssub{Zone tampon et pouvoir tampon sur une courbe de dosage}

\noindent
Lors du dosage d'un acide faible par une base forte (ou inversé), la courbe $\mathrm{pH} = f(V_{\text{versé}})$ présente une zone de pente faible : la \cle{zone tampon}.

\begin{propriete}[Demi-équivalence et zone tampon]
Lors du dosage d'un acide faible \ce{AH} par une base forte \ce{HO-} :
\begin{itemize}
  \item À la \cle{demi-équivalence} ($V = V_E/2$), la quantité d'acide neutralisé égale la quantité d'acide restant : $n(\ce{A-}) = n(\ce{AH})$.
  \item Le pH à la demi-équivalence vaut le $\mathrm{p}K_a$ de l'acide : $\mathrm{pH}_{V_E/2} = \mathrm{p}K_a$.
  \item La \cle{zone tampon} s'étend approximativement de $V_E/10$ à $9V_E/10$ (soit $\mathrm{p}K_a \pm 1$). Dans cette zone, le pH varie peu pour un volume ajouté donné.
\end{itemize}
\end{propriete}

\begin{definition}[Pouvoir tampon ($\beta$)]
Le \cle{pouvoir tampon} $\beta$ quantifie l'efficacité d'un tampon à résister au changement de pH :
\[ \beta = \left| \frac{\Delta n_{\text{acide/base fort ajouté}}}{\Delta \mathrm{pH}} \right| \quad (\text{en } \mathrm{mol \cdot pH^{-1} \cdot L^{-1}}) \]
\begin{itemize}
  \item $\beta$ est \textbf{maximal} lorsque $\mathrm{pH} = \mathrm{p}K_a$ (demi-équivalence, quantités égales).
  \item $\beta$ augmente avec la \textbf{concentration totale} du couple ($C_{\text{tot}} = [\ce{AH}] + [\ce{A-}]$).
  \item $\beta$ devient négligeable hors de la zone $\mathrm{p}K_a \pm 1$.
\end{itemize}
\end{definition}

\begin{popfigure}
\centering
\begin{tikzpicture}
\begin{axis}[
    width=\linewidth, height=8cm,
    axis lines=middle,
    xlabel={Volume de base forte versé $V$ (mL)}, ylabel={pH},
    xmin=0, xmax=22, ymin=0, ymax=14,
    xtick={5,10,15,20}, ytick={0,2,4,6,7,8,10,12,14},
    xticklabels={$V_E/2$, $V_E$, $3V_E/2$, $2V_E$},
    yticklabels={0,2,4,6,7,8,10,12,14},
    domain=0.1:21, samples=300,
    clip=false,
    enlargelimits=0.05,
    x label style={at={(axis description cs:0.5,-0.05)}, anchor=north},
    y label style={at={(axis description cs:-0.05,0.5)}, anchor=south},
]
% Courbe de dosage typique acide faible / base forte (simulation simplifiée)
% pH = pKa + log(V/(Ve-V)) approximatif pour la zone tampon, mais on trace une sigmoïde réaliste
\addplot[popcurve, very thick, domain=0.5:19.5] {7 + 6/(1+exp(-1.5*(x-10)))}; % Sigmoïde centrée sur Ve=10
% Zone tampon
\addplot[popgreen!50!popblue, fill opacity=0.2, domain=2:18, samples=100] {7 + 6/(1+exp(-1.5*(x-10)))} \closedcycle;
% Droite pH = pKa
\addplot[poporange, dashed, thick, domain=0:22] {7};
% Droite Ve
\addplot[popgray, dashed, domain=0:14] coordinates {(10,0) (10,14)};
% Annotations
\node[poporange, anchor=south west] at (2,7.2) {$\mathrm{pH} = \mathrm{p}K_a$};
\node[popgray, anchor=south] at (10, -0.5) {$V_E$};
\node[popgreen!50!popblue, anchor=north] at (10, 13.5) {\textbf{Zone tampon}};
\node[popdark, anchor=north] at (5, 3) {Pente forte};
\node[popdark, anchor=north] at (15, 11) {Pente forte};
\node[popdark, anchor=south] at (10, 7.5) {Pente faible};
\end{axis}
\end{tikzpicture}
\legende{Courbe de dosage d'un acide faible par une base forte : zone tampon (verte) autour de la demi-équivalence ($V_E/2$, $\mathrm{pH}=\mathrm{p}K_a$).}
\end{popfigure}

\begin{savaistu}[Pourquoi la zone d'efficacité est-elle $\mathrm{p}K_a \pm 1$ ?]
La relation $\mathrm{pH} = \mathrm{p}K_a + \log\frac{[\ce{A-}]}{[\ce{AH}]}$ montre que :
\begin{itemize}
  \item Si $\mathrm{pH} = \mathrm{p}K_a + 1$, alors $\log\frac{[\ce{A-}]}{[\ce{AH}]} = 1 \Rightarrow \frac{[\ce{A-}]}{[\ce{AH}]} = 10$. L'acide est minoritaire (9\%).
  \item Si $\mathrm{pH} = \mathrm{p}K_a - 1$, alors $\frac{[\ce{A-}]}{[\ce{AH}]} = 0,1$. La base est minoritaire (9\%).
\end{itemize}
Au-delà, l'espèce minoritaire est trop peu concentrée pour neutraliser efficacement l'ajout d'acide ou de base : le pouvoir tampon s'effondre. C'est pourquoi on choisit un couple dont le $\mathrm{p}K_a$ est le plus proche possible du pH visé.
\end{savaistu}

\courssub{Applications des solutions tampons}

\begin{savaistu}[Le tampon sanguin : une question de vie ou de mort]
Le pH du sang artériel est maintenu à \textbf{7,40 $\pm$ 0,05} (acidose si $< 7,35$, alcalose si $> 7,45$). Le système tampon principal est le couple \ce{H2CO3}/\ce{HCO3-} (acide carbonique / hydrogénocarbonate), avec $\mathrm{p}K_a' \approx 6,1$ (valeur apparente à 37°C, pression partielle \ce{CO2}).
\[ \mathrm{pH} = 6,1 + \log\frac{[\ce{HCO3-}]}{[\ce{H2CO3}]} \]
Le rapport $[\ce{HCO3-}]/[\ce{H2CO3}] \approx 20/1$ assure le pH 7,4.
\begin{itemize}
  \item \textbf{Régulation respiratoire :} Les poumons éliminent le \ce{CO2} (forme volatile de \ce{H2CO3}) $\Rightarrow$ ajustement rapide (minutes).
  \item \textbf{Régulation rénale :} Les reins réabsorbent/excrètent \ce{HCO3-} et excrètent \ce{H+} $\Rightarrow$ ajustement lent (heures/jours).
\end{itemize}
C'est un exemple parfait de \textbf{homéostasie} : un tampon chimique couplé à des systèmes physiologiques.
\end{savaistu}

\begin{savaistu}[Autres applications industrielles et courantes]
\begin{itemize}
  \item \textbf{Étalonnage des pH-mètres :} Solutions tampons de référence (pH 4,01 ; 7,00 ; 9,21 ou 10,00) traçables aux standards internationaux (NIST, AFNOR). Indispensable au laboratoire de la SONARA ou des brasseries (bière, bili-bili).
  \item \textbf{Industrie pharmaceutique :} Stabilité des médicaments (ex: aspirine, vaccins), formes injectables (sérum physiologique tamponné), collyres. Un pH mal contrôlé = principe actif dégradé ou douleur à l'injection.
  \item \textbf{Agroalimentaire :} Conservation (jus de bissap, sauces, conserves), texture (fromages, yaourts), fermentation contrôlée (vin de palme, bière). Le pH contrôle la croissance bactérienne.
  \item \textbf{Cosmétiques :} Crèmes, shampoings (pH $\approx$ 5,5 pour respecter le film hydrolipidique de la peau), colorations capillaires.
  \item \textbf{Chimie analytique :} Dosages par titration (maintenir pH pour indicateur), électrophorèse (tampons Tris, TAE, TBE pour séparation ADN).
\end{itemize}
\end{savaistu}

\courssub{Bilan synthétique}

\textbf{Définitions clés}
\begin{itemize}
  \item Une \cle{solution tampon} est une solution dont le \cle{pH varie peu} lors de l'addition \emph{modérée} d'un acide ou d'une base, ou lors d'une \emph{dilution modérée}.
  \item Elle contient un couple acide faible / base conjuguée \ce{AH}/\ce{A-} en quantités de matière \cle{voisines} (ni l'un ni l'autre négligeable).
  \item Le \cle{pouvoir tampon} $\beta$ mesure l'efficacité du tampon : il est d'autant plus grand que les concentrations sont élevées et que $\mathrm{pH} \approx \mathrm{p}K_a$.
  \item La \cle{zone tampon} d'une courbe de dosage est la région de pente faible autour de la \cle{demi-équivalence} ($V = V_E/2$), où $\mathrm{pH} = \mathrm{p}K_a$.
\end{itemize}

\textbf{Équations-types à connaître par cœur}
\begin{itemize}
  \item Équilibre du couple tampon : \[ \ce{AH + H2O <=> A- + H3O+} \]
  \item Action sur un acide ajouté : \[ \ce{A- + H3O+ -> AH + H2O} \]
  \item Action sur une base ajoutée : \[ \ce{AH + HO- -> A- + H2O} \]
  \item Préparation par demi-neutralisation : \[ \ce{AH + HO- -> A- + H2O} \quad \text{avec } n(\ce{HO-}) = \tfrac{1}{2}\,n(\ce{AH}) \]
\end{itemize}

\textbf{Formules essentielles}
\begin{center}
\begin{tabularx}{\linewidth}{|l|X|l|}
\hline
\rowcolor{popblueL} \textbf{Grandeur} & \textbf{Relation} & \textbf{Condition} \\
\hline
pH du mélange tampon & $\mathrm{pH} = \mathrm{p}K_a + \log\dfrac{[\ce{A-}]}{[\ce{AH}]} = \mathrm{p}K_a + \log\dfrac{n(\ce{A-})}{n(\ce{AH})}$ & couple \ce{AH}/\ce{A-} \\
\hline
Tampon idéal (pH = pKa) & $\mathrm{pH} = \mathrm{p}K_a$ & $n(\ce{AH}) = n(\ce{A-})$ \\
\hline
Mélange équimolaire & $C_a V_a = C_b V_b$ & même couple \\
\hline
Demi-neutralisation & $C_b V_b = \dfrac{C_a V_a}{2}$ & acide faible + base forte \\
\hline
Demi-équivalence & $\mathrm{pH}_{V_E/2} = \mathrm{p}K_a$ & courbe de dosage \\
\hline
Zone d'efficacité & $\mathrm{p}K_a - 1 \leq \mathrm{pH} \leq \mathrm{p}K_a + 1$ & $0{,}1 \leq \dfrac{[\ce{A-}]}{[\ce{AH}]} \leq 10$ \\
\hline
\end{tabularx}
\end{center}

\textbf{Méthodes express}
\begin{itemize}
  \item \cle{Préparer un tampon de pH $= \mathrm{p}K_a$} : (1) choisir le couple dont le $\mathrm{p}K_a$ est le plus proche du pH visé ; (2) mélanger l'acide et la base conjuguée en quantités égales, \emph{ou} neutraliser l'acide faible à moitié par une base forte ; (3) vérifier au pH-mètre.
  \item \cle{Calculer le pH d'un mélange} : dresser le bilan des quantités après réaction éventuelle, puis appliquer $\mathrm{pH} = \mathrm{p}K_a + \log\dfrac{n(\ce{A-})}{n(\ce{AH})}$.
  \item \cle{Lire une courbe de dosage} : repérer $V_E$ (saut de pH), puis $V_E/2$ ; le pH à la demi-équivalence donne directement le $\mathrm{p}K_a$ du couple.
\end{itemize}

\begin{aretenir}[Checklist avant le Bac]
\begin{itemize}
  \item $\square$ Je sais définir une solution tampon et citer ses trois propriétés (acide, base, dilution).
  \item $\square$ Je sais qu'une solution tampon contient un couple \ce{AH}/\ce{A-} en quantités voisines.
  \item $\square$ Je connais par cœur la relation $\mathrm{pH} = \mathrm{p}K_a + \log\dfrac{[\ce{A-}]}{[\ce{AH}]}$.
  \item $\square$ Je sais qu'à la demi-équivalence, $\mathrm{pH} = \mathrm{p}K_a$.
  \item $\square$ Je sais calculer les volumes pour préparer un tampon de $\mathrm{pH} = \mathrm{p}K_a$ (mélange équimolaire ou demi-neutralisation).
  \item $\square$ Je sais écrire les équations de réaction du tampon avec \ce{H3O+} et avec \ce{HO-}.
  \item $\square$ Je sais repérer la zone tampon (pente faible) sur une courbe $\mathrm{pH} = f(V)$.
  \item $\square$ Je sais que le pouvoir tampon est maximal quand $\mathrm{pH} = \mathrm{p}K_a$ et les concentrations élevées.
  \item $\square$ Je sais que la zone d'efficacité d'un tampon est $\mathrm{p}K_a \pm 1$.
  \item $\square$ Je connais le tampon du sang : couple \ce{H2CO3}/\ce{HCO3-}, $\mathrm{pH} \approx 7{,}4$.
  \item $\square$ Je sais citer deux applications : étalonnage des pH-mètres, industrie pharmaceutique ou agroalimentaire.
  \item $\square$ Je vérifie toujours l'homogénéité de mes calculs de quantités de matière avant d'appliquer la formule du pH.
\end{itemize}
\end{aretenir}

\begin{exercice}[Entraînement : Tampon ammoniacal]
On souhaite préparer \SI{500,0}{mL} d'une solution tampon \ce{NH4+}/\ce{NH3} de pH = 9,20. On dispose d'ammoniaque \SI{0,50}{mol.L^{-1}} ($pK_a(\ce{NH4+}) = 9,20$) et d'acide chlorhydrique \SI{0,50}{mol.L^{-1}}.
\begin{enumerate}
  \item Justifier le choix du couple.
  \item Calculer les volumes $V_{\ce{NH3}}$ et $V_{\ce{HCl}}$ à mélanger (méthode demi-neutralisation).
  \item On ajoute \SI{5,0}{mL} de \ce{HCl} \SI{0,10}{mol.L^{-1}} à \SI{100,0}{mL} de ce tampon. Calculer le nouveau pH. Comparer à l'ajout dans l'eau pure.
\end{enumerate}
\end{exercice}

\begin{exercice}[Analyse de courbe de dosage]
On dose \SI{20,0}{mL} d'un acide faible \ce{AH} (\SI{0,10}{mol.L^{-1}}) par une solution de soude \SI{0,10}{mol.L^{-1}}. La courbe $\mathrm{pH} = f(V_{\ce{NaOH}})$ montre un saut de pH à $V_E = \SI{20,0}{mL}$. À $V = \SI{10,0}{mL}$, le pH vaut 4,70.
\begin{enumerate}
  \item Déterminer la concentration de l'acide \ce{AH} (vérification).
  \item Déterminer le $\mathrm{p}K_a$ de \ce{AH}.
  \item Écrire l'équation de la réaction de dosage.
  \item Préciser les espèces majoritaires aux points : $V=0$, $V=10$ mL, $V=20$ mL, $V=25$ mL.
\end{enumerate}
\end{exercice}

\begin{corrige}[Exercice 1 : Tampon ammoniacal]
\begin{enumerate}
  \item Le couple \ce{NH4+}/\ce{NH3} a $\mathrm{p}K_a = 9,20$. Le pH visé est 9,20. Comme $\mathrm{pH}_{\text{visé}} = \mathrm{p}K_a$, ce couple est idéal (pouvoir tampon maximal).
  \item Méthode demi-neutralisation : $\ce{NH3 + H3O+ -> NH4+}$.
  On veut $n(\ce{NH4+}) = n(\ce{NH3})$ dans le mélange final (volume total \SI{500}{mL}).
  Soit $V$ le volume de \ce{NH3} \SI{0,50}{M} pris. $n_{\text{init}}(\ce{NH3}) = 0,50 V$.
  Volume de \ce{HCl} \SI{0,50}{M} à ajouter : $V_{\ce{HCl}}$. $n(\ce{HCl}) = 0,50 V_{\ce{HCl}}$.
  Condition demi-neutralisation : $n(\ce{HCl}) = \frac{1}{2} n_{\text{init}}(\ce{NH3}) \Rightarrow 0,50 V_{\ce{HCl}} = 0,25 V \Rightarrow V_{\ce{HCl}} = V/2$.
  Volume total $\approx V + V/2 = 1,5 V = \SI{500}{mL} \Rightarrow V \approx \SI{333,3}{mL}$.
  \textbf{Réponse :} $V_{\ce{NH3}} = \SI{333,3}{mL}$, $V_{\ce{HCl}} = \SI{166,7}{mL}$. (Compléter à \SI{500}{mL} à l'eau).
  \item \textbf{Dans le tampon (\SI{100}{mL}) :}
  $n(\ce{NH3}) = n(\ce{NH4+}) = \frac{0,50 \times 333,3}{5} = \SI{33,3}{mmol}$ (dans \SI{500}{mL}) $\Rightarrow$ dans \SI{100}{mL} : \SI{6,67}{mmol} chacun.
  Ajout \SI{5,0}{mL} \ce{HCl} \SI{0,10}{M} $\Rightarrow$ \SI{0,50}{mmol} \ce{H3O+}.
  Réaction : \ce{NH3 + H3O+ -> NH4+}.
  Nouveau bilan : $n(\ce{NH3}) = 6,67 - 0,50 = \SI{6,17}{mmol}$ ; $n(\ce{NH4+}) = 6,67 + 0,50 = \SI{7,17}{mmol}$.
  $\mathrm{pH} = 9,20 + \log\frac{6,17}{7,17} = 9,20 - 0,066 = \mathbf{9,13}$.
  Variation $\Delta\mathrm{pH} = -0,07$.
  \textbf{Dans l'eau pure (\SI{100}{mL}) :}
  $n(\ce{H3O+}) = \SI{0,50}{mmol}$ dans \SI{105}{mL} $\Rightarrow$ $[\ce{H3O+}] \approx \SI{4,76e-3}{M}$.
  $\mathrm{pH} = -\log(4,76\times10^{-3}) \approx \mathbf{2,32}$.
  Variation $\Delta\mathrm{pH} \approx 7 - 2,32 = 4,68$ unités !
  \textbf{Conclusion :} Le tampon résiste efficacement ($\Delta\mathrm{pH} \approx 0,07$), l'eau pas du tout ($\Delta\mathrm{pH} \approx 4,7$).
\end{enumerate}
\end{corrige}

\begin{corrige}[Exercice 2 : Analyse de courbe de dosage]
\begin{enumerate}
  \item À l'équivalence, $n(\ce{AH}) = n(\ce{HO-})$.
  $C_{\ce{AH}} \times \SI{20,0}{mL} = \SI{0,10}{mol.L^{-1}} \times \SI{20,0}{mL} \Rightarrow C_{\ce{AH}} = \SI{0,10}{mol.L^{-1}}$. (Cohérent avec l'énoncé).
  \item À la demi-équivalence ($V = V_E/2 = \SI{10,0}{mL}$), $\mathrm{pH} = \mathrm{p}K_a = \mathbf{4,70}$.
  \item $\ce{AH + HO- -> A- + H2O}$.
  \item
  \begin{center}
  \begin{tabular}{|c|c|c|}
  \hline
  \textbf{Point} & \textbf{Volume} & \textbf{Espèces majoritaires} \\
  \hline
  Début & $V=0$ & \ce{AH}, \ce{H3O+} (peu), \ce{A-} (peu) \\
  \hline
  Demi-éq. & $V=10$ mL & \ce{AH} et \ce{A-} (quantités égales) \textbf{\textless-- Zone tampon} \\
  \hline
  Équivalence & $V=20$ mL & \ce{A-}, \ce{Na+}, \ce{H2O} (hydrolyse \ce{A-} $\Rightarrow$ pH $>7$) \\
  \hline
  Après éq. & $V=25$ mL & \ce{HO-} (excès), \ce{A-}, \ce{Na+} \\
  \hline
  \end{tabular}
  \end{center}
\end{enumerate}
\end{corrige}

\findecours

\end{document}
